%% Manuel utilisateur Kesh
%% Version canonique française
%% À mettre à jour à chaque release Kesh (cf. ../README.md)

\documentclass[a4paper,11pt,oneside]{article}

\usepackage{polyglossia}
\setdefaultlanguage{french}
\usepackage{../shared/kesh-style}

%% Captures d'écran : graphicx est chargé par kesh-style.sty.
%% Les PNG sont dans ./screenshots/. Tant qu'une capture réelle n'est pas
%% fournie, le macro retombe sur _placeholder.png (approche texte-d'abord).
\graphicspath{{./screenshots/}}
\newcommand{\keshscreenshot}[2]{%
  \begin{figure}[htbp]
    \centering
    \IfFileExists{screenshots/#1}%
      {\includegraphics[width=0.85\linewidth]{#1}}%
      {\includegraphics[width=0.70\linewidth]{_placeholder}}%
    \caption{#2}
  \end{figure}%
}

\title{Manuel utilisateur Kesh}
\author{\keshAuthor}
\date{\keshReleaseDate}

\begin{document}

%%─────────────────────────────────────────────────────────────────────────────
%% Page de garde
%%─────────────────────────────────────────────────────────────────────────────

\keshcover{Manuel utilisateur}{Comment utiliser Kesh au quotidien : authentification, comptabilité, factures, import bancaire, rapports}{Documentation utilisateur}

%%─────────────────────────────────────────────────────────────────────────────
%% Table des matières
%%─────────────────────────────────────────────────────────────────────────────

\tableofcontents
\clearpage

%%─────────────────────────────────────────────────────────────────────────────
\section{Bienvenue dans Kesh}

\subsection{Qu'est-ce que Kesh ?}

Kesh est un logiciel de comptabilité et de gestion conçu pour les PME, indépendants et associations en Suisse. Il vous permet de tenir votre comptabilité en partie double, d'émettre des factures avec QR Bill conformes au standard suisse, d'importer vos relevés bancaires (CAMT.053 et CSV), de réconcilier automatiquement vos paiements et de générer tous les rapports légaux nécessaires.

\subsection{Pour qui est ce manuel ?}

Ce manuel s'adresse aux utilisateurs finaux : comptables internes d'une PME, indépendants tenant leur propre comptabilité, fiduciaires gérant des dossiers clients, ou trésoriers d'association. Il ne suppose aucune connaissance technique préalable, mais une familiarité avec les concepts comptables de base (débit, crédit, plan comptable, exercice) est nécessaire.

\subsection{Conventions du manuel}

\begin{itemize}
    \item Les \emph{noms d'écrans} et \emph{libellés de boutons} apparaissent en italique.
    \item Les chemins de navigation utilisent → : \emph{Paramètres} → \emph{Utilisateurs} signifie « cliquer sur Paramètres dans le menu, puis sur Utilisateurs ».
    \item \keshkey{Ctrl+S} indique une combinaison de touches clavier.
    \item Les boîtes colorées : \emph{Note} (information), \emph{Astuce}, \emph{Attention} (avertissement).
\end{itemize}

\begin{keshnote}
Ce manuel décrit Kesh version \keshVersion. Les captures d'écran et libellés exacts peuvent varier légèrement selon la version utilisée et les traductions activées. Pour la dernière version du manuel, consultez \keshWebsite.
\end{keshnote}

%%─────────────────────────────────────────────────────────────────────────────
\section{Premiers pas}

\subsection{Création de votre compte}

\keshscreenshot{setup-admin.png}{Écran de création du compte administrateur, au tout premier démarrage d'une installation vierge.}  % TODO capture: ecran /setup/admin (donnees factices : admin / Pierre Exemple)

L'accès à Kesh nécessite un compte utilisateur. Deux scénarios sont possibles :

\subsubsection{Création par un administrateur (cas courant)}

L'administrateur de l'installation crée votre compte : il choisit votre \textbf{nom d'utilisateur} et un \textbf{mot de passe initial} (12 caractères au moins par défaut), qu'il vous transmet. Pour le changer, passez par la procédure \emph{Mot de passe oublié} (ci-dessous) --- si l'administrateur a activé la récupération par email et renseigné votre adresse --- ou demandez-le à l'administrateur : aucun écran de changement de mot de passe en session n'existe encore (cf. section \emph{Gestion de votre compte}).

\subsubsection{Onboarding initial (premier utilisateur d'une instance)}

Si vous êtes le tout premier utilisateur d'une instance Kesh fraîchement installée, vous serez accueilli par un écran d'onboarding initial vous demandant de créer le compte administrateur. Voir le \emph{Manuel administrateur} pour les détails.

\subsection{Connexion}

Pour vous connecter :

\begin{enumerate}
    \item Ouvrir l'URL de votre instance Kesh dans un navigateur (par exemple \keshpath{https://compta.exemple.ch}).
    \item Saisir votre identifiant (nom d'utilisateur) et votre mot de passe.
    \item Cliquer sur \emph{Se connecter}.
\end{enumerate}

La session reste active 15 minutes, prolongée automatiquement à chaque action grâce au système de \emph{refresh token}. Au-delà de 30 jours d'inactivité totale, vous devrez vous reconnecter.

\subsection{Récupération de mot de passe}

Si vous avez oublié votre mot de passe :

\begin{enumerate}
    \item Sur l'écran de connexion, cliquez sur \emph{Mot de passe oublié ?} (le lien n'apparaît que si votre administrateur a activé la récupération par email).
    \item Saisissez votre nom d'utilisateur \textbf{ou} votre adresse email, puis validez. Un message de confirmation générique s'affiche dans tous les cas — par sécurité, Kesh ne révèle jamais si un compte existe ou non.
    \item Ouvrez l'email reçu et cliquez sur le lien de réinitialisation. Le lien est valable \textbf{30 minutes} et ne fonctionne qu'\textbf{une seule fois}.
    \item Choisissez votre nouveau mot de passe (12 caractères minimum) et confirmez-le.
    \item Reconnectez-vous avec le nouveau mot de passe. Par sécurité, toutes vos sessions ouvertes (autres onglets, autres appareils) ont été déconnectées.
\end{enumerate}

\begin{keshwarning}
Si vous ne recevez pas l'email dans les 5 minutes, vérifiez votre dossier spam/indésirables. Si l'email n'arrive pas du tout : soit aucune adresse email n'est renseignée sur votre compte, soit l'envoi d'emails n'est pas configuré sur l'instance — dans les deux cas, contactez votre administrateur, qui peut réinitialiser votre mot de passe manuellement.
\end{keshwarning}

\begin{keshtip}
Un lien expiré ou déjà utilisé affiche « lien invalide ou expiré » : cliquez sur \emph{Refaire une demande} pour recevoir un nouveau lien. Après 5 demandes en 15 minutes, patientez 30 minutes avant de réessayer. Pour changer votre mot de passe \textbf{sans} l'avoir oublié, refaites cette procédure (si la récupération par email est active), ou demandez à votre administrateur de le réinitialiser : aucun écran de changement en session n'existe encore.
\end{keshtip}

\subsection{Déconnexion}

Ouvrez le \emph{menu utilisateur} (coin supérieur droit, il affiche votre rôle) → \emph{Se déconnecter}. Tous les refresh tokens de votre session actuelle sont révoqués.

\begin{keshtip}
Si vous utilisez Kesh sur un ordinateur partagé, déconnectez-vous toujours après usage. Le mode \emph{HttpOnly Secure} des cookies de session protège contre certaines attaques, mais ne remplace pas une déconnexion explicite.
\end{keshtip}

\subsection{Vue d'ensemble de l'interface}

L'interface Kesh est organisée comme suit :

\begin{itemize}
    \item \textbf{Barre latérale gauche} : navigation principale, organisée en trois
    groupes dépliables — \emph{Quotidien} (accueil, carnet d'adresses, catalogue,
    facturation, échéancier, avoirs, factures et paiements fournisseurs, import de
    factures, import bancaire), \emph{Mensuel} (écritures, réconciliation, postes ouverts, rapports)
    et \emph{Administration} (plan comptable, exercices, comptes et profils
    bancaires, règles d'affectation, export global, paramètres ; et, pour les
    administrateurs, utilisateurs, facturation, sauvegarde et restauration).
    \item \textbf{Zone centrale} : contenu de la section active.
    \item \textbf{Barre supérieure} : le logo, un champ de recherche (à venir) et le
    \emph{menu utilisateur} (affiche votre rôle ; déconnexion et bascule du mode
    Guidé/Expert).
\end{itemize}

\subsection{Choix de la langue}

Kesh gère le français, l'allemand, l'italien et l'anglais pour les messages du
système. La langue est choisie lors de l'\textbf{onboarding}.

\begin{keshnote}
Un sélecteur permettant de changer la langue de l'interface en cours de session
est prévu pour une version ultérieure : dans la version actuelle, la langue est
fixée à l'onboarding et n'est pas modifiable depuis un menu.
\end{keshnote}

\begin{keshwarning}
\textbf{Les traductions allemande, italienne et anglaise n'ont pas été relues
par des locuteurs natifs.} Depuis la version 0.11.0, les quatre catalogues
portent rigoureusement les mêmes libellés, et des contrôles automatiques
empêchent qu'un écran retombe silencieusement sur le français. Un glossaire fixe
la terminologie comptable suisse (MWST, \emph{Geschäftsjahr}, \emph{aliquota})
et des contrôles systématiques ont écarté les confusions de sens --- mais
\textbf{rien ne garantit le naturel de la langue}. Si vous relevez une
formulation maladroite ou fautive, signalez-la sur \keshWebsite.

\emph{Avant la version 0.11.0}, une partie des écrans affichait du
\textbf{français} aux utilisateurs ayant choisi une autre langue --- non pas un
message d'erreur, mais du français correct, sans qu'aucun avertissement ne le
signale.
\end{keshwarning}

%%─────────────────────────────────────────────────────────────────────────────
\section{Onboarding}

À votre première connexion sur une nouvelle installation, Kesh vous guide via un assistant d'onboarding. Deux parcours sont proposés :

\keshscreenshot{onboarding-welcome.png}{Écran d'accueil de l'assistant d'onboarding : choix entre le mode exploration et le mode production.}  % TODO capture: 1er ecran onboarding

\subsection{Chemin A — Mode exploration}

Si vous souhaitez \textbf{tester Kesh sans données réelles}, choisissez le mode exploration. Kesh prépare :

\begin{itemize}
    \item Une société fictive (« Démo SA », ou sa traduction dans la langue de l'installation), avec une adresse et un numéro IDE d'exemple.
    \item Le plan comptable PME, dans la langue de l'installation.
    \item L'exercice de l'année en cours.
    \item Les quatre taux de TVA suisses et les réglages de facturation par défaut.
\end{itemize}

\textbf{Rien d'autre} : aucun contact, aucun produit, aucune écriture. Les listes sont vides~; c'est à vous de saisir les données sur lesquelles vous voulez essayer Kesh.

Vous pouvez librement créer, modifier, supprimer des données dans ce mode. L'installation ne porte qu'une société~: pour passer en mode production, réinitialisez-la par le bouton \emph{Réinitialiser pour la production} de la bannière de démonstration, qui ramène l'assistant au début. ⚠️ \emph{Ce bouton n'agit que si l'administrateur de l'installation l'a autorisé au préalable} (variable \texttt{KESH\_PRODUCTION\_RESET}, voir le manuel d'administration)~: sans cette autorisation, la réinitialisation est refusée, \textbf{même à un administrateur}. Elle vide \textbf{toutes} les données de la société --- y compris ce que vous avez saisi pendant l'exploration --- et \textbf{conserve} les utilisateurs et les clés d'API.

\subsection{Chemin B — Mode production}

Si vous souhaitez \textbf{tenir votre comptabilité réelle}, choisissez le mode production. L'assistant vous guide à travers une série d'écrans :

\keshscreenshot{onboarding-production.png}{Saisie des coordonnées légales en mode production (raison sociale, adresse, IDE).}  % TODO capture: ecran coordonnees production (donnees factices : ABC Web Sarl, CHE-123.456.789)

\begin{enumerate}
    \item \textbf{Type d'organisation} : indépendant, association ou PME. Ce choix détermine le plan comptable que Kesh met en place (voir ci-dessous) ; il ne se modifie plus après l'onboarding.
    \item \textbf{Langue comptable} : langue des comptes et des rapports, indépendante de la langue de l'interface (utile si vous travaillez en français mais produisez des rapports en allemand, par exemple).
    \item \textbf{Coordonnées légales} : raison sociale (ou nom pour un indépendant), adresse, et IDE (CHE-XXX.XXX.XXX), facultatif quel que soit le type d'organisation --- une entité inscrite au registre du commerce a intérêt à le renseigner.
    \item \textbf{Compte bancaire} : nom de la banque, IBAN et QR-IBAN éventuel (cette étape peut être passée et complétée plus tard).
\end{enumerate}

Dès que vous avez choisi la langue comptable, Kesh met en place le plan comptable de votre type d'organisation, libellé dans cette langue. À la finalisation, il crée votre premier exercice, l'année civile en cours. Vous êtes alors prêt à saisir vos premières écritures.

\begin{keshnote}
Vous pourrez à tout moment ajouter ou modifier vos comptes bancaires après l'onboarding (\emph{Administration} → \emph{Comptes bancaires}) et adapter votre plan comptable (cf. section \emph{Plan comptable}).
\end{keshnote}

\begin{keshtip}
Le mode guidé pas-à-pas est recommandé pour les nouveaux utilisateurs. Les utilisateurs expérimentés peuvent basculer en \emph{Mode expert} via le \emph{menu utilisateur} (coin supérieur droit), ce qui réduit les confirmations et accélère la saisie.
\end{keshtip}

%%─────────────────────────────────────────────────────────────────────────────
\section{Tableau de bord}

Le tableau de bord est votre page d'accueil après connexion. Il présente trois
tuiles de synthèse :

\keshscreenshot{dashboard.png}{Tableau de bord : dernières écritures, factures ouvertes et comptes bancaires.}  % TODO capture: dashboard avec donnees demo

\begin{itemize}
    \item \textbf{Dernières écritures} — les cinq plus récentes (date, numéro, libellé,
    montant), chacune menant à son détail.
    \item \textbf{Factures ouvertes} — le nombre de factures validées non soldées et leur
    \textbf{reste dû} total, dont la part échue. S'y ajoute, pour un comptable ou un
    administrateur, le nombre de factures à rappeler.
    \item \textbf{Comptes bancaires} — pour chaque compte, son \textbf{solde comptable} et,
    si un relevé a été importé, le solde de ce relevé et l'écart entre les deux
    (détails ci-dessous).
\end{itemize}

Un rôle \emph{Consultation} voit les mêmes chiffres, sans les boutons de saisie. Si une
tuile ne peut pas se charger, elle le dit — elle n'affiche jamais « Aucune écriture » sur
une comptabilité qui en contient.

\subsection{Solde comptable, relevé et écart}

Le chiffre affiché sous chaque compte bancaire est le \textbf{solde comptable} : il vient
du compte du grand livre lié au compte bancaire (\emph{Administration} → \emph{Comptes
bancaires}), \textbf{pas de la banque}. Ce compte lié ne peut être ni le \textbf{compte
débiteurs} ni le \textbf{compte créanciers} désignés dans \emph{Paramètres} →
\emph{Facturation}~: chaque encaissement écrirait sinon une écriture nulle. Kesh ne les
propose pas et refuse ce lien~; reliez le compte bancaire à son propre compte de banque.
Ce solde ne reflète un encaissement qu'une fois celui-ci comptabilisé, par un rapprochement
ou un règlement. Le total ne compte qu'une fois un
compte du grand livre partagé par plusieurs comptes bancaires.

Quand un relevé a été importé, la tuile affiche aussi \emph{Relevé du \dots} : le solde de
clôture du \textbf{dernier relevé} importé, à sa date. Seuls les relevés \textbf{CAMT.053}
portent un solde de clôture — un import CSV n'en a pas (\S\ref{sec:import-bancaire}). Si
le solde comptable à cette date diffère du relevé, la tuile affiche l'\textbf{écart}
(solde comptable − relevé) : typiquement un mouvement bancaire pas encore rapproché.

L'écart tient compte des \textbf{dates de valeur} : la banque arrête son relevé par date de
comptabilisation, Kesh date les écritures de rapprochement à la date de valeur. Un
mouvement comptabilisé par la banque le dernier jour du relevé, mais rapproché à une date
de valeur postérieure, ne crée donc pas d'écart. Trois limites :
\begin{itemize}
    \item seule une écriture \textbf{issue d'un rapprochement} est ainsi corrigée : une
    écriture saisie à la main pour un mouvement bancaire, à une autre date, apparaît dans
    l'écart — rien ne la relie au relevé ;
    \item quand plusieurs comptes bancaires partagent un même compte du grand livre,
    l'écart n'est pas calculable — la tuile l'indique : le solde de ce compte ne
    s'attribue à aucun d'eux ;
    \item après l'\textbf{annulation d'un rapprochement}, un mouvement comptabilisé par la
    banque après la date du relevé mais dont l'écriture d'origine la précédait peut laisser
    un écart fictif — l'écriture de contre-passation est datée du jour de l'annulation.
\end{itemize}

Un compte bancaire sans compte du grand livre lié n'a pas de solde comptable ; la tuile
affiche tout de même son dernier relevé. Si les comptes bancaires ne peuvent pas être
chargés, la tuile le dit au lieu de disparaître.

\begin{keshnote}
Les indicateurs plus riches (graphique de chiffre d'affaires, liste des
transactions à rapprocher, tâches en attente, personnalisation des tuiles) sont
prévus pour une version ultérieure.
\end{keshnote}

%%─────────────────────────────────────────────────────────────────────────────
\section{Gestion de votre compte}

Les \emph{Paramètres} regroupent, selon votre rôle, la configuration de
l'organisation, la facturation, les taux de TVA, les exercices comptables, les
comptes bancaires, les clés API et (pour les administrateurs) les utilisateurs.

\begin{keshnote}
Une page de \textbf{profil personnel} (nom, email, téléphone), un écran de
\textbf{changement de mot de passe en session} et la \textbf{liste des sessions
actives} ne sont pas encore exposés dans l'interface : ils sont prévus pour une
version ultérieure. En attendant, le mot de passe se réinitialise via la
procédure « Mot de passe oublié » ci-dessous, et un administrateur peut recréer
ou réinitialiser un compte.
\end{keshnote}

\subsection{Récupération de mot de passe oublié}

Si vous avez oublié votre mot de passe, utilisez le lien \emph{Mot de passe oublié ?} de l'écran de connexion (cf. section \emph{Premiers pas → Récupération de mot de passe}). Cette fonctionnalité nécessite qu'un serveur de messagerie (SMTP) ait été configuré par votre administrateur et qu'une adresse email soit associée à votre compte.

\subsection{Clés API (accès programmatique)}

Pour connecter Kesh à un logiciel tiers ou à une intégration (script, outil d'IA, etc.), vous pouvez créer une \textbf{clé API} (jeton d'accès personnel) depuis \emph{Paramètres} → \emph{Clés API} (\keshpath{/settings/api-keys}). Chaque clé porte un périmètre (\emph{lecture seule} ou \emph{lecture-écriture}) et une date d'expiration optionnelle ; le secret n'est affiché qu'une seule fois à la création. La gestion des clés est accessible aux rôles \emph{Comptable} et \emph{Administrateur}.

Une clé peut paraître « révoquée » sans que personne l'ait révoquée : Kesh l'a fait en réparant une installation touchée par le défaut \#528 (entrée \emph{Réparation de l'installation} du journal d'audit). Elle n'authentifie plus : créez-en une neuve.

\begin{keshnote}
La documentation complète de l'API externe (authentification, points d'accès, exemples) se trouve dans \keshpath{docs/api-external.md}. Les aspects d'administration (rotation, durcissement) sont décrits dans le manuel administrateur.
\end{keshnote}

%%─────────────────────────────────────────────────────────────────────────────
\section{Plan comptable}

Le plan comptable est la colonne vertébrale de votre comptabilité. Kesh fournit trois plans comptables, dont la structure suit le plan comptable PME suisse : un pour les PME, un pour les indépendants, un pour les associations. Le plan est choisi par le type d'organisation que vous indiquez à l'onboarding ; vous l'adaptez ensuite en ajoutant, modifiant ou archivant des comptes.

\subsection{Visualisation}

\emph{Administration} → \emph{Plan comptable}. La page affiche tous les comptes organisés selon leur structure hiérarchique :

\keshscreenshot{plan-comptable.png}{Plan comptable PME affiché par classes hiérarchiques.}  % TODO capture: page plan comptable

\begin{itemize}
    \item Classe 1 : Actifs (1000 Caisse, 1020 Banque, 1100 Débiteurs, 1171 Impôt préalable, etc.).
    \item Classe 2 : Passifs (2000 Créanciers, 2200 TVA due, 2800 Capital social, etc.).
    \item Classe 3 : Produits (3000 Ventes de marchandises, 3200 Prestations de services, etc.).
    \item Classe 4 : Charges de marchandises (4000 Achats de marchandises, 4200 Achats de matières premières, etc.).
    \item Classe 5 : Charges de personnel (5000 Salaires, 5700 Charges sociales, etc.).
    \item Classe 6 : Autres charges d'exploitation (6000 Loyers, 6100 Entretien et réparations, 6500 Administration, etc.).
    \item Classe 7 : Résultat hors exploitation et impôts.
\end{itemize}

\subsection{Ajout d'un compte}

\emph{Administration} → \emph{Plan comptable} → \emph{Nouveau compte}. Saisir :

\begin{itemize}
    \item \textbf{Numéro} : par exemple 1020.1 pour un sous-compte de la banque.
    \item \textbf{Libellé} : par exemple « UBS compte courant ».
    \item \textbf{Compte parent} : compte de niveau supérieur (1020 dans l'exemple).
    \item \textbf{Type} : Actif / Passif / Produit / Charge / Bilan / Résultat.
    \item \textbf{Rôle} : fonction du compte dans l'application (\emph{Créances clients}, \emph{TVA due}, \emph{Capital}, \emph{Résultat de l'exercice}\ldots) ou \emph{Aucun}. Voir la section \emph{Rôles des comptes} ci-dessous.
    \item \textbf{Postable} : décocher pour un compte de regroupement (compte titre) sur lequel on ne saisit pas d'écriture.
    \item \textbf{Devise} : CHF par défaut.
    \item \textbf{Compte de banque} : cocher si ce compte représente un compte bancaire (active la liaison avec l'import CAMT.053).
\end{itemize}

\subsection{Modification d'un compte}

Cliquer sur un compte ouvre son détail. Vous pouvez modifier le libellé, le compte parent, le \textbf{rôle} et l'indicateur \textbf{postable}. Le numéro n'est pas modifiable une fois des écritures saisies (pour préserver le journal d'audit). Le \textbf{type} se modifie librement tant que le compte ne porte aucune écriture ; au-delà, Kesh vous avertit avant de le laisser faire (voir ci-dessous).

\begin{keshwarning}
\textbf{Changer le type d'un compte qui porte des écritures reclasse tout son historique.} Le type — Actif, Passif, Produit, Charge — ne décrit pas seulement le compte : c'est lui qui décide dans quel état ses montants apparaissent. Les rapports le lisent \textbf{au moment où vous les demandez}, jamais au moment où l'écriture a été passée.

Concrètement : un compte de charge portant 1\,200.-- d'écritures de l'an dernier, passé en Actif, fait sortir ces 1\,200.-- du compte de résultat pour les faire entrer au bilan. \textbf{Le résultat de l'exercice écoulé change}, donc le report à nouveau, donc le bilan d'ouverture de l'exercice courant --- et \textbf{aucune écriture n'explique ce mouvement}.

C'est pourquoi Kesh \textbf{refuse le changement au premier essai} et vous indique combien d'écritures sont concernées, en nommant les \textbf{exercices clos} touchés s'il y en a. Le geste reste possible --- un compte mal typé à la création doit pouvoir être corrigé --- mais il exige une confirmation explicite, et il est journalisé sous son propre intitulé dans le journal d'audit.

\textbf{Si l'historique est important et que le type est faux}, préférez un compte neuf correctement typé : le reclassement rétroactif d'un exercice clos est difficile à justifier devant un contrôleur.
\end{keshwarning}

\subsection{Suppression d'un compte}

Un compte ne peut être supprimé que s'il n'a aucune écriture ni sous-compte. Sinon, vous pouvez l'\textbf{archiver} (le compte reste visible dans les écritures historiques mais n'est plus proposé à la saisie).

Pour retrouver un compte archivé, cocher \emph{Afficher les archivés} en haut de la page. Un bouton \textbf{Réactiver} permet de le remettre en service. La réactivation est refusée dans deux cas : si son \textbf{compte parent} est lui-même archivé (il faut le réactiver d'abord), ou si son \textbf{rôle} a été attribué à un autre compte pendant son archivage — le message d'erreur nomme alors le compte concerné, dont il faut retirer le rôle au préalable.

\label{sec:plan-comptable-archivage}

\begin{keshwarning}
\textbf{Avant d'archiver un compte de produit, pensez aux avoirs.} Un compte de produit imputé par des factures \textbf{déjà validées} reste nécessaire tant qu'un avoir peut être émis sur ces factures — ce qui n'a pas de limite de temps. Il en va de même du \textbf{compte débiteurs} et du \textbf{compte de différences d'arrondi} que ces factures ont mouvementés, ainsi que du \textbf{compte de TVA due} désigné dans les réglages --- celui du moment, non celui que la facture a crédité --- si la facture porte de la TVA. Une fois le compte archivé, l'émission d'un avoir sur ces factures est \textbf{refusée} (le motif est expliqué au \S\ref{sec:avoirs}).

Ce n'est pas bloquant : la marche à suivre est de \textbf{réactiver} le compte, d'émettre l'avoir, puis de le \textbf{ré-archiver}. Les deux mouvements étant journalisés, la manœuvre reste traçable.
\end{keshwarning}

\subsection{Rôles des comptes}

Chaque compte peut porter un \textbf{rôle} qui indique à Kesh à quoi il sert : compte de créances clients, de TVA due, de capital, de résultat de l'exercice, etc. Ce rôle est \textbf{explicite et modifiable} : Kesh ne devine jamais la fonction d'un compte à partir de son \emph{numéro}. C'est ce qui vous permet de renuméroter ou de restructurer votre plan comptable sans rien casser — le plan comptable suisse est un usage, pas une obligation légale.

Certains rôles sont \textbf{uniques} : un seul compte actif à la fois peut être le compte de créances clients, de TVA due, de résultat de l'exercice, etc. D'autres peuvent être portés par plusieurs comptes (capital, autres fonds propres). Si vous attribuez un rôle unique déjà pris, Kesh refuse la modification et vous indique quel compte le détient.

\begin{keshwarning}
Lors de la mise à jour vers cette version, Kesh a \textbf{pré-attribué les rôles} d'après les numéros du plan comptable standard (1100 → créances clients, 2200 → TVA due, etc.). Si vous aviez \textbf{renuméroté} vos comptes, vérifiez la colonne \emph{Rôle} de votre plan comptable et corrigez les attributions le cas échéant.
\end{keshwarning}

\begin{keshnote}
Dans cette version, l'indicateur \emph{postable} \textbf{bloque la saisie manuelle} : Kesh refuse une écriture saisie à la main dont une ligne vise un compte non imputable — un compte de regroupement, le compte de résultat de l'exercice ou un compte de clôture. Les autres chemins d'écriture le vérifient aussi~: la validation d'une facture contrôle le compte de produit de chaque ligne, l'enregistrement d'un règlement son compte de contrepartie, la saisie d'une facture fournisseur son compte de charge~; le rapprochement bancaire — manuel, éclaté, ou par acceptation d'une proposition — refuse un compte de contrepartie non imputable, et une règle d'affectation ne peut ni être créée sur un tel compte, ni y être déplacée, ni être réactivée s'il l'est devenu. Les réglages de facturation refusent un compte non imputable \textbf{au moment où vous le désignez}, et la fiche d'un compte bancaire refuse d'y lier un tel compte. Un compte \textbf{désigné dans les réglages de facturation} et devenu non imputable \emph{après} sa désignation --- la créance, la TVA due, le compte créanciers, la TVA récupérable --- est contrôlé \textbf{au moment où une écriture automatique veut l'utiliser}~: la validation d'une facture et la saisie (ou la complétion) d'une facture fournisseur sont alors refusées, avec un message qui nomme le compte et renvoie à \emph{Paramètres} → \emph{Facturation}, où un administrateur doit désigner à sa place un compte imputable --- l'un de ses sous-comptes, si vous l'avez scindé. Deux exceptions, voulues~: l'\textbf{avoir} n'est pas soumis à ce contrôle, pour qu'une facture émise reste toujours annulable --- il crédite la créance que la facture a débitée, lue sur l'écriture de la facture et non dans les réglages, et relit la TVA due dans les réglages au moment de l'avoir, sans la contrôler~; et le \textbf{compte de produit par défaut} n'est pas contrôlé à l'usage.

Trois cas échappent encore à ce contrôle, parce qu'un compte peut devenir non imputable \emph{après} avoir été choisi — typiquement quand vous lui créez des sous-comptes~: (1)~le compte comptable d'un compte bancaire \textbf{déjà lié} reste accepté, et continue de servir au rapprochement et aux règlements~; (2)~une règle d'affectation active dont le compte est devenu non imputable reste affichée «~Active~», mais Kesh ne la propose plus~; (3)~le compte de produit d'une \textbf{fiche article} n'est pas contrôlé sur la fiche — il l'est sur la ligne de facture, sauf s'il s'agit du compte de produit par défaut.

Une écriture déjà enregistrée \textbf{se modifie} tant que son exercice est ouvert, sauf si l'une de ses lignes est \textbf{lettrée}~: délettrez-la d'abord (voir «~Modifier une écriture~») --- et la modification applique la \textbf{même} garde, \textbf{à toutes ses lignes, y compris celles que vous ne touchez pas}~: si un compte est devenu non imputable depuis la saisie, l'enregistrement est refusé avec un message qui \textbf{nomme} ce compte, et il faut le remplacer.
\end{keshnote}

\subsection{Les comptes de clôture 9000, 9100 et 9200}

Les plans comptables livrés avec Kesh comportent une classe~9 — \emph{Bilan d'ouverture} (9000), \emph{Compte de résultat} (9100) et \emph{Bilan de clôture} (9200). Dans les logiciels qui passent des \textbf{écritures de clôture}, ces comptes reçoivent le virement des soldes en fin d'exercice.

\textbf{Kesh n'en passe aucune}~: le bilan, le compte de résultat et le report à-nouveau sont \textbf{recalculés en permanence} à partir de vos écritures. Ces trois comptes n'ont donc, ici, aucun usage — et Kesh \textbf{refuse toute écriture saisie à la main} qui les vise, exactement comme pour le compte \emph{Résultat de l'exercice}, dont il calcule aussi le solde.

\begin{keshnote}
\textbf{Cette protection ne se limite pas à la saisie manuelle}~: comme pour tout compte non imputable (voir «~Rôles des comptes~»), la validation d'une facture refuse un compte de produit non imputable — sauf s'il s'agit du compte de produit \emph{par défaut} que vous avez configuré —, l'enregistrement d'un règlement refuse un compte de contrepartie non imputable, la saisie d'une facture fournisseur refuse un compte de charge non imputable, et le rapprochement bancaire (manuel, éclaté, propositions et règles d'affectation), les réglages de facturation et le compte comptable d'un compte bancaire refusent eux aussi un compte de clôture. Ce qui se passe quand un compte devient non imputable \emph{après} avoir été choisi --- le contrôle à l'usage des comptes des réglages de facturation, et les cas qui échappent encore au contrôle --- est décrit à la section «~Rôles des comptes~».
\end{keshnote}

\begin{keshwarning}
\textbf{Si vous arrivez d'un autre logiciel, n'utilisez pas le 9000 pour votre bilan d'ouverture.} La reprise se fait par l'écran \emph{Administration} → \emph{Soldes de départ} (voir «~Reprise de comptabilité~»), et votre report accumulé se pose sur le compte \emph{Bénéfice/perte reporté} (2970 dans les plans standards).

Avant la version qui introduit cette règle, un à-nouveau passé par le 9000 partait \textbf{en charges} sans le moindre avertissement~: le bilan restait équilibré, mais le résultat de l'exercice était faux du montant du bilan, et le report à-nouveau l'était aussi.
\end{keshwarning}

\begin{keshnote}
La mise à jour \textbf{ferme ces comptes pour l'avenir, elle ne corrige pas le passé}. Si des écritures y ont déjà été passées, elles y restent~: les déplacer d'office serait une correction \textbf{invisible dans vos livres}, ce que la loi n'admet pas (art.~958f~CO). Consultez le \textbf{grand livre} du compte concerné, puis \textbf{contre-passez} les écritures fautives et ressaisissez-les sur les bons comptes.

Si votre plan a été renuméroté et que le numéro 9000 y désigne un compte que vous utilisez, la page \emph{Plan comptable} vous permet de le rouvrir~: ouvrez sa fiche et rétablissez l'indicateur \emph{postable}.
\end{keshnote}

\subsection{Adapter le plan à votre secteur}

Le plan PME couvre la plupart des besoins, mais chaque PME l'adapte à son activité. Quelques exemples concrets :

\begin{itemize}
    \item \textbf{Commerce de détail} (ex. \emph{Boutique Pierre Exemple Sàrl}) : créez des sous-comptes de stock sous la classe 1 (ex. 1200.1 « Stock textile », 1200.2 « Stock accessoires ») et ventilez les ventes par rayon en sous-comptes de 3000 (ex. 3000.1 « Ventes textile », 3000.2 « Ventes accessoires ») pour suivre vos marges par catégorie.
    \item \textbf{Cabinet médical ou prestataire de services} (ex. \emph{Dr Marie Exemple}) : peu de stock, l'essentiel est dans les produits de services (3200 « Prestations de services ») et les charges de personnel (classe 5). Créez un sous-compte 3200.1 par type de prestation si vous voulez analyser votre activité.
    \item \textbf{Agence web / indépendant} (ex. \emph{ABC Web Sàrl}) : facturation au forfait ou à l'heure via le catalogue produits ; suivez les sous-traitances en sous-compte de la classe 4 (ex. 4400 « Prestations de tiers »).
    \item \textbf{Artisan} (ex. \emph{Menuiserie Exemple}) : distinguez achats de matières premières (classe 4) et outillage (classe 1, immobilisations) ; un sous-compte 4000.1 « Bois » / 4000.2 « Quincaillerie » facilite le suivi des coûts de matière.
    \item \textbf{Association} (ex. \emph{Association Exemple}) : le plan association porte déjà les comptes de produits propres à une association --- 3000 « Cotisations des membres », 3100 « Dons reçus », 3200 « Subventions », 3300 « Produits des manifestations » ; créez des sous-comptes (ex. 3300.1 « Fête d'été ») si vous voulez suivre chaque manifestation à part. Pas de TVA si l'association n'est pas assujettie.
\end{itemize}

\begin{keshtip}
Archivez (plutôt que supprimer) les comptes du plan que vous n'utilisez pas : ils disparaissent de la saisie tout en restant disponibles pour l'historique. C'est particulièrement utile pour un indépendant ou une petite structure dont le plan complet est surdimensionné.
\end{keshtip}

%%─────────────────────────────────────────────────────────────────────────────
\section{Écritures comptables}

\subsection{Saisie d'une écriture en partie double}

\emph{Mensuel} → \emph{Écritures} → \emph{Nouvelle écriture}. Renseigner :

\keshscreenshot{ecriture-saisie.png}{Formulaire de saisie d'une écriture en partie double avec ses lignes débit/crédit.}  % TODO capture: form nouvelle ecriture

\begin{itemize}
    \item \textbf{Date} : date de la transaction (entre la date de début et la date de fin de l'exercice).
    \item \textbf{Libellé} : description de l'écriture (par exemple « Vente facture F-2026-042 »).
    \item \textbf{Référence externe} : numéro de facture, référence bancaire, etc. (optionnel).
    \item \textbf{Lignes débit / crédit} : minimum 2 lignes, somme débits = somme crédits.
\end{itemize}

Pour chaque ligne :

\begin{itemize}
    \item \textbf{Compte} : sélection dans le plan comptable.
    \item \textbf{Montant} : positif uniquement (la colonne débit/crédit indique le sens).
    \item \textbf{Description ligne} : précision supplémentaire (optionnel).
\end{itemize}

\subsection{Exemple de saisie}

Vente d'une marchandise à crédit pour CHF 1'000 + TVA 8.1\% :

\begin{itemize}
    \item Débit 1100 (Débiteurs) : CHF 1'081.00
    \item Crédit 3000 (Ventes de marchandises) : CHF 1'000.00
    \item Crédit 2200 (TVA due) : CHF 81.00
\end{itemize}

Total débit = Total crédit = CHF 1'081.00. L'écriture peut être validée.

\subsection{Cycle de vie d'une écriture}

Une écriture enregistrée est \textbf{immédiatement comptabilisée}. Kesh ne connaît pas
d'état « brouillon » : il n'y a rien à valider, et une écriture apparaît dans les rapports
dès son enregistrement. Elle reste \textbf{modifiable} --- et \textbf{supprimable} ---
tant que son exercice est ouvert, dans le cadre décrit ci-dessous --- sauf si l'une de ses
lignes est \textbf{lettrée}~: délettrez-la d'abord (\S\ref{sec:lettrage}).

\subsection{Modifier ou supprimer une écriture}\label{sec:modifier-ecriture}

Une erreur de saisie --- débit et crédit inversés, un mauvais compte, une date ou un libellé
faux --- se corrige \textbf{là où elle est}. Sur la fiche d'une écriture (depuis la liste des
écritures, cliquez sur la ligne), le bouton \textbf{Modifier} ouvre le formulaire de saisie,
pré-rempli~: corrigez, puis \emph{Valider}. La fiche se recharge avec l'écriture corrigée.

Le bouton n'est offert que si \textbf{toutes} ces conditions tiennent~:

\begin{itemize}
    \item l'\textbf{exercice} de l'écriture est \textbf{ouvert}~;
    \item \textbf{aucun exercice postérieur n'est clôturé}~: le bilan étant cumulatif depuis
          l'origine, un exercice suivant clôturé reprend cette écriture dans ses soldes ---
          la modifier réécrirait son bilan. Depuis la v0.13.0, la même garde vaut pour
          \textbf{toute} écriture de cet exercice~: rien ne s'y saisit, ne s'y règle, ne s'y
          dévalide ni ne s'y supprime tant qu'un exercice postérieur est clôturé
          (\S\ref{sec:exercices-desordre})~;
    \item sa date est \textbf{après} la période verrouillée (\S\ref{sec:verrou-periode})~---
          et la nouvelle date aussi~;
    \item \textbf{aucune pièce} ne la possède~: l'écriture d'une facture, d'un avoir, d'une
          facture fournisseur, d'un règlement ou d'un solde, d'un rapprochement bancaire se
          corrige par le chemin de sa pièce~;
    \item elle n'est \textbf{ni contre-passée, ni elle-même une contre-passation}~;
    \item \textbf{aucune de ses lignes n'est lettrée} --- sinon, délettrez-la d'abord~: le
          motif de la fiche porte le code du groupe, en lien vers le groupe, où se fait le
          délettrage (\S\ref{sec:lettrage}), puis modifiez~; quand c'est le seul motif, le
          délettrage aboutit, sauf geste concurrent sur les mêmes lignes.
\end{itemize}

En pratique, ce sont les écritures \textbf{saisies à la main}, l'\textbf{écriture
d'ouverture} et les \textbf{compléments de soldes de départ}
(\S\ref{sec:soldes-depart}). Sinon, la fiche affiche le motif à la place du bouton, et la
correction passe par une contre-passation (ci-dessous) --- ou, pour une pièce, par le chemin
de la pièce~; ou, pour une écriture lettrée, par le délettrage.

Ce qui ne change pas en modifiant~:

\begin{itemize}
    \item le \textbf{numéro} de l'écriture~; la \textbf{date} peut changer, mais reste dans
          l'exercice de l'écriture (l'ordre des numéros peut alors cesser de suivre celui
          des dates)~;
    \item les \textbf{contrôles de la saisie}~: équilibre, comptes actifs et imputables ---
          sur \textbf{toutes} les lignes, même celles que vous ne touchez pas. Une ligne
          dont le compte a été archivé, ou est devenu non imputable, est signalée dans le
          formulaire~: remplacez le compte.
\end{itemize}

\textbf{Supprimer.} Dans le \textbf{même cadre}, la fiche offre aussi \textbf{Supprimer}, à
côté de \emph{Modifier}~: une écriture saisie par erreur --- un doublon, une écriture qui
n'aurait pas dû exister --- disparaît sans écriture inverse. Une confirmation vous rappelle
deux choses avant d'agir~: le \textbf{numéro} de l'écriture \textbf{ne sera pas réattribué}
(la numérotation garde un trou, voir «~Numérotation~»), et la suppression est \textbf{inscrite au
journal d'audit} avec le contenu complet de l'écriture, lignes comprises, et son auteur. Quand
l'écriture ne se modifie pas, elle ne se supprime pas non plus~: la fiche n'offre alors ni
\emph{Modifier} ni \emph{Supprimer}, et affiche le motif.

\textbf{La correction reste apparente.} Chaque modification est inscrite au journal
d'audit (\S\ref{sec:tracabilite}), avec l'écriture \textbf{avant et après}, lignes
comprises, l'auteur --- ou la clé d'API --- et l'heure. C'est l'exigence de l'art.~958f CO
et de l'Olico~: \emph{on ne demande pas de ne jamais se tromper, on demande que la
correction soit apparente}. Une fois l'exercice clôturé, ou la période verrouillée,
l'écriture est figée et ne se corrige plus que par contre-passation. La fiche d'une écriture
modifiée porte la mention \textbf{«~Modifiée~»}, et son lien \textbf{Historique} ouvre le
journal d'audit filtré sur elle --- sa création, chacune de ses modifications, avant et après
(rôles Administrateur et Comptable).

\begin{keshwarning}
\textbf{Modifier une écriture change les rapports déjà tirés} sur sa période --- c'est le
sens même de la modification, les rapports se calculant à la volée. Pour figer un trimestre
déclaré (TVA), \textbf{verrouillez la période}~: c'est le geste qui fige
(\S\ref{sec:verrou-periode}).
\end{keshwarning}

\begin{keshnote}
Si quelqu'un a modifié l'écriture entre le moment où vous avez ouvert le formulaire et
celui où vous enregistrez, Kesh vous le dit et recharge la fiche~: reprenez votre
correction sur l'état présent. Il en va de même pour une suppression refusée en cours de
route~: le message dit pourquoi, et la fiche se recharge.
\end{keshnote}

\begin{keshtip}[title=\faLightbulb\ Pourquoi il n'y a pas de «~brouillon~»]
Kesh ne connaît pas d'état intermédiaire~: ce que vous validez part directement aux livres
--- et reste corrigeable tant que l'exercice est ouvert (une écriture lettrée se délettre
d'abord). La contrepartie est qu'il faut se
relire avant de clôturer --- et que l'écran de saisie vous y aide en refusant toute
écriture déséquilibrée.
\end{keshtip}

\subsection{Arrêter les livres à une date : le verrou de période}
\label{sec:verrou-periode}

La clôture d'exercice est annuelle. Or on remet ses comptes bien plus souvent que cela~: un
décompte de TVA trimestriel est déposé, un exercice est remis au fiduciaire, un mois est arrêté.
Rien n'empêchait alors de \textbf{saisir après coup} une écriture datée d'une période déjà
déclarée --- et comme les rapports se recalculent à la volée, le chiffre changeait \textbf{sans
que rien ne le signale}.

Le \textbf{verrou de période} ferme cela. Dans \emph{Paramètres} $\rightarrow$ \emph{Verrou de
période}, indiquez une date~: «~les écritures jusqu'au 31 mars 2026 inclus sont verrouillées~». Aucune écriture ne pourra plus y être datée, \textbf{par aucun chemin} ---
ni la saisie manuelle, ni la validation d'une facture antérieure, ni un rapprochement bancaire.
Et aucune écriture déjà passée dans la période ne pourra plus en \textbf{disparaître} ni
être \textbf{modifiée}~: \textbf{verrouiller est le geste qui fige} un trimestre déclaré. La
\textbf{dévalidation} d'une facture, qui emporte son écriture, est refusée si cette écriture est
datée d'une période verrouillée.
Le \textbf{lettrage} des lignes de la période se fige avec elles~: un groupe dont toutes les
lignes sont dans la période verrouillée (ou dans des exercices clôturés) ne se lettre ni ne se
délettre. L'annulation d'un règlement, d'un solde, d'un rapprochement, d'un paiement ou d'une
facture fournisseur dont le lettrage s'est figé avec la période est donc \textbf{refusée}~: un
administrateur fait reculer le verrou avant la date la plus récente du lettrage --- en général
celle du dernier règlement --- pour la permettre, et rouvre aussi les exercices clôturés si
l'exercice de cette date l'est \textbf{ou s'il est suivi d'un exercice clôturé} --- jusqu'à
celui-ci, en commençant par le plus récent.

\begin{keshtip}[title={\faLightbulb\ Deux gestes, deux exigences différentes}]
\textbf{Verrouiller} est un geste d'hygiène~: un comptable comme un administrateur peut le
faire, aussi souvent qu'il le souhaite, sans justification. \textbf{Déverrouiller} défait une
garantie~: c'est réservé à un \emph{administrateur}, un \textbf{motif est obligatoire}, et
l'opération est inscrite au journal d'audit --- exactement comme la réouverture d'un exercice
clôturé.
\end{keshtip}

\begin{keshwarning}[title={\faExclamationTriangle\ Ce que le verrou n'empêche pas, et c'est voulu}]
\textbf{Une écriture d'une période verrouillée reste corrigeable par contre-passation.} La
contre-passation est datée du \textbf{jour}, donc après le verrou~: elle passe. Le verrou arrête
ce qui \textbf{entre dans le passé}, jamais ce qui \textbf{corrige depuis le présent}. C'est
cette distinction qui empêche le verrou de rendre vos livres incorrigibles --- et c'est pourquoi
Kesh refuse de poser une borne à la date du jour ou dans le futur.
\textbf{Une exception}~: l'annulation d'un règlement (ou d'un paiement) dont le
\textbf{lettrage} s'est figé avec la période est refusée --- elle défairait un lettrage de la
période verrouillée~; le verrou doit reculer pour la permettre (ci-dessus).
\end{keshwarning}

Le verrou et la clôture d'exercice sont \textbf{indépendants}~: clôturer ne touche pas la borne,
et la borne n'empêche ni la clôture ni la réouverture. Le premier est un état d'exercice, le
second une borne mobile à l'intérieur d'un exercice ouvert.

\subsection{Corriger une écriture : la contre-passation}\label{sec:contre-passation}

Sur la fiche d'une écriture, le bouton \textbf{Contre-passer} crée l'\textbf{écriture
inverse} --- mêmes comptes, débit et crédit échangés --- à la date du jour, dans l'exercice
ouvert. \textbf{L'écriture d'origine n'est pas modifiée} --- ni ses montants, ni ses comptes,
ni sa date, ni son libellé~: les deux restent visibles au grand livre, et chacune renvoie à
l'autre depuis sa fiche. Seule marque posée sur l'origine~: ses lignes sur un compte
lettrable encore ouvertes sont \textbf{lettrées} avec leur miroir dans l'écriture inverse
(groupe «~contre-passation~») --- la ligne et sa correction se soldent l'une l'autre~; une
ligne déjà lettrée garde son groupe, et son miroir reste ouvert. C'est ce que l'art.~958f CO et l'Olico
exigent --- non pas qu'on ne se trompe jamais, mais que \textbf{la correction soit
apparente}.

Le lien figure aussi dans l'export comptable complet (colonne
\texttt{reverses\_entry\_id} du fichier \texttt{journal\_entries.csv}), afin qu'une
fiduciaire ou un organe de révision le retrouve sans ouvrir le logiciel.

Une écriture n'est pas toujours contre-passable, et l'écran vous dit alors pourquoi plutôt
que de vous laisser essayer~: une écriture qui appartient à une \textbf{pièce} se corrige
par le chemin de cette pièce (une facture client par un \textbf{avoir}, une facture
fournisseur par son \textbf{annulation}, même payée (\S\ref{sec:annuler-facture-fournisseur}), un \textbf{règlement} --- client ou fournisseur --- par
l'annulation du règlement depuis la fiche de la facture), une écriture déjà contre-passée ne l'est pas deux
fois, et une écriture dont un \textbf{compte a été archivé} attend que vous réactiviez ce
compte --- son numéro vous est indiqué.

\begin{keshnote}
La contre-passation est \textbf{le} chemin de correction d'une écriture désormais
figée~: écriture d'un exercice clôturé ou suivi d'un exercice clôturé, d'une période
verrouillée, ou d'une pièce. Tant que l'exercice est ouvert, une écriture saisie à la main
--- délettrée d'abord si l'une de ses lignes est lettrée ---
se corrige plus simplement par \textbf{modification}, ou disparaît par \textbf{suppression}
(ci-dessus) --- dans le même cadre. Une écriture disparaît donc par deux voies, toutes deux
tracées~: sa \textbf{suppression} depuis sa fiche, inscrite au journal d'audit avec son contenu
complet~; et la \textbf{dévalidation} d'une facture, qui la repasse en brouillon avec l'écriture
que sa validation avait créée (voir «~Supprimer ou dévalider une facture~») --- une décision de
facturation, inscrite au journal d'audit sous son propre nom. \emph{Exception~: l'annulation
d'un règlement (ou d'un paiement) dont le lettrage s'est figé avec la période est refusée~; le
verrou doit reculer pour la permettre (\S\ref{sec:verrou-periode}).} \emph{La restauration d'une
sauvegarde, qui remet les livres dans l'état de l'archive, fait exception à cette règle.}

L'historique complet des corrections apportées à vos écritures se \textbf{consulte à
l'écran}~: \emph{Administration} → \emph{Journal d'audit} (\S\ref{sec:tracabilite}), avec
ses filtres et son export CSV, dans votre langue --- et par l'interface de programmation. Un
verrou \textbf{plus fin} que l'exercice existe aussi~: le verrou de période arrête les livres
à la date de votre choix --- le trimestre TVA déclaré, par exemple
(\S\ref{sec:verrou-periode}).
\end{keshnote}

\subsection{Numérotation}

Chaque écriture reçoit un numéro séquentiel par société et par exercice, tiré d'un
\textbf{compteur qui ne redescend jamais}.

\textbf{Un numéro attribué ne sera jamais réattribué à une autre écriture.} C'est la
garantie qui compte devant un contrôle. Si la dernière écriture d'un exercice disparaît --- par
sa suppression ou par la dévalidation de sa facture ---, la suivante prendra le numéro d'après --- pas celui qui vient d'être libéré.
Jusqu'à cette version, ce numéro était repris~: une écriture au contenu tout différent
portait alors un numéro déjà employé, \textbf{sans aucun message et sans trace visible}.

\textbf{Cette numérotation peut en revanche comporter des trous}, et c'est assumé. Deux chemins
en creusent~: la \textbf{suppression} d'une écriture depuis sa fiche (ci-dessus), et la
\textbf{dévalidation} d'une facture, qui emporte son écriture (voir «~Supprimer ou dévalider une
facture~»). \emph{Un trou se voit et s'explique} --- le journal d'audit nomme ce qui a été
supprimé, quand et par qui. C'est la réattribution, elle, qui ne se voyait pas~; le combler
exigerait de renuméroter des écritures existantes, ce qu'une comptabilité ne fait pas.

\subsection{Recherche d'écritures}

\emph{Mensuel} → \emph{Écritures}. Filtres disponibles :

\begin{itemize}
    \item Période (date début / date fin).
    \item Compte (un ou plusieurs).
    \item Montant (>=, <=).
    \item Libellé (recherche full-text).
    \item Référence externe.
    \item Statut (brouillon / validée / annulée).
\end{itemize}

\begin{keshtip}
La recherche full-text utilise l'index MariaDB FULLTEXT (Story 7-4). Pour rechercher un mot incomplet, ajoutez \texttt{*} : \texttt{Marc*} matche \texttt{Marcel}, \texttt{Marcus}, \texttt{Marchand}.
\end{keshtip}

%%─────────────────────────────────────────────────────────────────────────────
\section{Exercices comptables}

\subsection{Concept}

Un exercice comptable est une période sur laquelle s'établissent les comptes annuels (bilan + compte de résultat). En Suisse, l'exercice est généralement civil (1\textsuperscript{er} janvier -- 31 décembre), mais d'autres périodes sont possibles (décalé, court, long).

\subsection{Création d'un nouvel exercice}

\emph{Administration} → \emph{Exercices comptables} → \emph{Nouvel exercice}. Renseigner :

\begin{itemize}
    \item \textbf{Date de début} et \textbf{Date de fin}.
    \item \textbf{Libellé} : par exemple « Exercice 2026 ».
    \item \textbf{Type} : Standard / Court / Long / Décalé.
\end{itemize}

Le nouvel exercice s'ouvre vide. Les comptes du bilan sont automatiquement reportés depuis le bilan de clôture de l'exercice précédent (si clôturé).

Un exercice \textbf{ne se crée pas avant un exercice clôturé} (depuis la v0.13.0)~: si un exercice dont la date de début est postérieure à celle demandée est déjà clôturé, Kesh refuse la création et nomme cet exercice --- son bilan, cumulatif, reprendrait tout ce que le nouvel exercice contiendrait. Pour créer malgré tout un exercice antérieur, un administrateur rouvre d'abord les exercices clôturés postérieurs, en commençant par le plus récent.

\subsection{Clôture d'un exercice}

\keshscreenshot{exercice-cloture.png}{Écran de clôture d'un exercice comptable.}  % TODO capture: ecran cloture exercice

\emph{Administration} → \emph{Exercices comptables} → sélectionner l'exercice → \emph{Clôturer}.

Les exercices se clôturent \textbf{dans l'ordre} (depuis la v0.13.0)~: un exercice ne se clôture qu'après \textbf{tous ceux qui le précèdent}. Tant qu'un exercice antérieur est encore ouvert, le bouton \emph{Clôturer} est désactivé, et son infobulle nomme l'exercice à clôturer d'abord --- le plus ancien encore ouvert~; Kesh refuse de même la clôture demandée par l'API. Le bilan étant cumulatif depuis l'origine, clôturer un exercice en laissant ouvert un exercice antérieur laisserait son bilan changer après coup.

La clôture est un simple \textbf{verrou} : elle ne génère aucune écriture. Le report des soldes de bilan d'un exercice à l'autre est \textbf{calculé en temps réel} depuis l'origine (report à-nouveau automatique), sans écriture de clôture ni de report à passer. La clôture :

\begin{itemize}
    \item Verrouille l'exercice : plus aucune écriture ne peut y être \textbf{saisie} tant qu'il reste clôturé (immutabilité, CO art. 957-964). C'est aussi la clôture qui \textbf{ferme la modification}~: une écriture se modifie tant que son exercice est ouvert (sauf si l'une de ses lignes est lettrée~: elle se délettre d'abord). Un exercice ne se clôturant qu'après tous ceux qui le précèdent, aucun exercice postérieur n'est clôturé tant que le sien est ouvert. Si des données antérieures à la v0.13.0 (installation mise à jour, sauvegarde restaurée) portent déjà un exercice clôturé après un exercice ouvert, \textbf{rien} ne s'enregistre plus dans cet exercice ouvert --- ni saisie, ni règlement, ni modification, ni suppression~: le bilan étant cumulatif, toute écriture y changerait le bilan clos. L'écran des exercices le signale et indique comment rétablir l'ordre (\S\ref{sec:exercices-desordre}). La clôture ferme de même la \textbf{suppression} d'une écriture --- celle de l'exercice clôturé, et celle de tout exercice antérieur.
    \item Inscrit l'opération dans le journal d'audit (\texttt{audit\_log}).
    \item N'affecte pas le report à-nouveau, qui reste calculé automatiquement — les rapports finaux (bilan, compte de résultat, balance) demeurent consultables et exportables (CSV/PDF).
\end{itemize}

\begin{keshwarning}
La clôture d'un exercice est une opération critique. Effectuez d'abord une sauvegarde complète. Une fois clôturé, l'exercice ne peut plus être modifié tant qu'il reste clôturé ; seul un administrateur peut le rouvrir (voir ci-dessous).
\end{keshwarning}

\subsection{Réouverture d'un exercice clôturé}\label{sec:reouverture-exercice}

En cas d'erreur découverte après clôture, un \textbf{administrateur} peut rouvrir l'exercice directement depuis l'interface : \emph{Paramètres} → \emph{Exercices comptables} → bouton \emph{Réouvrir} (visible uniquement pour les administrateurs sur un exercice clôturé). Un \textbf{motif est obligatoire} ; il est conservé, avec l'auteur et l'horodatage, dans le journal d'audit (\texttt{audit\_log}, conservation 10 ans). L'exercice rouvert redevient modifiable --- ses écritures saisies à la main se modifient de nouveau, si aucun exercice postérieur n'est clôturé et après délettrage pour celles dont une ligne est lettrée ---, puis pourra être re-clôturé.

\begin{keshtip}
On rouvre les exercices \textbf{du plus récent vers l'ancien} : si un exercice postérieur est encore clôturé, Kesh refuse la réouverture et indique lequel rouvrir d'abord. C'est l'ordre inverse de la clôture, qui va du plus ancien vers le plus récent~: ensemble, ils garantissent qu'aucun exercice ouvert ne précède un exercice clôturé. Effectuer cette opération avec parcimonie pour préserver l'intégrité légale.
\end{keshtip}

\subsubsection{Exercices dans le désordre --- données antérieures à la v0.13.0}\label{sec:exercices-desordre}

Les exercices se clôturant désormais dans l'ordre, un exercice ouvert ne peut plus précéder un exercice clôturé. Des données plus anciennes --- une installation mise à jour, une sauvegarde restaurée --- peuvent pourtant porter cet état. Kesh le tolère sans le réparer d'office, et l'écran \emph{Exercices comptables} l'affiche dans un \textbf{bandeau d'avertissement}, visible de tous les rôles~: il nomme l'exercice ouvert le plus ancien, l'exercice clôturé qui le suit, et l'exercice clôturé le plus récent.

Tant que l'ordre n'est pas rétabli, \textbf{rien ne s'enregistre} dans l'exercice ouvert concerné~: ni saisie, ni validation ou règlement de facture, ni paiement fournisseur, ni complément des soldes de départ, ni dévalidation, ni modification, ni suppression. Le bilan étant cumulatif, chacune de ces écritures changerait le bilan de l'exercice clôturé qui le suit. Le refus nomme cet exercice. Une contre-passation, datée du jour, n'est refusée que si l'exercice du jour est lui-même suivi d'un exercice clôturé.

\textbf{Des boutons offerts, puis refusés au clic.} Dans cet état, l'écran ne prévoit pas toujours le refus~: les boutons \emph{Annuler} d'un règlement client, d'un rapprochement bancaire, d'un paiement ou d'une facture fournisseur, et \emph{Contre-passer} sur une écriture, restent affichés, et le complément des soldes de départ reste proposé. Le clic est alors refusé, avec un message qui nomme l'exercice clôturé~: rien n'a été enregistré. Rétablissez l'ordre ci-dessous, puis recommencez (issue \#568).

Pour rétablir l'ordre, \textbf{dans cet ordre}~:

\begin{enumerate}
    \item \textbf{Si les comptes de l'exercice ouvert sont arrêtés}, clôturez-le --- c'est la réparation directe, ouverte au Comptable~: aucun exercice antérieur n'étant ouvert, Kesh l'accepte. S'il reste d'autres exercices ouverts suivis d'un exercice clôturé, clôturez-les à leur tour, \textbf{du plus ancien au plus récent}~; le bandeau disparaît quand l'ordre est rétabli.
    \item \textbf{Sinon}, un administrateur rouvre les exercices clôturés postérieurs, \textbf{en commençant par le plus récent} (celui que le bandeau nomme en dernier) et jusqu'à celui qui suit l'exercice ouvert (\S\ref{sec:reouverture-exercice}). Rouvrir d'abord l'exercice clôturé le plus proche est refusé dès qu'un exercice plus récent est clôturé.
\end{enumerate}

\subsection{Reprise de comptabilité — soldes de départ}\label{sec:soldes-depart}

Si vous migrez vers Kesh depuis un autre logiciel comptable, l'écran \emph{Administration} → \emph{Soldes de départ} (réservé aux rôles Comptable et Administrateur) vous permet de reprendre le bilan de clôture de votre ancienne comptabilité en une seule saisie :

\begin{enumerate}
    \item \textbf{Préparez le bilan de clôture} de votre ancien logiciel (actifs, passifs, capitaux propres, report à-nouveau accumulé).
    \item \textbf{Ouvrez l'écran} : la grille liste vos comptes de bilan actifs et postables (actifs et passifs — jamais les comptes de produits ou de charges, qui fausseraient le résultat de l'exercice).
    \item \textbf{Saisissez les soldes} colonne Débit ou Crédit, ligne par ligne. Le total débit, le total crédit et la différence s'affichent en direct.
    \item \textbf{Équilibrez via le report à-nouveau} : posez votre report accumulé sur le compte de rôle \emph{Bénéfice/perte reporté} (2970 dans les plans standards) — au crédit pour un bénéfice cumulé, au \textbf{débit} pour une perte reportée.
    \item \textbf{Générez} : le bouton \emph{Générer l'écriture d'ouverture} (actif seulement une fois la saisie équilibrée) crée \textbf{une écriture OD équilibrée datée au 1\textsuperscript{er} jour de votre premier exercice}. Votre bilan est juste dès le premier jour.
\end{enumerate}

Sous la grille, Kesh affiche aussi le total des \textbf{Actifs}, celui des \textbf{Passifs et capitaux} (compte de report exclu), et le \textbf{montant à porter au compte de report} pour équilibrer — au crédit s'il est positif, au débit s'il est négatif. Quand la saisie est équilibrée \textbf{sans rien sur le compte de report}, un avertissement le signale~: soit l'écart a été mis sur un autre compte, soit votre report vaut réellement zéro. Il ne bloque pas la génération. Deux variantes le remplacent quand aucun compte en service ne porte le rôle \emph{Bénéfice/perte reporté}, ou quand ce compte n'est pas imputable.

\begin{keshwarning}
\textbf{Le montant à porter aide la saisie, il ne contrôle rien une fois la saisie équilibrée.} Si le report a été mis par erreur sur un autre compte de passif, ce montant est compté dans les passifs, et le montant à porter vaut zéro. \textbf{Le vrai contrôle} est de comparer les totaux Actifs et Passifs affichés à ceux du bilan de votre ancienne comptabilité~: ils doivent être identiques.
\end{keshwarning}

\begin{keshnote}
La génération n'est disponible que tant que la société ne contient \textbf{aucune écriture} (dans aucun exercice). Ensuite, l'écriture d'ouverture se \textbf{modifie} comme toute écriture saisie à la main~: si vous avez inversé débit et crédit, ouvrez sa fiche (\emph{Écritures}, première écriture du premier exercice), \emph{Modifier}, corrigez et validez --- tant que l'exercice est ouvert et hors période verrouillée, et après l'avoir délettrée si l'une de ses lignes est lettrée (ses lignes de bilan sont lettrables~: par exemple le solde d'ouverture d'un compte de passif avec l'écriture, saisie à la main, qui le règle~; les encaissements des factures, eux, se lettrent d'eux-mêmes avec leur facture, une ligne de pièce ne se lettrant pas à la main). L'écran reste en mode \emph{compléter}. Elle se \textbf{supprime} aussi, dans le même cadre~: si elle est la \textbf{seule} écriture de la société, la société redevient vierge et la \textbf{génération est de nouveau proposée} --- la nouvelle écriture d'ouverture reçoit alors le numéro~2, pas le~1, car un numéro n'est jamais réattribué (le trou s'explique par la trace de la suppression au journal d'audit). Si d'autres écritures existent, la génération reste refusée~: les comptes de l'ouverture \textbf{sans autre mouvement} redeviennent complétables, mais un compte \textbf{déjà mouvementé ailleurs} --- typiquement la banque --- ne l'est pas~; son solde de départ se ressaisit à la main, par une écriture d'OD. Un complément se supprime comme l'ouverture. \emph{Limite~: rien ne la distingue plus d'une autre écriture~; rien n'empêche donc d'y ajouter une ligne de résultat, ce que la génération refuse.}
\end{keshnote}

\subsubsection{Compléter un compte oublié}

Une fois l'ouverture générée, l'écran passe en mode \emph{compléter}~: sous le rappel de ce qui a été généré, il propose les comptes de bilan --- actifs et passifs --- \textbf{jamais mouvementés}~: en service, imputables, sans aucune ligne d'écriture dans aucun exercice. Il faut pour cela un compte de rôle \emph{Bénéfice/perte reporté} en service et imputable et, pour une régularisation, un exercice ouvert qui couvre la date du jour~; à défaut, l'écran dit ce qui manque. Saisissez le solde des seuls comptes oubliés, puis \emph{Compléter}~: Kesh crée une écriture OD «~Complément des soldes de départ~», dont il \textbf{calcule lui-même la contrepartie} sur le compte de rôle \emph{Bénéfice/perte reporté}~: elle équilibre l'ensemble des lignes saisies --- au crédit quand leurs débits l'emportent (un actif oublié, saisi au débit), au débit dans le cas contraire (un passif, saisi au crédit). Si les lignes saisies s'équilibrent entre elles, il n'y a pas de contrepartie.

\begin{itemize}
    \item \textbf{La date.} Si votre premier exercice est encore ouvert et que son premier jour n'est pas dans la période verrouillée, le complément est daté de ce premier jour~: le compte oublié faisait partie de l'ouverture, et les soldes d'ouverture --- donc tout rapport déjà tiré sur cet exercice --- en sont modifiés. Si un exercice suivant est déjà clôturé --- ce qui ne se rencontre plus qu'avec des données antérieures à la v0.13.0, les exercices se clôturant désormais dans l'ordre ---, le complément est \textbf{refusé} tant qu'il l'est (\S\ref{sec:exercices-desordre})~: il changerait son bilan reporté. ⚠️ L'écran ne le prévoit pas et propose le complément~; le refus vient à l'enregistrement, et nomme l'exercice clôturé. Sinon, le complément est une \textbf{régularisation datée du jour}, dans l'exercice ouvert qui couvre cette date. La date annoncée par l'écran est indicative~; elle est confirmée à l'enregistrement.
    \item \textbf{Le jour est celui du serveur, en temps universel}~: entre minuit et une ou deux heures du matin, heure suisse, c'est encore la veille --- le 1\textsuperscript{er} janvier, une régularisation peut tomber dans l'exercice précédent, ou être refusée s'il est clôturé.
    \item \textbf{La numérotation.} Daté du premier jour, le complément porte pourtant le numéro suivant du journal~: dans cet exercice, l'ordre des numéros cesse de suivre l'ordre des dates.
\end{itemize}

\begin{keshnote}
\textbf{Un montant faux sur un compte déjà saisi ne se complète pas}~: le compte est mouvementé. On le corrige dans le journal --- d'abord en \textbf{modifiant} l'écriture fautive depuis sa fiche, tant que l'exercice est ouvert (après l'avoir délettrée si l'une de ses lignes est lettrée)~; sinon en la \textbf{contre-passant}, ou par une écriture de correction. Une modification qui \textbf{retire} la ligne d'un compte le rend de nouveau complétable. Une écriture contre-passée garde ses lignes~: ses comptes restent mouvementés. En particulier, \textbf{après la contre-passation de toute l'écriture d'ouverture, plus aucun de ses comptes n'est complétable}~; les soldes justes se saisissent alors par une écriture manuelle. À l'inverse, la \textbf{dévalidation d'une facture} supprime son écriture~: un compte dont c'était le seul mouvement redevient complétable.
\end{keshnote}

\begin{keshtip}
Si votre premier exercice Kesh ne doit pas commencer au 1\textsuperscript{er} janvier de l'année courante (exercice décalé ou reprise antérieure), créez d'abord l'exercice voulu via \emph{Administration} → \emph{Exercices comptables} — l'écriture d'ouverture est datée au premier jour du \textbf{premier} exercice existant. Un exercice antérieur créé \emph{après} la génération devient à son tour le premier~: un complément serait alors daté de son premier jour, et non de celui de l'écriture d'ouverture. (Cette création est refusée si un exercice postérieur est déjà clôturé --- voir la création d'un nouvel exercice.)
\end{keshtip}

%%─────────────────────────────────────────────────────────────────────────────
\section{Carnet d'adresses}

\subsection{Contacts}

\emph{Contacts} ou \emph{Carnet d'adresses}. La page affiche tous vos contacts (clients, fournisseurs, autres).

\subsubsection{Création d'un contact}

\keshscreenshot{contact-nouveau.png}{Création d'un contact client (société) avec IDE et conditions de paiement.}  % TODO capture: form nouveau contact (donnees factices : Boutique Pierre Exemple Sarl)

\emph{Nouveau contact}. Renseigner :

\begin{itemize}
    \item \textbf{Type} : \textbf{Entreprise} (raison sociale) ou \textbf{Personne} (prénom + nom séparés).
    \item \textbf{Nom} : raison sociale pour une entreprise ; \textbf{prénom et nom} pour une personne.
    \item \textbf{IDE} (si société suisse soumise à inscription au RC).
    \item \textbf{Numéro de client} (optionnel) : la référence sous laquelle \emph{vous} identifiez ce client. Elle est \textbf{libre} --- Kesh ne la génère pas et ne l'impose pas ---, doit être unique dans votre société, et figure sur le PDF de facture (voir \ref{sec:numero-client}).
    \item \textbf{Adresse structurée} : rue, numéro, NPA, localité, pays en \emph{champs séparés}. Ce format « structuré » (type S) est celui exigé par les banques suisses sur les QR-factures — le NPA et la localité sont requis pour générer une facture. Le NPA reste libre (adresses étrangères possibles).
    \item \textbf{Email}, \textbf{Téléphone} (optionnels).
    \item \textbf{Compte comptable par défaut} : par exemple 1100 (Débiteurs) pour un client.
    \item \textbf{Conditions de paiement} : par exemple « 30 jours net ».
\end{itemize}

\begin{keshnote}
\textbf{Pendant que vous tapez, Kesh vous signale les contacts qui pourraient déjà
exister.} Deux avertissements distincts, qui n'ont pas la même force :

\begin{itemize}
    \item \textbf{Un nom proche} --- Kesh liste les contacts déjà enregistrés qui
    pourraient correspondre, avec de quoi les reconnaître : localité, numéro de
    client, à défaut l'email. C'est une \emph{suggestion} : deux clients peuvent
    légitimement porter des noms très proches --- deux sociétés d'un même groupe,
    deux homonymes, un père et son fils.
    \item \textbf{Un numéro IDE déjà pris} --- l'avertissement est franc, car un
    IDE est un identifiant attribué par l'État : deux entités ne le partagent
    jamais. Si le contact qui le porte est \textbf{archivé}, Kesh le dit : un IDE
    reste réservé même après archivage, et l'enregistrement sera refusé.
\end{itemize}

\textbf{Ces avertissements n'empêchent jamais d'enregistrer.} Vous restez juge :
si vous savez qu'il s'agit bien d'un autre contact, enregistrez. Kesh vous
informe, il ne décide pas à votre place.

\textbf{Ce que cet avertissement ne garantit pas.} Kesh examine les vingt
premiers contacts que la recherche remonte, et vous en montre les cinq plus
proches. Lorsqu'un très grand nombre de contacts partagent un mot du nom que
vous saisissez --- une vingtaine de sociétés en \emph{Sàrl}, par exemple ---
le doublon que vous cherchiez à éviter peut se trouver \emph{au-delà} de ces
vingt, et l'avertissement ne le mentionnera pas. \textbf{L'absence
d'avertissement ne prouve donc pas l'absence de doublon.} Dans ce cas, la
recherche du carnet, elle, les trouve tous.
\end{keshnote}

\begin{keshnote}
Pour un contact de type \textbf{Entreprise}, vous pouvez ajouter une ou plusieurs
\textbf{personnes de contact} (interlocuteurs : prénom, nom, fonction, email,
téléphone) depuis sa fiche. Ces personnes sont \emph{purement informatives} —
elles n'apparaissent jamais sur les factures ni les paiements.
\end{keshnote}

\subsubsection{Recherche de contacts}

La barre de recherche utilise l'index FULLTEXT MariaDB sur le \textbf{nom}, complété par une recherche sur l'\textbf{email}, le \textbf{numéro de client} et le \textbf{numéro IDE} --- y compris sur un fragment. Le wildcard \texttt{*} est supporté en suffixe du nom.

C'est ce qui permet de remonter d'une facture papier au contact : saisissez le numéro de client imprimé sur le document.

\begin{keshwarning}
\textbf{Le numéro IDE se cherche sans ses séparateurs.} Kesh le conserve sous sa
forme compacte, \texttt{CHE109322551} : c'est donc celle-là qu'il faut saisir
dans la recherche. La forme imprimée sur les documents officiels,
\texttt{CHE-109.322.551}, ne remonte \emph{rien}.
\end{keshwarning}

\subsection{Catalogue produits}

\emph{Produits}. Les produits sont les biens ou services que vous vendez régulièrement et que vous voulez réutiliser dans vos factures.

\subsubsection{Création d'un produit}

\keshscreenshot{produit-nouveau.png}{Création d'un produit/service du catalogue avec prix unitaire HT et compte produit par défaut.}  % TODO capture: form nouveau produit

\emph{Nouveau produit}. Renseigner :

\begin{itemize}
    \item \textbf{Code interne} : identifiant unique (ex. PROD-001).
    \item \textbf{Désignation} : libellé qui apparaîtra sur les factures.
    \item \textbf{Prix unitaire HT}.
    \item \textbf{Taux TVA} : 0\%, 2.6\%, 3.8\%, 8.1\% --- les taux suisses en vigueur depuis le 1\textsuperscript{er} janvier 2024, que Kesh pose à la création de la société (\emph{Paramètres} → \emph{Taux de TVA}).
    \item \textbf{Unité} : pièce, heure, mètre, kilo, etc.
    \item \textbf{Compte de produit} (facultatif) : le compte du plan comptable sur lequel le chiffre d'affaires de cet article sera crédité — par exemple 3200 \emph{Prestations de services} pour du conseil, ou un sous-compte 3200.1 \emph{Maintenance}, à créer d'abord dans le plan, pour un abonnement. Quand vous montez une facture depuis le catalogue, la ligne créée reprend ce compte automatiquement ; vous pouvez encore le changer sur la ligne. \textbf{Laissé vide}, l'article n'impose rien : la ligne suit alors le compte de produit par défaut de la société.
\end{itemize}

\subsection{Conditions de paiement}

\emph{Conditions de paiement}. Définir les conditions standard utilisées dans vos factures :

\begin{itemize}
    \item « Comptant » (0 jour).
    \item « 30 jours net ».
    \item « 60 jours net ».
    \item « 10 jours 2\% escompte sinon 30 jours net ».
    \item Personnalisées selon vos besoins.
\end{itemize}

%%─────────────────────────────────────────────────────────────────────────────
\section{Facturation QR Bill}

\subsection{Création d'une facture}

\emph{Factures} → \emph{Nouvelle facture}. Renseigner :

\keshscreenshot{facture-nouvelle.png}{Création d'une facture multi-lignes avec sélection du client et du compte bancaire de réception.}  % TODO capture: form nouvelle facture

\begin{itemize}
    \item \textbf{Client} : sélection dans le carnet d'adresses (ou création à la volée).
    \item \textbf{Date de facture} : par défaut aujourd'hui.
    \item \textbf{Date d'échéance} : calculée automatiquement selon les conditions de paiement du client (modifiable).
    \item \textbf{Numéro} : généré automatiquement selon le pattern configuré (cf. \emph{Numérotation}).
    \item \textbf{Lignes de facture} : un ou plusieurs produits avec quantité, prix, TVA, et \textbf{compte de produit}.
    \item \textbf{Compte de produit (par ligne)} : le compte du plan comptable sur lequel le chiffre d'affaires de cette ligne sera crédité. Quand vous ajoutez une ligne \textbf{depuis le catalogue}, Kesh y \textbf{recopie} le compte de produit de l'article, s'il en porte un ; vous restez libre de le changer ensuite. Laissé vide, le champ suit le \textbf{compte de produit par défaut de la société} (\emph{Paramètres} → \emph{Facturation}) — il affiche alors ce compte entre parenthèses, pour que vous voyiez d'un coup d'œil où ira le produit. \textbf{Attention} : la recopie a lieu \textbf{une seule fois}, au moment où vous choisissez l'article ; la ligne ne reste pas liée à lui. Vider le champ d'une ligne venue du catalogue la ramène donc au compte par défaut de la société, \textbf{pas} au compte de l'article. Renseigné, il permet de \textbf{ventiler une même facture sur plusieurs comptes} : honoraires, marchandises, prestations, etc. Le compte reste \textbf{modifiable tant que la facture est un brouillon} ; il est figé à la validation.
    \item \textbf{Remarques} : texte libre apparaissant en bas de facture (mention de remerciement, conditions particulières, etc.).
    \item \textbf{Compte bancaire de réception} : le compte sur lequel vous attendez le paiement (utilisé pour générer la référence QR Bill).
\end{itemize}

\subsubsection{Exemple : facture d'une agence web}

\emph{ABC Web Sàrl} facture une prestation à son client \emph{Boutique Pierre Exemple Sàrl} :

\begin{itemize}
    \item Client : Boutique Pierre Exemple Sàrl (sélectionné dans le carnet d'adresses, conditions « 30 jours net »).
    \item Ligne 1 : « Refonte site web » — 1 × CHF 4'500.00.
    \item Ligne 2 : « Maintenance mensuelle » — 3 × CHF 150.00 = CHF 450.00.
    \item Total : CHF 4'950.00. Échéance calculée à 30 jours.
    \item Compte de réception : le compte bancaire QR-IBAN de l'agence.
\end{itemize}

À la validation, Kesh attribue le numéro \texttt{F-2026-0042} et crée l'écriture comptable automatique (Débit 1100 Débiteurs / Crédit 3200 Prestations de services, désigné par l'agence comme compte de produit par défaut). Le PDF QR Bill est produit et \textbf{figé} au premier téléchargement ou au premier envoi (cf. \S\ref{sec:pdf-fige}).

\medskip
\noindent\emph{Variante ventilée.} Si l'agence souhaite distinguer ses activités, elle crée d'abord dans son plan le sous-compte \texttt{3200.1 Maintenance}, puis affecte la ligne 1 au compte \texttt{3200 Prestations de services} et la ligne 2 au compte \texttt{3200.1 Maintenance}. L'écriture générée porte alors \textbf{deux lignes de crédit distinctes} — Crédit 3200 pour CHF 4'500.00 et Crédit 3200.1 pour CHF 450.00 — au lieu d'un crédit unique. Le compte de résultat devient exploitable sans reclassement manuel a posteriori. Les lignes laissées vides suivent le compte de produit par défaut de la société et s'y regroupent — y compris celles qui viennent du catalogue, dont le compte a été vidé après la recopie. Si les lignes de la facture portent un taux de TVA, Kesh comptabilise aussi automatiquement la TVA due au compte de TVA (p.~ex. 2200) ; ici l'agence n'est pas assujettie, l'écriture se limite donc à CHF 4'950.00 (cf. \emph{Décompte TVA}).

\subsection{Validation d'une facture}

Une facture passe par trois états :

\begin{enumerate}
    \item \textbf{Brouillon} : modification libre, n'apparaît pas dans les rapports. Un brouillon n'a en général \emph{pas} de numéro --- sauf s'il vient d'être \textbf{dévalidé}, auquel cas il \textbf{conserve} celui qu'il portait, et le reprendra à la revalidation (cf. \S\ref{sec:suppression-facture}).
    \item \textbf{Validée} : numéro définitif attribué, PDF QR Bill disponible (figé au premier téléchargement ou envoi), écriture comptable créée automatiquement (Débit 1100 Débiteurs / Crédit 3000 Ventes de marchandises + 2200 TVA due), et \textbf{comptes de produit figés}. Une facture validée n'est plus modifiable : le compte de chaque ligne reste néanmoins \textbf{consultable} dans la colonne « Compte de produit » de l'écran de détail, où les lignes suivant le compte par défaut de la société l'affichent avec la mention \emph{(défaut)}. C'est le moyen de contrôler la ventilation sans avoir à ouvrir l'écriture comptable et à la relire à l'envers.
    \item \textbf{Payée} : \textbf{soldée} par un ou plusieurs règlements --- enregistrés à la main
          depuis la fiche, ou par le rapprochement d'une transaction bancaire entrante
          (\S\ref{sec:reglement-client}). Une facture réglée en partie reste \emph{validée}, avec
          son reste dû.
\end{enumerate}

\paragraph{Le total arrondi à 5 centimes.} Comme le veut l'usage suisse, le total d'une facture est \textbf{arrondi à 5~centimes} à sa validation : 123.44 devient 123.45, 234.52 devient 234.50. La fiche de la facture et son PDF portent une ligne \textbf{« Arrondi »} entre la TVA et le total, avec l'écart signé (+0.01, -0.02) ; le bulletin de versement QR, lui, porte directement le montant arrondi ; l'écart, sans TVA, passe sur le compte de différences d'arrondi. Sur un \textbf{brouillon}, la fiche montre déjà la ligne, en \emph{« Arrondi (estimé) »}, avec un total estimé : l'arrondi n'est figé qu'à la validation. Une facture validée garde son arrondi, même si le réglage change ensuite ; un avoir le reprend et l'annule. Le réglage se trouve dans \emph{Paramètres} → \emph{Facturation}, case \emph{« Arrondir le total des factures émises à 5~centimes »} --- cochée par défaut. Sans compte de différences d'arrondi désigné, une facture dont le total n'est pas déjà rond ne peut pas être validée.

\paragraph{Le montant minimum d'une facture.} Une facture d'un très petit montant coûte souvent plus à traiter qu'elle ne rapporte. Dans \emph{Paramètres} → \emph{Facturation}, section \emph{Montant minimum}, vous pouvez fixer un seuil : une facture dont le total --- arrondi, celui que le client paiera --- lui est \textbf{inférieur} ne peut pas être validée, et le message donne les deux montants. Le brouillon reste enregistrable. Un total égal au seuil est accepté. Laissé vide (c'est le défaut), aucun seuil ne s'applique. Les \textbf{avoirs} ne sont pas concernés : ils annulent une facture déjà émise, quel que soit son montant.

\paragraph{Un compte des réglages devenu non imputable.} La validation écrit sur la \textbf{créance} et, si la facture porte de la TVA, sur la \textbf{TVA due} désignées dans \emph{Paramètres} → \emph{Facturation}. Si l'un de ces comptes est devenu non imputable depuis sa désignation (sous-comptes créés, case \emph{imputable} décochée), la validation est refusée avec un message qui \textbf{nomme le compte} et renvoie à ces paramètres~; la facture reste en brouillon. Un administrateur y désigne un compte imputable, puis la validation se relance telle quelle. L'\textbf{avoir}, lui, n'est \textbf{pas} soumis à ce contrôle, et c'est voulu~: une facture émise doit rester annulable, même si sa créance ou sa TVA due est devenue non imputable depuis (voir «~Rôles des comptes~»).

\subsection{Vos coordonnées sur la facture}\label{sec:coordonnees-emetteur}

Le bloc en haut à gauche du PDF porte votre \textbf{nom}, votre \textbf{adresse} et votre \textbf{numéro IDE}. Vous pouvez y ajouter vos \textbf{coordonnées de contact} --- téléphone, e-mail et site web --- pour qu'un client qui a une question sur la facture puisse vous joindre depuis le document lui-même.

\textbf{Où les saisir} : \emph{Paramètres} $\rightarrow$ section \emph{Organisation}, bouton \textbf{Modifier} en face de « Téléphone · Site web ». La saisie est réservée aux administrateurs, comme celle de l'adresse de réponse.

\begin{keshnote}
\textbf{Chaque champ est facultatif, et un champ vide ne laisse aucun blanc sur le document} : la ligne correspondante n'est simplement pas imprimée. Si vous ne renseignez rien, votre facture est exactement celle que vous aviez avant --- aucune mise en page ne bouge.
\end{keshnote}

\textbf{L'e-mail que vous avez peut-être déjà renseigné} comme adresse de réponse des envois apparaît désormais lui aussi sur le PDF. Si vous ne le souhaitez pas, il suffit de le vider --- mais notez qu'il sert aussi de \emph{Reply-To} aux factures envoyées par e-mail (cf. \S\ref{sec:envoi-email}).

Les \textbf{avoirs} portent les mêmes coordonnées, à l'identique.

\textbf{Une facture déjà émise ne change plus} : vos coordonnées s'impriment sur les factures dont le PDF n'est pas encore figé. Celui d'une facture déjà téléchargée ou envoyée reste tel qu'il a été émis (cf. \S\ref{sec:pdf-fige}).

\begin{keshwarning}
\textbf{Une coordonnée très longue est raccourcie à l'affichage} (suivie de «~\ldots~») pour ne pas empiéter sur le bloc de droite, qui porte le numéro et la date de la facture. La valeur complète reste intacte dans vos réglages.

\textbf{Et si l'en-tête entier ne tient plus sur la page, le document est refusé} --- facture comme avoir --- avec un message explicite. C'est délibéré : l'en-tête et le tableau des lignes se partagent une hauteur fixe, et mieux vaut un refus qu'un document où l'un s'imprime par-dessus l'autre. La cause peut être vos coordonnées --- mais aussi une adresse de destinataire comptant beaucoup de lignes.

Le message vous invite à \textbf{supprimer} une coordonnée, et non à la raccourcir : chaque coordonnée occupe une ligne entière quelle que soit sa longueur, si bien que l'abréger ne libère aucune place. Vider le champ, en revanche, fait disparaître la ligne.
\end{keshwarning}

\subsection{Le numéro de client sur la facture}\label{sec:numero-client}

Le bloc en haut à \textbf{droite} du PDF porte le titre du document, son numéro et sa date. Vous pouvez y faire figurer le \textbf{numéro de client} --- la référence sous laquelle vous identifiez ce client, ou celle que \emph{lui} vous impose.

Elle rend deux services : votre client rapproche la facture de son propre dossier fournisseur, et vous-même remontez d'une facture papier au contact en saisissant ce numéro dans la \textbf{recherche} du carnet d'adresses.

\textbf{Où le saisir} : sur la \textbf{fiche du contact}, champ \emph{Numéro de client}, sous le numéro IDE. La valeur apparaît ensuite sur toutes les factures adressées à ce contact dont le PDF n'est pas encore figé, \textbf{avoirs compris} --- une facture déjà émise ne change plus (cf. \S\ref{sec:pdf-fige}).

\begin{keshnote}
\textbf{Kesh ne génère aucun numéro de client}, et n'impose aucun format : \texttt{42}, \texttt{CLI-2026-00042} ou une référence dictée par votre client conviennent également. Le champ est facultatif ; un contact sans numéro n'affiche simplement aucune ligne sur le PDF, et rien ne bouge dans la mise en page.
\end{keshnote}

\begin{keshnote}
Le numéro doit être \textbf{unique parmi vos contacts actifs} --- c'est ce qui lui permet d'identifier. Si vous en saisissez un déjà pris, Kesh vous le signale et refuse l'enregistrement. \textbf{Deux numéros identiques à l'œil sont le même numéro} : la casse ne distingue pas (\texttt{CLI-1} et \texttt{cli-1}), la façon dont un accent a été composé par votre système non plus, et les caractères invisibles qu'un copier-coller depuis un courriel ou un document Word peut coller au passage sont ignorés. En revanche, \textbf{un accent distingue} : \texttt{CLI-É1} et \texttt{CLI-E1} restent deux numéros différents --- deux clients peuvent légitimement ne différer que par là. Ce que vous avez tapé est conservé tel quel : c'est ce qui s'affiche sur la fiche et s'imprime sur la facture.

En revanche, \textbf{archiver un contact libère son numéro} : vous pouvez le réattribuer à un nouveau contact. L'ancien contact, lui, \textbf{conserve le sien} --- ses factures passées restent ainsi lisibles. Si vous réattribuez un numéro puis cochez «~\emph{Inclure archivés}~» dans la recherche, ce numéro remonte donc \textbf{deux} contacts : l'actif et l'archivé. C'est voulu, mais gardez-le en tête avant de recycler un numéro encore présent sur des factures récentes.
\end{keshnote}

\begin{keshwarning}
\textbf{Un numéro très long est raccourci à l'impression} (suivi de «~\ldots~») : le bloc de droite ne dispose que d'environ 70 mm, soit une trentaine de caractères \emph{libellé compris} --- donc une vingtaine pour la valeur elle-même. La valeur complète reste intacte sur la fiche du contact et dans la recherche. En pratique, une référence de la longueur de \texttt{CLI-2026-00042} s'imprime entière.
\end{keshwarning}

\subsection{Numérotation des factures}

Le numéro de facture est généré automatiquement selon un pattern configurable dans \emph{Paramètres} → \emph{Facturation} :

\begin{itemize}
    \item Pattern par défaut : \texttt{F-AAAA-NNNN} (ex. F-2026-0042).
    \item Patterns alternatifs : \texttt{INV-AAAAMM-NNN}, \texttt{F\{company\_id\}-AAAA-NNNN}, etc.
    \item La numérotation est séquentielle dans un exercice (CO Art. 957a al. 2 ch. 5 « traçabilité ») : les numéros sont attribués dans l'ordre et \textbf{jamais réattribués}. Elle peut en revanche comporter des \textbf{trous} --- supprimer un brouillon déjà numéroté en laisse un, le compteur ne redescendant pas. Un trou se voit et s'explique ; c'est la réattribution qui serait fautive, et elle n'a pas lieu.
    \item Deux validations simultanées reçoivent toujours deux numéros distincts : le compteur de numérotation est \textbf{verrouillé} le temps de la validation, et une contrainte d'unicité de la base interdit d'attribuer deux fois le même numéro. Une validation qui croise une autre opération (un rapprochement bancaire, par exemple) est \textbf{rejouée automatiquement} ; si le conflit se répète trois fois de suite, la validation échoue sans rien enregistrer et peut simplement être relancée.
\end{itemize}

\subsection{Génération du PDF avec QR Bill}

\keshscreenshot{facture-qrbill.png}{Exemple de PDF de facture avec section de paiement QR Bill 2.2 détachable en bas de page.}  % TODO capture: PDF facture QR Bill (donnees factices)

Le PDF d'une facture validée est conforme au standard \href{https://www.six-group.com/en/products-services/banking-services/standardization/qr-bill.html}{QR Bill 2.2} (norme suisse 2020+) :

\begin{itemize}
    \item Format A4, page principale + section de paiement séparable en bas.
    \item QR code suisse contenant toutes les informations de paiement (créancier, débiteur, montant, IBAN, référence structurée, communication).
    \item Facture et section de paiement dans la \textbf{langue de correspondance du client} (fiche contact --- sinon la langue de l'instance).
    \item Compatible scan par les e-banking suisses et applications mobiles (Raiffeisen, PostFinance, UBS, Credit Suisse, Migros Bank, etc.).
\end{itemize}

Le PDF est téléchargeable depuis la facture validée — pour l'imprimer (envoi postal) ou l'archiver. Il est \textbf{identique d'un téléchargement à l'autre} : c'est le document émis, figé (cf. \S\ref{sec:pdf-fige}). Depuis la version v0.6, Kesh peut aussi l'\textbf{envoyer directement par e-mail} à votre client, sans passer par votre messagerie : voir la sous-section suivante.

\begin{keshtip}
La référence structurée du QR Bill est calculée selon la norme ISO 11649 (RF Reference) ou le format suisse QRR (Référence Kesh). Le choix du format dépend du compte bancaire de réception : les comptes IBAN suisses standard utilisent IS0 11649, les comptes QR-IBAN utilisent QRR.
\end{keshtip}

\subsubsection{Le PDF émis est figé}\label{sec:pdf-fige}

Le PDF d'une facture est produit \textbf{une seule fois} : au premier téléchargement ou au premier envoi par e-mail qui suit la validation, dans la langue de correspondance du client. Kesh le conserve alors tel quel, avec son \textbf{empreinte} (SHA-256), et c'est ce même fichier, octet pour octet, que rendent ensuite chaque téléchargement et chaque envoi. La fiche facture l'indique par la mention \emph{Document figé le} suivie de la date.

\begin{itemize}
    \item \textbf{Modifier ensuite} vos coordonnées, l'adresse, le numéro ou la langue du client, ou la mise en page, \textbf{ne change pas} une facture déjà émise : seules les factures non encore figées en tiennent compte. C'est ce qui permet de reproduire, des années plus tard, la copie exacte de la pièce envoyée.
    \item \textbf{Qui que ce soit} qui télécharge le premier --- y compris un utilisateur en consultation ou une clé d'API en lecture --- fige le document : une pièce sortie de Kesh est figée. Le gel est inscrit au \textbf{journal d'audit}, avec l'empreinte et le nom de qui l'a déclenché.
    \item Une facture \textbf{annulée par un avoir} garde le PDF qu'elle a émis, et la fiche continue de le proposer. Ce PDF ne porte pas de mention d'annulation : c'est l'avoir qui la porte. Une facture annulée dont le PDF n'a \textbf{jamais été figé} --- jamais téléchargée ni envoyée, ou envoyée seulement avant la v0.12.1 --- n'a pas de PDF.
    \item Une facture validée \textbf{avant la version v0.12.1} est figée à son prochain téléchargement, avec les données du moment : ses PDF antérieurs n'avaient pas été conservés. ⚠️ Si elle avait déjà été envoyée, le document figé \textbf{n'est pas forcément celui que le client a reçu} --- la date de la mention \emph{Document figé le} est alors postérieure à celle de l'envoi.
    \item La \textbf{restauration d'une sauvegarde} antérieure au gel ramène la facture à l'état non figé : son prochain téléchargement figera un \textbf{nouveau} document. L'ancien reste dans le répertoire des documents, et son empreinte au journal d'audit.
    \item La \textbf{dévalidation} détache le document figé (son empreinte reste au journal d'audit) ; la revalidation en figera un nouveau au premier téléchargement.
    \item Les \textbf{rappels}, eux, sont produits à chaque envoi (cf. \S\ref{sec:relances-user}) : seul le PDF de la facture est figé.
\end{itemize}

\begin{keshwarning}
\textbf{Si le fichier figé a disparu} --- typiquement après une restauration qui n'a pas repris le répertoire des documents ---, le téléchargement est \textbf{refusé} avec un message qui nomme le fichier manquant (\texttt{<empreinte>.pdf}). Kesh ne le régénère jamais de lui-même : ce serait fabriquer une pièce qui n'a pas été émise. Le remède est de \textbf{restaurer ce fichier} depuis la sauvegarde du répertoire des documents --- c'est le seul moyen de retrouver l'original. À défaut, un \textbf{administrateur} voit apparaître sur la fiche, après ce refus, le bouton \emph{Refiger le document} : il produit un \textbf{nouveau} document avec les données et la langue du client d'aujourd'hui, \textbf{qui n'est pas l'original}, et le journal d'audit garde l'empreinte de l'ancien et du nouveau. Une facture annulée ne peut pas être refigée : seule la restauration du fichier la répare.
\end{keshwarning}

\subsection{Envoyer une facture par e-mail}\label{sec:envoi-email}

\keshscreenshot{facture-envoi-email.png}{Modale d'envoi : destinataire verrouillé, objet et message pré-remplis dans la langue du client, modifiables avant l'envoi.}  % TODO capture: modale envoi e-mail (facture validee, contact avec e-mail)

Sur une facture \textbf{validée}, le bouton \emph{Envoyer par e-mail} ouvre
une fenêtre pré-remplie :

\begin{itemize}
    \item \textbf{Destinataire} : l'adresse e-mail de la fiche contact,
          affichée mais \textbf{non modifiable} (sécurité). Pour envoyer à
          une autre adresse, modifiez d'abord la fiche contact
          (\emph{Contacts}).
    \item \textbf{Objet et message} : pré-remplis dans la \textbf{langue de
          correspondance du client} (fiche contact — sinon la langue de
          l'instance), avec une formule d'appel adaptée à sa civilité
          (« Chère Madame Muster », « Sehr geehrter Herr… »).
          Les deux restent modifiables avant l'envoi : ce que vous voyez est
          ce qui part.
    \item \textbf{Pièce jointe} : la QR-facture PDF, attachée
          automatiquement --- le document \textbf{figé}, octet pour octet
          celui du téléchargement (\S\ref{sec:pdf-fige}). Le premier envoi
          d'une facture jamais téléchargée la fige, même si l'envoi échoue.
          Si la langue du client a changé depuis le gel, le message suit la
          nouvelle langue et la pièce jointe garde la sienne.
\end{itemize}

Après l'envoi, la fiche facture affiche \emph{Envoyée le} suivi de la date
et de l'adresse du destinataire. Le bouton devient \emph{Renvoyer par
e-mail} — le renvoi est autorisé (nouvelle date, nouvelle entrée au journal d'audit).

\begin{keshnote}
« Envoyée » signifie que le message a été \textbf{remis au serveur
d'envoi} — ce n'est pas un accusé de réception. Si l'adresse est invalide ou
la boîte pleine, le rebond n'apparaît pas dans Kesh. En cas d'échec
\emph{immédiat} (serveur d'envoi injoignable), la facture n'est pas marquée
envoyée et un message d'erreur s'affiche : vous pouvez réessayer.
\end{keshnote}

\begin{keshtip}
Bouton grisé avec une infobulle ? L'envoi d'e-mails n'est pas configuré sur
votre instance (variables \texttt{KESH\_SMTP\_*}) — voyez avec votre
administrateur (manuel administrateur, section « Envoi de factures par
e-mail »).
\end{keshtip}

\subsubsection{Personnaliser les modèles d'e-mail}

\begin{sloppypar}
\emph{Paramètres → Modèles d'e-mail} (réservé à l'administrateur) : l'objet
et le corps du message se personnalisent par langue (onglets FR/DE/IT/EN),
avec des variables remplacées automatiquement à l'envoi
(\texttt{\{salutation\}}, \texttt{\{contactName\}},
\texttt{\{invoiceNumber\}}, \texttt{\{amount\}}, \texttt{\{dueDate\}},
\texttt{\{companyName\}} — la liste est affichée à côté de l'éditeur). Les
modèles de \textbf{rappel} ont en plus : \texttt{\{reminderLevel\}} (le niveau),
\texttt{\{reminderFee\}} (les frais du seul niveau envoyé),
\texttt{\{totalDue\}} (le \textbf{reste dû} de la facture \textbf{plus} les frais
de rappel cumulés), \texttt{\{daysOverdue\}} (les jours de retard) et
\texttt{\{feeNotice\}} (une phrase qui dit les frais compris dans
\texttt{\{totalDue\}}, \textbf{vide} quand il n'y en a pas) ;
\texttt{\{amount\}} y reste le montant total de la facture. Un
modèle par défaut est fourni pour chaque langue ; \emph{Restaurer le défaut}
y revient à tout moment. Sans aucune personnalisation, les envois utilisent
les modèles par défaut — rien à configurer.
\end{sloppypar}

\subsubsection{Langue et civilité du contact, adresse de réponse}

\begin{itemize}
    \item \textbf{Fiche contact} : le champ \emph{Langue de correspondance}
          (héritée de l'instance, ou FR/DE/IT/EN) détermine la langue de
          l'e-mail \textbf{et} celle du PDF de la facture --- téléchargé ou
          joint ---, figée au premier rendu (\S\ref{sec:pdf-fige}) ; le champ \emph{Civilité}
          (Neutre/Monsieur/Madame, contacts de type Personne) personnalise la
          formule d'appel. Une entreprise reçoit toujours la formule neutre.
    \item \textbf{Adresse de réponse} : \emph{Paramètres → Organisation →
          E-mail (adresse de réponse)} (administrateur) — les réponses de vos
          clients arrivent alors dans votre boîte habituelle.
\end{itemize}

\subsection{Échéancier des factures}

\emph{Factures} → \emph{Échéancier}. Vue d'ensemble des factures :

\keshscreenshot{facture-echeancier.png}{Échéancier des factures : à échoir, échues impayées, et payées.}  % TODO capture: page echeancier

\begin{itemize}
    \item Factures à échoir (dans les 30 prochains jours).
    \item Factures échues impayées (en retard).
    \item Factures payées (historique).
\end{itemize}

Chaque ligne montre le \textbf{total} de la facture (TTC) et son \textbf{reste dû}, qui tient compte des
règlements partiels ; une facture réglée en partie porte le statut «~Partiellement payée~». Les
totaux du haut de page additionnent les restes dus. Le bouton \emph{Régler} ouvre le dialogue de
règlement \textbf{pré-rempli au reste dû} --- modifiable pour un nouvel acompte. L'export CSV porte
la même colonne «~Reste dû~».

Pour relancer les factures en retard, Kesh dispose d'un \textbf{écran de rappels} dédié — voir \S\ref{sec:relances-user}. Une facture suspendue de rappels reste visible dans l'échéancier (rien ne se cache).

\subsection{Enregistrer et annuler un règlement}\label{sec:reglement-client}

\textbf{Enregistrer.} Sur la fiche d'une facture validée, \emph{Enregistrer un règlement}~: le
mode (virement sur un compte bancaire, ou espèces et un autre compte de contrepartie), le montant
et la date. Kesh passe l'écriture \emph{Débit contrepartie / Crédit débiteurs} et calcule ce qui
reste dû. La contrepartie \textbf{ne peut pas être le compte débiteurs de la facture}~: l'écriture
serait nulle, et le reste dû baisserait sans que rien ne bouge au grand livre. Kesh le refuse dans
tous les cas~--- de même qu'un compte bancaire lié au compte débiteurs~--- et ne le propose pas dans
les listes, \textbf{à condition} qu'il porte, au plan comptable, le \textbf{rôle} de compte de
créances clients~: un compte qui l'a perdu reste proposé, et c'est l'enregistrement qui est
refusé. Si le compte en cause
est celui que vous \textbf{choisissez} (compte interne, compte bancaire), choisissez-en un autre
ou reliez le compte bancaire à son propre compte de banque~; si c'est un compte \textbf{désigné
dans les réglages} --- le compte de différences d'arrondi ---, un administrateur doit en désigner
un autre dans \emph{Paramètres} → \emph{Facturation}.
Un règlement \textbf{partiel} est permis~: la facture reste validée, avec son reste dû, et
redevient «~payée~» quand le solde atteint zéro. Un montant supérieur au reste dû est
\textbf{refusé}, de même qu'un montant à plus de deux décimales.

\textbf{Au centime.} Une facture dont le total porte plus de deux décimales (quantité ou prix
fractionnaire) se règle au \textbf{reste dû arrondi au centime}~: c'est le montant que le
dialogue propose. Payer 10.01 une facture de 10.0050 la \textbf{solde}~: l'écriture porte alors
une \textbf{troisième ligne}, l'écart d'un demi-centime au plus, sur le compte de différences
d'arrondi désigné dans \emph{Paramètres} → \emph{Facturation}. Sans ce compte, un tel paiement est
refusé --- un acompte, lui, passe toujours. Le rapprochement bancaire enregistre, lui, ses règlements tout seul.

\textbf{La liste des règlements.} La fiche liste chaque règlement --- date, montant, mode --- avec
un lien vers son écriture.

\textbf{Solder le reste.} Quand le client ne paiera pas ce qui reste --- il a déduit un escompte,
sa banque a retenu des frais, ou il ne paiera plus rien ---, le bouton \emph{Solder le reste} de la
fiche (comptable ou administrateur) éteint \textbf{tout} le reste dû et le passe au compte de la
\textbf{nature} choisie~: \emph{escompte accordé}, \emph{frais bancaires}, \emph{perte sur
débiteur}, ou \emph{reste d'arrondi}. Le montant n'est pas saisi~: c'est le reste dû, exactement
--- affiché aux quatre décimales s'il porte une fraction de centime. L'escompte et la perte \textbf{corrigent la TVA due} au prorata des taux de la facture --- le
décompte TVA de la période du solde la retranche (\S\ref{sec:decompte-tva})~; les
frais et l'arrondi, non. Chaque nature a son compte, choisi par un administrateur dans
\emph{Paramètres} → \emph{Facturation}, section \emph{Solde du reste}~; sans lui, le dialogue le
signale. Le \emph{reste d'arrondi} n'est proposé que pour un reste de moins de 5~centimes~--- c'est
ainsi que se solde un reste trop petit pour être payé, inférieur au demi-centime. Si le reste porte
une fraction de centime, elle va au compte de différences d'arrondi. La facture est alors
«~payée~»~; la liste des règlements montre le solde («~Solde --- escompte accordé~»), et la fiche le
distingue de ce qui a été encaissé (\emph{Déjà réglé} / \emph{Soldé}). Un solde est
\textbf{refusé}, sans rien écrire, si le compte de la nature, le compte de différences d'arrondi
ou le compte de TVA due est le \textbf{compte débiteurs de la facture} (un compte débiteurs
changé de type depuis, puis désigné)~: un administrateur doit en désigner un autre dans
\emph{Paramètres} → \emph{Facturation}. Un solde s'\textbf{annule}
comme un règlement --- bouton \emph{Annuler le solde}, contre-passation datée du jour, TVA
comprise.

\textbf{Annuler.} Un règlement enregistré par erreur ne se supprime pas~: il s'\textbf{annule par
contre-passation}. Le bouton \emph{Annuler le règlement} de sa ligne (comptable ou
administrateur) passe une écriture inverse \textbf{datée du jour}, retire le règlement de la
facture, et le montant redevient dû --- la facture repasse «~validée~» si plus rien ne la solde.
L'écriture d'origine et sa contre-passation restent toutes deux au grand livre, et le journal
d'audit nomme le geste.

Quand l'annulation n'est pas possible, le bouton est remplacé par la raison~:

\begin{itemize}
    \item \textbf{la facture a été créditée} par un avoir~: le règlement reste ouvert au compte
          débiteurs --- il ne s'annule pas. Ce cas ne se présente plus
          que sur une facture créditée \emph{avant} la version 0.12.1, ou restaurée d'une sauvegarde
          antérieure~: un avoir est désormais refusé sur une facture réglée (\S\ref{sec:avoirs})~;
    \item \textbf{le reste de la facture a été soldé}~: annulez d'abord le \textbf{solde} --- lui
          reste annulable --- puis le règlement~;
    \item \textbf{l'exercice du règlement est clôturé}~: un administrateur doit
          d'abord \textbf{rouvrir} les exercices clôturés jusqu'à celui du règlement, en
          commençant par le plus récent~;
    \item \textbf{le lettrage de la facture est figé par une période close}~: ce règlement est
          lettré avec sa facture, et toutes les lignes de ce lettrage sont dans une période
          close. Pour pouvoir l'annuler, prenez la date la plus récente du lettrage (en général
          celle du dernier règlement)~: si elle est sous le verrou de période, un administrateur
          doit faire reculer le verrou avant elle~; et si son exercice est suivi d'un exercice
          clôturé, ou clôturé lui-même, il doit rouvrir les exercices clôturés jusqu'à
          celui-ci, en commençant par le plus récent (\S\ref{sec:verrou-periode})~;
    \item \textbf{le règlement est rapproché} d'une transaction bancaire~: il faut d'abord
          annuler ce rapprochement --- le bouton \emph{Annuler le rapprochement} apparaît à
          côté du motif (\S\ref{sec:annuler-rapprochement})~;
    \item \textbf{un compte du règlement a été archivé}~: réactivez-le (son numéro est indiqué)~;
    \item \textbf{aucun exercice ouvert ne couvre la date du jour}, où la contre-passation serait
          datée~: créez-le.
\end{itemize}

\textbf{Lettrage.} Une facture \textbf{soldée} --- par ses règlements, par le solde de son reste
ou par son avoir --- est \textbf{lettrée} d'office, au moment du geste qui la solde~: sa ligne au
compte débiteurs et celles de ses règlements, de son solde ou de son avoir sur ce compte forment
un groupe de lettrage («~pièce~»)~: la colonne «~Lettrage~» du grand livre porte son code sur
chacune de ces lignes, et la facture quitte les postes ouverts du compte (\S\ref{sec:lettrage}) ---
ce que disent aussi la fiche, l'échéancier et les rappels. Un règlement partiel ne lettre rien. L'\textbf{annulation} d'un règlement --- ou du
solde, ou le dé-rapprochement (\S\ref{sec:annuler-rapprochement}) --- \textbf{délettre} la
facture~: sa créance et ses autres règlements redeviennent ouverts, et le règlement annulé est
lettré avec sa contre-passation.

\begin{keshnote}
\textbf{Mise à jour vers la version 0.13.} Les pièces déjà soldées --- factures clientes soldées,
factures fournisseurs payées --- et les contre-passations déjà passées sont \textbf{lettrées par
la mise à jour elle-même}, comme les gestes l'auraient fait. Deux réserves. Une pièce ou une
contre-passation dont \textbf{toutes les lignes sont en période close} (exercices clôturés,
période verrouillée) n'est pas lettrée~: elle reste ouverte, et ne se lettre pas à la main. Et ce
lettrage n'a \textbf{pas d'entrée au journal d'audit}~: une mise à jour n'a pas d'auteur.
\end{keshnote}

Kesh ne lettre une facture que si ses lignes sur le compte débiteurs se soldent \textbf{au grand
livre}~: un \textbf{avoir} émis avant la version 0.13 et porté sur un autre compte que le compte
débiteurs de la facture, ou une facture à la fois \textbf{créditée et réglée} (avant la version
0.12.1), ne sont pas lettrés --- leurs lignes ne se soldent pas. Un compte débiteurs qui n'est
\textbf{plus lettrable} (relié depuis à un compte bancaire, ou changé en compte de charge) n'est
plus lettré par un nouveau règlement --- celui-ci n'en est pas moins enregistré ---, mais un
lettrage posé avant le changement subsiste, et l'annulation d'un règlement le défait. Une
annulation qui défairait un lettrage figé par une période close est refusée (motif ci-dessus).

\begin{keshnote}
Les factures réglées avant la version qui a introduit les règlements (Kesh 0.12) portent la
mention «~payée~» sans règlement enregistré~: il n'y a rien à annuler sur elles.
\end{keshnote}

\subsection{Relancer les débiteurs}\label{sec:relances-user}

Depuis la version \textbf{0.7}, l'écran \emph{Factures → Rappels} (réservé aux rôles \emph{Comptable} et \emph{Administrateur}) liste les factures à relancer, groupées par débiteur, avec leur \textbf{prochain niveau} de rappel. Les niveaux, délais et frais se configurent par l'administrateur (\emph{Paramètres → Rappels débiteurs}).

\keshscreenshot{rappels-liste.png}{Écran Rappels : factures à relancer groupées par débiteur, prochain niveau, envoi unitaire ou par lot.}  % TODO capture: page /invoices/reminders

\begin{itemize}
    \item \textbf{Envoi unitaire} : un aperçu du courrier (dans la langue et le ton du niveau — 1er rappel courtois jusqu'à la mise en demeure) s'affiche, modifiable avant envoi ; un \textbf{rappel PDF} est joint (voir ci-dessous). Le destinataire est \textbf{verrouillé} sur l'e-mail de la fiche contact (sécurité).
    \item \textbf{Envoi par lot} : sélectionnez plusieurs factures (jusqu'à 20) et envoyez-les d'un coup, chacune à son prochain niveau ; un compte-rendu liste les envois réussis et les échecs par facture.
    \item \textbf{Rappel manuel} : enregistrez un rappel déjà envoyé hors Kesh (courrier, recommandé) — utile pour tracer un envoi papier. Le saut direct à un niveau élevé (mise en demeure) est autorisé.
    \item Un contact \textbf{sans adresse e-mail} reste visible (badge « sans e-mail ») : l'envoi e-mail est impossible, mais le rappel papier reste enregistrable.
    \item Une \textbf{protection anti-double-envoi} empêche qu'un double-clic n'envoie deux fois le même rappel. Une facture arrivée au \textbf{dernier niveau} reste affichée avec le badge « Dernier niveau atteint » (poursuite à envisager) — elle ne disparaît jamais silencieusement.
\end{itemize}

\paragraph{Ce que réclame un rappel.} Un rappel réclame ce que le client \textbf{doit encore} : le \textbf{reste dû} de la facture, après ses règlements partiels --- le montant total s'il n'a rien payé. Le courrier annonce ce reste, augmenté des frais de rappel s'il y en a ; à frais nuls, il ne parle pas de frais. Si un règlement est enregistré \textbf{entre l'ouverture de l'aperçu et l'envoi}, l'envoi est \textbf{refusé} : le courrier validé annoncerait un autre montant que la QR-facture jointe. Fermez puis rouvrez l'aperçu, qui reprend le nouveau reste. Le PDF joint n'est pas la facture réémise, c'est un \textbf{rappel}, produit à chaque envoi --- il n'est pas figé comme le PDF de la facture (\S\ref{sec:pdf-fige}) : il en porte le titre, rappelle le numéro de la facture d'origine et, sous le total, ce qui est \textbf{déjà réglé} puis le \textbf{reste à payer} (ces deux lignes n'apparaissent qu'après un règlement). Sa \textbf{QR-facture porte le reste dû}, avec la même référence que la facture. La réconciliation bancaire ne propose toutefois la facture que si le paiement arrive dans les \textbf{30 jours de la date de la facture} --- ce qui, après un rappel, est rarement le cas~: enregistrez alors le règlement depuis la fiche de la facture (\S\ref{sec:reglement-client}). Les \textbf{frais de rappel n'y sont jamais inclus} : ils figurent sur une ligne à part, avec la mention qu'ils ne sont pas compris dans le bulletin de versement.

\begin{keshwarning}
Les frais de rappel ne sont \textbf{pas comptabilisés} par Kesh. Un client qui paie le montant du courrier (reste dû \emph{et} frais) ne verra pas son versement \textbf{proposé} au rapprochement, puisqu'il dépasse ce qui reste à payer ; saisi comme règlement manuel, il est \textbf{refusé comme trop-perçu}. Seul le montant de la QR-facture s'enregistre (issue \#401).
\end{keshwarning}

\begin{keshnote}
Le rappel PDF ne part qu'avec un rappel \textbf{envoyé par e-mail}. Pour une mise en demeure envoyée hors de Kesh, ou pour un contact sans adresse e-mail, le seul PDF imprimable est celui de la \textbf{facture}, dont la QR porte toujours le \textbf{montant total} : ne le joignez pas tel quel à une facture déjà réglée en partie (issue \#477).
\end{keshnote}

\begin{keshtip}
Pour le \textbf{dernier niveau} (mise en demeure), privilégiez un envoi \textbf{recommandé} : il constitue une preuve de la mise en demeure en cas de poursuite. Envoyez le recommandé hors Kesh, puis enregistrez-le via \textbf{« Rappel manuel »} pour en garder la trace dans l'historique.
\end{keshtip}

\paragraph{Suspendre et reprendre les rappels d'une facture.} Sur la fiche d'une facture, vous pouvez \textbf{suspendre} ses rappels (avec un \textbf{motif} facultatif : litige, arrangement de paiement) puis les \textbf{reprendre}. Un badge \textbf{« Suspendu »} signale l'état (le motif apparaît au survol).

\begin{keshwarning}
Une facture suspendue \textbf{reste dans l'échéancier et la balance âgée} — elle sort seulement de la liste « à rappeler ». Rien n'est caché : c'est un garde-fou anti-dissimulation.
\end{keshwarning}

\paragraph{Historique des rappels.} La fiche de chaque facture affiche l'\textbf{historique de ses rappels} : date, niveau, canal (e-mail ou manuel), frais et destinataire. Un rappel annulé (par un administrateur) apparaît barré.

\subsection{Avoirs et notes de crédit}
\label{sec:avoirs}

Un avoir (note de crédit) annule une facture validée : c'est le moyen \textbf{correct} d'annuler une facture \emph{réellement émise} et déjà comptabilisée — l'écriture d'origine reste intacte, hors la marque de lettrage, qui la rattache à l'avoir (principe d'intangibilité comptable). Pour une facture validée \emph{par erreur} (jamais envoyée, ou créée pour un essai), il faut la \textbf{dévalider} d'abord --- elle redevient alors un brouillon, que l'on corrige ou que l'on supprime : voir \S\ref{sec:suppression-facture}. Pour créer un avoir, ouvrez la facture validée concernée et cliquez sur \textbf{« Créer un avoir »} (réservé aux rôles \emph{Comptable} et \emph{Administrateur}). En une seule étape, Kesh :

\begin{itemize}
    \item crée un avoir reprenant toutes les lignes de la facture, avec son propre numéro issu d'une séquence dédiée (par défaut \keshcommand{AV-2026-0001}, configurable dans \emph{Paramètres → Facturation}) ;
    \item génère automatiquement l'\textbf{écriture de contre-passation} : l'inverse exact de l'écriture de la facture (la créance client, le produit et la TVA due — une ligne par taux — sont contre-passés), de sorte que le \textbf{solde du client revient à zéro}. L'avoir crédite le \textbf{compte débiteurs que la facture a débité}, et contre-passe l'éventuel arrondi à 5~centimes sur le \textbf{compte que la facture a mouvementé} --- même si les comptes par défaut ont été changés dans \emph{Paramètres} → \emph{Facturation} depuis la validation : c'est ce qui ramène ces comptes à zéro. \textbf{Limite connue de cette version}~: la TVA, elle, est contre-passée sur le compte de TVA due \textbf{actuellement désigné} dans les réglages~; si ce réglage a changé entre la facture et son avoir, la TVA de la facture reste sur l'ancien compte et l'avoir débite le nouveau ;
    \item bascule la facture d'origine au statut \textbf{« annulée »} (sa propre écriture reste intacte, hors la marque de lettrage, qui la rattache à l'avoir).
\end{itemize}

L'avoir est consultable dans \emph{Avoirs} (menu \emph{Quotidien}), avec un \textbf{PDF « Avoir »} téléchargeable (mention « Avoir », référence à la facture d'origine, sans section QR de paiement). Le \emph{Décompte TVA} de la période tient compte de l'annulation (la TVA de la facture annulée est exclue).

\begin{keshwarning}
\textbf{Avoir refusé : « la facture porte un règlement ».} Un avoir annule \emph{toute} la facture. Si
le client en a déjà réglé une partie, l'avoir lui laisserait un crédit du montant encaissé, que
Kesh ne sait encore ni afficher, ni rembourser, ni imputer sur une autre facture. L'avoir est donc
refusé dès qu'un règlement existe --- même partiel ---, et Kesh le dit dans le dialogue.
\textbf{Annulez d'abord le règlement} depuis la fiche de la facture (\S\ref{sec:reglement-client}),
puis émettez l'avoir. Si le client ne paiera simplement pas le solde, ce n'est pas un avoir qu'il
faut, mais \textbf{solder le reste} en perte sur débiteur (\S\ref{sec:reglement-client}).
\end{keshwarning}

\begin{keshwarning}
\textbf{Avoir refusé : « compte de produit archivé ».} Si une ligne de la facture impute un compte de produit qui a été \textbf{archivé depuis la validation}, Kesh refuse d'émettre l'avoir et nomme la ligne et le compte concernés.

Ce refus est \textbf{volontaire}. Se replier sur le compte de produit par défaut du jour ferait extourner un compte que la facture n'a jamais crédité : le total resterait juste et le bilan équilibré, mais deux comptes de produit garderaient un solde résiduel se compensant, sans qu'aucun message ne le signale. Mieux vaut vous demander un geste explicite que fausser silencieusement votre compte de résultat.

\medskip
\noindent\textbf{Procédure quand l'avoir est néanmoins nécessaire} — trois étapes, à enchaîner \textbf{sans interruption} :

\begin{enumerate}
    \item \textbf{Réactiver} le compte dans le plan comptable (\S\ref{sec:plan-comptable-archivage}) ;
    \item \textbf{Émettre} l'avoir depuis la facture ;
    \item \textbf{Ré-archiver} le compte.
\end{enumerate}

Deux précisions utiles. D'abord, la réactivation peut elle-même être refusée dans deux cas (compte parent archivé, ou rôle repris entre-temps) : voir \S\ref{sec:plan-comptable-archivage}, qui explique quoi faire. Ensuite, tant que le compte est réactivé, il \textbf{réapparaît dans tous les sélecteurs de compte} de l'application — en cabinet, une écriture pourrait y être imputée par un collègue. D'où la recommandation d'enchaîner les trois étapes d'affilée plutôt que de laisser le compte ouvert le temps de finir la journée.

\medskip
\noindent Cette manœuvre n'a rien d'officieux : la réactivation \emph{et} le ré-archivage sont tous deux inscrits au \textbf{journal d'audit}. Un réviseur voit donc exactement pourquoi le compte a été rouvert quelques minutes, et dans quel ordre — ce qui préserve la valeur probatoire de votre comptabilité (CO art. 957a).
\end{keshwarning}

\begin{keshwarning}
\textbf{Avoir refusé : compte débiteurs, d'arrondi ou de TVA archivé.} L'avoir écrit aussi sur jusqu'à trois comptes qui ne sont pas des comptes de produit : le \textbf{compte débiteurs que la facture a débité}, le \textbf{compte qui a reçu son arrondi} à 5~centimes, s'il y en a un, et, si la facture porte de la TVA, le \textbf{compte de TVA due désigné dans les réglages}. Si l'un d'eux a été archivé, Kesh refuse d'émettre l'avoir avec le message «~Impossible d'émettre l'avoir --- compte(s) archivé(s)~: \ldots{} Réactivez le ou les comptes concernés.~», qui \textbf{nomme le compte}.

La marche à suivre est la \textbf{même procédure en trois étapes} que ci-dessus~: réactiver le compte, émettre l'avoir, ré-archiver. Si vous aviez entre-temps donné le rôle \emph{Créances clients} à un nouveau compte débiteurs, retirez-le d'abord de ce nouveau compte~: la réactivation est refusée tant qu'un autre compte porte le rôle (\S\ref{sec:plan-comptable-archivage}).

Pour le compte de \textbf{TVA due}, il existe une autre issue, et c'est souvent la bonne~: ce compte est celui \emph{des réglages}, non celui que la facture a crédité (limite connue, ci-dessus). Un administrateur peut donc désigner dans \emph{Paramètres} → \emph{Facturation} un compte de TVA due actif --- de préférence celui que la facture a crédité ---, plutôt que de réactiver un compte que la facture n'a peut-être jamais mouvementé.

Ce refus est distinct du précédent~: celui du compte de produit nomme une \textbf{ligne} de la facture, celui-ci un \textbf{compte}. Les deux peuvent se suivre~: le compte débiteurs, d'arrondi ou de TVA est contrôlé d'abord.
\end{keshwarning}

\begin{keshnote}
En version actuelle, l'avoir est \textbf{total} (il annule la facture entière). L'avoir \emph{partiel} et le remboursement d'une facture \emph{déjà encaissée} (un avoir est refusé sur une facture \textbf{réglée, même en partie}) sont prévus pour une version ultérieure. Les paiements sortants par fichier \texttt{pain.001.001.09} sont, eux, \textbf{disponibles} (cf. \S\ref{sec:paiements-fournisseurs}).
\end{keshnote}

\subsection{Supprimer ou dévalider une facture}
\label{sec:suppression-facture}

Une facture \textbf{brouillon} se supprime librement, par le bouton \textbf{« Supprimer »} de sa
fiche ou depuis la liste des factures. Cette suppression est réservée à l'\emph{Administrateur}.

\begin{keshwarning}
\textbf{Un brouillon peut déjà porter un numéro} --- c'est le cas d'une facture que vous venez de
\emph{dévalider} (voir ci-dessous). Le supprimer laisse alors un \textbf{trou définitif} dans la
séquence de numérotation~: le compteur ne redescend pas. Kesh vous le dit dans la fenêtre de
confirmation, en nommant le numéro concerné.

Ce trou est \textbf{assumé} et sans gravité~: il se voit et s'explique. Ce qui serait grave, c'est
la \emph{réattribution} --- deux factures différentes portant le même numéro ---, et elle n'a pas
lieu.
\end{keshwarning}

\medskip
\noindent\textbf{Une facture validée, elle, ne se supprime plus.} Elle se \textbf{dévalide}
d'abord.

Le bouton \textbf{« Dévalider »} de la fiche (rôles \emph{Comptable} et \emph{Administrateur})
repasse la facture en brouillon et \textbf{supprime son écriture comptable}, en une seule
transaction. Trois effets, que la fenêtre de confirmation rappelle~:

\begin{itemize}
    \item l'\textbf{écriture comptable} de la facture est supprimée~;
    \item la facture redevient un \textbf{brouillon}, modifiable~;
    \item son \textbf{numéro est conservé}. En la revalidant, elle reprendra \emph{le même} --- le
          compteur n'est pas consommé une seconde fois.
\end{itemize}

Deux sorties s'ouvrent alors, et c'est tout l'intérêt du geste~: \textbf{corriger puis revalider}
--- la facture retrouve son numéro et une écriture neuve ---, ou \textbf{supprimer} le brouillon.

\begin{keshtip}[title={\faLightbulb\ Pourquoi deux gestes plutôt qu'un}]
Jusqu'à la version 0.12, supprimer une facture validée effaçait son écriture comptable \emph{par
effet de bord} du mot « supprimer ». La destruction d'une écriture mérite d'être nommée~: c'est ce
que fait la dévalidation. Le geste est explicite, il est \textbf{journalisé sous son propre nom}
dans le journal d'audit, et il laisse la facture en place --- on peut la relire avant de décider.
\end{keshtip}

Kesh \textbf{refuse} la dévalidation dans les cas suivants, pour préserver la cohérence comptable.
Le refus \textbf{nomme toujours son motif}~:

\begin{itemize}
    \item facture \textbf{réglée}, même \textbf{partiellement}~: annulez d'abord le règlement
          (\S\ref{sec:reglement-client}), un paiement rapproché ne devant pas pointer dans le
          vide~;
    \item facture déjà \textbf{créditée par un avoir}~: l'avoir référence la facture, et il
          \emph{est} déjà la correction~;
    \item facture ayant un \textbf{historique de rappels}~: la dévalider effacerait la preuve du
          recouvrement~;
    \item facture \textbf{envoyée au client} par Kesh~: refus \textbf{sec}, que rien ne lève. Une
          facture que le client détient se corrige par un \textbf{avoir} (\S\ref{sec:avoirs}).
          ⚠️ Kesh ne connaît que ce qu'il a envoyé lui-même~: un PDF téléchargé puis transmis à la
          main ne bloque pas la dévalidation. Son premier téléchargement est bien inscrit au
          journal d'audit (gel du document, cf. \S\ref{sec:pdf-fige}), et la dévalidation
          \textbf{détache} ce document~: celui que vous avez transmis reste valable, mais n'est plus
          celui que Kesh conserve~;
    \item écriture \textbf{rapprochée} d'une transaction bancaire~;
    \item écriture d'un \textbf{exercice clos}~;
    \item écriture d'un exercice \textbf{suivi d'un exercice clôturé}~: son bilan, cumulatif,
          reprend l'écriture (\S\ref{sec:exercices-desordre})~;
    \item écriture déjà \textbf{contre-passée}~;
    \item écriture datée d'une \textbf{période verrouillée} (cf. \S\ref{sec:verrou-periode})~: la
          supprimer changerait les totaux d'une période déjà déclarée.
\end{itemize}

\begin{keshnote}
Toute dévalidation et toute suppression sont \textbf{inscrites au journal d'audit}~: Kesh
conserve qui a fait quoi et quand, ainsi qu'un instantané du contenu concerné --- la trace subsiste
même après la disparition de la facture.
\end{keshnote}

\begin{keshtip}
\textbf{Avoir ou dévalidation ?} Si la facture a été \emph{envoyée au client}, ou si elle appartient
à un exercice déjà déclaré, utilisez un \textbf{avoir} (\S\ref{sec:avoirs})~: l'historique reste
complet, et c'est la voie que la loi attend. Réservez la \textbf{dévalidation} aux erreurs de saisie
et aux essais qui n'auraient jamais dû être validés --- tant que l'exercice est ouvert et la facture
ni réglée, ni créditée, ni partie chez le client.
\end{keshtip}

%%─────────────────────────────────────────────────────────────────────────────
\section{Factures fournisseurs et paiements}\label{sec:factures-fournisseurs}
\label{sec:fournisseurs}

Le pendant « achats » de la facturation : enregistrer les factures reçues de vos
fournisseurs, les régler, et — pour les virements — produire un fichier de
paiement à charger dans votre e-banking. Ces écrans sont dans le menu
\emph{Quotidien}.

\subsection{Saisir une facture fournisseur}
\label{sec:saisie-facture-fournisseur}

\emph{Quotidien} → \emph{Factures fournisseurs} → \emph{Nouvelle facture}. Renseigner
le fournisseur (un contact marqué « fournisseur »), la date, l'échéance
éventuelle, les coordonnées de paiement (IBAN créancier, référence QR/ISR le cas
échéant) et les lignes (description, montant, compte de charge, taux de TVA
applicable — l'impôt préalable est comptabilisé automatiquement). À la validation,
Kesh passe l'écriture d'achat (charge + TVA préalable au débit, dette fournisseur
au crédit).

La dette fournisseur est portée au \textbf{compte créanciers} désigné dans
\emph{Paramètres} → \emph{Facturation} (section \emph{Comptes par défaut}), et
l'impôt préalable au compte de \textbf{TVA récupérable} de la même page. Si le
compte créanciers --- ou, pour une facture qui porte de la TVA, le compte de TVA
récupérable --- n'est pas imputable, la saisie est refusée avec un message qui
nomme le compte et renvoie à ces paramètres, où un administrateur doit le
remplacer~; rien n'est enregistré.

\subsection{Régler une facture fournisseur}\label{sec:regler-facture-fournisseur}

Sur la fiche d'une facture ouverte, le \textbf{règlement} est « binaire » — deux
cas selon le moyen de paiement :

\begin{itemize}
    \item \textbf{Virement bancaire} : vous choisissez le \emph{compte bancaire
    source}. Kesh solde la dette fournisseur par le crédit de ce compte bancaire.
    \item \textbf{Compte interne} (caisse, carte, compte de passage) : vous
    choisissez un \emph{compte de contrepartie}.
\end{itemize}

La contrepartie \textbf{ne peut pas être le compte créanciers de la facture}~: l'écriture serait
nulle, et la facture passerait «~payée~» sans paiement. Kesh le refuse dans tous les cas~--- de
même qu'un compte bancaire lié au compte créanciers~--- et ne le propose pas dans les listes,
\textbf{à condition} qu'il porte, au plan comptable, le \textbf{rôle} de compte de dettes
fournisseurs~: un compte qui l'a perdu reste proposé, et c'est l'enregistrement qui est refusé.
Choisissez un autre compte, ou reliez le compte bancaire à son propre compte de banque.

\textbf{Lettrage.} Une facture fournisseur \textbf{payée} est \textbf{lettrée} d'office~: sa
ligne d'achat au compte créanciers et celle du paiement sur ce compte forment un groupe de
lettrage («~pièce~»)~: la colonne «~Lettrage~» du grand livre porte son code sur ces deux lignes,
et la dette quitte les postes ouverts du compte créanciers (\S\ref{sec:lettrage}) --- ce que dit
aussi sa fiche. L'annulation du
règlement (ci-dessous) la \textbf{délettre}~: la dette redevient ouverte, et le paiement est
lettré avec sa contre-passation. Un compte créanciers \textbf{non lettrable} (relié à un compte
bancaire) ne forme aucun groupe --- le paiement n'en est pas moins enregistré. Une annulation qui
défairait un lettrage figé par une période close est refusée (motif ci-dessous). Les factures
déjà payées avant la mise à jour vers la version 0.13 sont lettrées par la mise à jour
(\S\ref{sec:reglement-client}, note).

\textbf{Annuler le règlement.} Un règlement enregistré par erreur ne se supprime pas~: sur la
fiche d'une facture payée, \emph{Annuler le règlement} (comptable ou administrateur) passe une
écriture inverse \textbf{datée du jour} et ramène la facture à \emph{ouverte}, à payer --- c'est
le geste à choisir quand c'est le \textbf{paiement} qui est faux et que vous voulez le refaire.
L'écriture d'origine et sa contre-passation restent toutes deux au grand livre, et le journal
d'audit nomme le geste. Un lien mène de la fiche à l'écriture de règlement. Si c'est la
\textbf{facture} elle-même qui est à écarter, payée ou non, voyez «~Annuler une facture
fournisseur~» ci-dessous~: il n'est pas nécessaire d'annuler d'abord son règlement.

Quand l'annulation n'est pas possible, le bouton est remplacé par la raison~: l'exercice du
règlement est \textbf{clôturé} (un administrateur doit d'abord rouvrir les exercices clôturés
jusqu'à celui-ci, en commençant par le plus récent), le \textbf{lettrage} du paiement avec sa
facture est figé par une période close (un administrateur fait reculer le verrou avant la date
la plus récente du lettrage et/ou rouvre les exercices clôturés, selon la cause ---
\S\ref{sec:verrou-periode}), un compte du
règlement a été \textbf{archivé} (réactivez-le), ou aucun exercice ouvert ne couvre la date du
jour (créez-le).

\subsection{Annuler une facture fournisseur}
\label{sec:annuler-facture-fournisseur}

Une facture fournisseur saisie par erreur --- doublon, mauvais fournisseur, facture qui n'était
pas due --- s'annule depuis sa fiche, \textbf{qu'elle soit ouverte ou déjà payée}~: bouton
\emph{Annuler la facture} (comptable ou administrateur). Kesh passe une écriture inverse de
l'achat, \textbf{datée du jour}, liée à l'écriture d'origine~: la charge et l'impôt préalable
reviennent à zéro, et la facture passe à \emph{annulée}. Le journal d'audit nomme le geste.

\textbf{Si la facture était payée}, son règlement \textbf{reste au grand livre}~: l'argent est
bien sorti. Il est \textbf{détaché} de la facture et devient un \emph{paiement sans facture}~---
au compte fournisseurs (2000), il laisse un solde débiteur, une somme que le fournisseur vous
doit ou qui attend d'être rattachée à la bonne facture. Ces paiements se retrouvent parmi les
\textbf{postes ouverts} du compte fournisseurs, s'il est lettrable (\emph{Mensuel} →
\emph{Postes ouverts}, \S\ref{sec:lettrage}), à son grand livre, et dans le journal d'audit, qui garde le lien vers
la facture annulée. La fiche de ce paiement n'offre ni «~Modifier~» ni «~Supprimer~»~: elle affiche le motif, et la contre-passation reste disponible. Si le paiement lui-même était une erreur, le plus simple est d'\textbf{annuler le
règlement d'abord}, puis la facture~: le lien avec la facture reste lisible au grand livre.

\textbf{Lettrage.} L'annulation d'une facture payée la \textbf{délettre}, et lettre l'achat avec
sa contre-passation. Le paiement détaché reste \textbf{ouvert} au compte créanciers~: il se
lettre à la main avec la ligne qui le solde, à l'écran \emph{Postes ouverts}
(\S\ref{sec:lettrage}), ou, contre-passé depuis sa fiche d'écriture, il est lettré
avec sa contre-passation.

Quand l'annulation n'est pas possible, le bouton est remplacé par la raison~: l'exercice de
l'\textbf{achat} est clôturé (un administrateur doit d'abord rouvrir les exercices clôturés
jusqu'à celui-ci, en commençant par le plus récent), le \textbf{lettrage} de la facture avec
son paiement est figé par une période close (un administrateur fait reculer le verrou avant la
date la plus récente du lettrage et/ou rouvre les exercices clôturés, selon la cause ---
\S\ref{sec:verrou-periode}), un compte de l'achat
a été \textbf{archivé} (réactivez-le), aucun exercice ouvert ne couvre la date du jour
(créez-le), ou la facture figure dans un \textbf{lot de paiement en cours} (annulez d'abord le
lot). Une facture annulée ne se \emph{désannule} pas~: si l'annulation était une erreur,
saisissez-la de nouveau.

\subsection{Paiements par fichier pain.001}
\label{sec:paiements-fournisseurs}

Pour régler un lot de factures par virement, Kesh génère un fichier de paiement
ISO 20022 \texttt{pain.001.001.09} (Swiss Payment Standards / SIX), à téléverser
dans l'e-banking. Le flux est en \textbf{deux temps} :

\begin{enumerate}
    \item \emph{Quotidien} → \emph{Paiements fournisseurs} : créer un lot en
    sélectionnant les factures à payer, puis télécharger le fichier
    \texttt{pain.001} généré et le déposer dans votre e-banking.
    \item Une fois le paiement exécuté par la banque, \textbf{confirmer} le lot
    dans Kesh : les écritures de règlement sont alors passées, et chaque facture
    payée est \textbf{lettrée} avec son achat, comme un paiement depuis sa fiche
    (un lot peut aussi être \textbf{annulé} tant qu'il n'est pas confirmé).
\end{enumerate}

Une facture dont le compte créanciers est le compte lié au compte bancaire du lot
\textbf{n'entre pas dans le lot}~: elle apparaît parmi les refus, avec son motif, et le fichier
ne la porte pas. Si le compte bancaire a été relié à ce compte \emph{après} la création du lot,
la \textbf{confirmation est refusée}, sans rien écrire~; deux issues~: relier de nouveau le compte
bancaire à son propre compte de banque puis confirmer, ou annuler le lot et régler la facture
depuis sa fiche. Le fichier \texttt{pain.001} a pu être exécuté par la banque~: c'est alors la
première qu'il faut choisir.

\begin{keshwarning}[title=\faExclamationTriangle\ Annuler le règlement d'une facture payée par lot]
Le règlement d'une facture payée par un lot \textbf{confirmé} peut être annulé --- le cas
d'usage est le \textbf{rejet} du virement par la banque. Le lot, lui, n'est pas modifié~: il
reste la trace de l'ordre transmis. La facture redevient \emph{ouverte} et peut repartir dans un
nouveau lot.

\textbf{Si la banque a réellement exécuté l'ordre}, la facture est déjà payée~: un nouveau lot
la paierait \textbf{deux fois}. Kesh vous le rappelle au moment d'annuler, en nommant le dernier
lot confirmé qui la contient. Dans ce cas, corrigez par un nouveau règlement direct depuis la
fiche, jamais par un nouveau lot.
\end{keshwarning}

\subsection{Importer des factures depuis un dossier}
\label{sec:import-factures}

Kesh peut lire, sans saisie de votre part, les factures fournisseurs suisses porteuses
d'un \textbf{QR-facture} (Swiss QR Code) déposées dans le dossier d'import (voir le
manuel administrateur pour sa configuration). L'import ne se déclenche pas tout seul~:
vous le lancez depuis l'écran décrit ci-dessous.

\emph{Quotidien} → \emph{Importer des factures}, puis \textbf{« Importer le
dossier »}. Kesh lit chaque fichier (PDF ou image), décode le QR côté serveur,
archive une copie du justificatif, et crée des factures fournisseurs
\textbf{« à compléter »} (coordonnées de paiement pré-remplies depuis le QR). Un
rapport indique les fichiers acceptés, en échec et les avertissements.

Pour chaque facture importée, cliquez sur \textbf{« Compléter »} : choisissez le
fournisseur et saisissez les lignes (compte de charge, TVA). Kesh vérifie que le
total TTC \textbf{correspond exactement} au montant du QR avant de créer la
facture fournisseur définitive (ensuite réglable comme ci-dessus). La complétion
crée cette facture par le même chemin que la saisie~: elle est refusée de la même
façon si le compte créanciers ou la TVA récupérable des réglages n'est pas
imputable (\S\ref{sec:saisie-facture-fournisseur}), et la facture importée reste
\emph{à compléter}. Vous pouvez
aussi \textbf{écarter} une entrée, et \textbf{consulter le justificatif d'origine}
depuis la fiche de la facture (« Voir la facture d'origine »).

\begin{keshnote}
Seules les factures en \textbf{CHF} sont acceptées à la complétion dans cette
version. Le justificatif importé reste consultable même après création de la
facture.
\end{keshnote}

%%─────────────────────────────────────────────────────────────────────────────
\section{Import bancaire}\label{sec:import-bancaire}

L'import bancaire est l'étape clé pour automatiser votre comptabilité. Kesh supporte deux formats.
Seul le CAMT.053 porte le \textbf{solde de clôture} du relevé, que le tableau de bord confronte
au solde comptable ; un import CSV n'en a pas :

\subsection{Import CAMT.053 (ISO 20022)}

\keshscreenshot{import-camt.png}{Prévisualisation d'un import de relevé bancaire CAMT.053 avant validation.}  % TODO capture: ecran import CAMT preview

Le format CAMT.053 est le standard suisse officiel pour les extraits de compte bancaire. Il est fourni par toutes les banques suisses en complément ou remplacement du format CSV propriétaire.

\subsubsection{Récupération du fichier CAMT.053}

\begin{itemize}
    \item \textbf{e-banking} : section « Téléchargement » ou « Extraits » → format CAMT.053 / ISO 20022 XML.
    \item \textbf{Multibanking} : si vous utilisez un agrégateur comme Multibanking ou Klippa, exportez en CAMT.053.
    \item \textbf{PostFinance Yellownet} : Menu Documents → Type de document : CAMT.053.
\end{itemize}

\subsubsection{Import dans Kesh}

\begin{enumerate}
    \item \emph{Quotidien} → \emph{Importer} → \emph{Nouvel import}.
    \item Sélectionner le compte bancaire cible (préalablement créé dans le plan comptable, marqué « Compte de banque »).
    \item Glisser-déposer le fichier XML CAMT.053 ou cliquer sur \emph{Parcourir}.
    \item Cliquer sur \emph{Importer}.
\end{enumerate}

Kesh parse le fichier, extrait les transactions, détecte les doublons (si le fichier contient des transactions déjà importées) et affiche un rapport :

\begin{itemize}
    \item Nombre de transactions importées.
    \item Nombre de doublons détectés (ignorés).
    \item Nombre de transactions rejetées (anomalies de format, montants invalides, etc.).
\end{itemize}

\subsection{Import CSV multi-encodage}

\keshscreenshot{import-csv.png}{Configuration d'un profil d'import CSV (mapping des colonnes, encodage, séparateur).}  % TODO capture: ecran profil CSV

\subsubsection{Profils banque}

Le format CSV n'étant pas standardisé, chaque banque propose son propre format de colonnes, séparateurs et encodage. Kesh permet de créer des \emph{profils banque} réutilisables :

\begin{enumerate}
    \item \emph{Administration} → \emph{Profils bancaires} → \emph{Nouveau profil}.
    \item Saisir le nom du profil (ex. « UBS CHF Compte courant »).
    \item Uploader un fichier CSV exemple.
    \item Mapper les colonnes : date, libellé, montant, référence, devise, etc.
    \item Configurer l'encodage : UTF-8, ISO-8859-1, Windows-1252, etc.
    \item Configurer le séparateur : virgule, point-virgule, tabulation.
    \item Sauvegarder.
\end{enumerate}

\subsubsection{Import avec un profil existant}

\begin{enumerate}
    \item \emph{Quotidien} → \emph{Importer} → \emph{Nouvel import CSV}.
    \item Sélectionner le profil banque adéquat.
    \item Sélectionner le compte bancaire cible.
    \item Uploader le fichier CSV.
    \item Cliquer sur \emph{Importer}.
\end{enumerate}

\subsection{Détection des doublons}

\keshscreenshot{import-doublons.png}{Avertissement de détection de doublons lors d'un ré-import de fichier bancaire.}  % TODO capture: ecran doublons

Kesh détecte automatiquement les transactions déjà importées en utilisant une empreinte unique calculée à partir de :

\begin{itemize}
    \item La date.
    \item Le montant.
    \item La référence bancaire (si fournie).
    \item Le libellé.
    \item Le compte bancaire cible.
\end{itemize}

Les doublons sont ignorés silencieusement (cf. Story Epic 8 « Détection doublons + rejet partiel »). Le rapport d'import indique le nombre de doublons trouvés.

\subsection{Rejet partiel}

Si un fichier contient un mélange de transactions valides et invalides, Kesh importe les valides et reporte les invalides dans un rapport. Vous pouvez corriger le fichier source et le réimporter (les valides déjà importées sont ignorées via détection de doublons).

\begin{keshtip}
Importez vos relevés bancaires \textbf{au moins une fois par semaine} pour faciliter la réconciliation. Plus le délai entre la transaction et l'import est court, plus il est facile de se souvenir du contexte et donc de rapprocher rapidement.
\end{keshtip}

%%─────────────────────────────────────────────────────────────────────────────
\section{Réconciliation bancaire}\label{sec:reconciliation}

La réconciliation consiste à passer en comptabilité les transactions bancaires importées. Kesh ne rapproche pas une transaction d'une écriture qui existerait déjà~: il \textbf{crée} l'écriture, soit en réglant une facture client, soit en imputant la transaction à un compte de contrepartie --- par une règle d'affectation, manuellement ou en l'éclatant sur plusieurs comptes.

\subsection{Réconciliation automatique}

\keshscreenshot{reconciliation-auto.png}{Propositions de réconciliation~: une ligne par transaction en attente, avec la candidate proposée, son score et la case à cocher.}  % TODO capture: ecran reconciliation auto (propositions)

\subsubsection{Algorithme de matching}

À l'ouverture de la page de réconciliation, Kesh cherche, pour chaque transaction bancaire en attente (créditrice, en CHF), les factures qui pourraient lui correspondre. Une facture est \textbf{candidate} si elle est validée et pas encore payée, si la transaction tombe dans les \textbf{30 jours qui précèdent ou suivent la date de la facture} --- sa date, pas son échéance --- et si \textbf{ce qui reste à payer} --- son \textbf{reste dû} : le montant total, moins l'avoir éventuel, moins les règlements déjà enregistrés --- égale le montant de la transaction à 5 centimes près. Un virement qui \textbf{solde} une facture déjà réglée en partie est donc reconnu : la proposition porte sur le reste, pas sur le montant d'origine. Un paiement reçu \textbf{plus de 30 jours après la date de la facture} --- une facture payée avec retard, ou après un rappel --- n'est en revanche \textbf{pas proposé}~: enregistrez alors le règlement depuis la fiche de la facture (\S\ref{sec:reglement-client}).

\paragraph{Une facture proposée peut être refusée à l'acceptation.}\label{sec:facture-posterieure} La fenêtre de 30 jours joue dans les deux sens, si bien qu'une facture datée \textbf{après} le paiement peut être proposée. Or un paiement ne peut pas précéder la facture qu'il règle de \textbf{plus d'un jour}~: à l'acceptation, Kesh compare la date de la facture à la \textbf{date de valeur} de la transaction --- à défaut, à sa date de comptabilisation --- et refuse la ligne, qui apparaît dans les \emph{Échecs partiels} avec le motif «~Le paiement est daté de plus d'un jour avant la facture~». Le même refus s'applique au règlement saisi depuis la fiche de la facture. Vérifiez la date de la facture~: si le paiement est réellement antérieur --- un acompte, par exemple ---, il ne règle pas cette facture.

Chaque candidate reçoit un score, somme de trois critères :

\begin{itemize}
    \item \textbf{Montant} (50\,\%) : le montant de la transaction égale exactement le reste dû \textbf{arrondi au centime}. Il n'y a pas de score intermédiaire : un écart d'un centime donne zéro sur ce critère.
    \item \textbf{Référence} (40\,\%) : la référence de la transaction contient le numéro de la facture, ou l'inverse (score plein) ; ou les deux commencent par les mêmes quatre caractères au moins (demi-score).
    \item \textbf{Contact} (10\,\%) : le nom de la contrepartie contient celui du client de la facture, ou l'inverse.
\end{itemize}

Une candidate dont le score vaut zéro n'est pas proposée. Quand aucune facture n'est candidate, la réconciliation propose à la place la \textbf{règle d'affectation} qui correspond à la transaction, s'il y en a une (\S\ref{sec:regles-affectation})~; elle s'affiche avec un score de 100\,\%. Rien n'est jamais passé en comptabilité sans votre acceptation : il n'y a ni seuil de confiance réglable, ni acceptation automatique.

\subsubsection{Accepter ou rejeter les propositions}

\emph{Mensuel} → \emph{Réconciliation}, puis choisir le compte bancaire. La page liste les transactions \textbf{en attente} de ce compte --- les 100 plus récentes au plus ---, une par ligne, avec leur date, leur montant, leur contrepartie, et \textbf{une} candidate~:

\begin{itemize}
    \item la candidate affichée est la \textbf{première} que propose Kesh~: la facture la mieux notée s'il y en a une, avec son numéro et le \textbf{montant qui reste à payer}, au centime --- si la facture est déjà réglée en partie, ce montant est suivi du total d'origine, par exemple «~600, reste dû sur 1000~»~;
    \item une \textbf{règle} n'est affichée, avec le badge \emph{Règle} et la mention «~libellé → compte~», que s'il n'y a \textbf{aucune} facture candidate --- même une facture de score faible passe devant elle~;
    \item sans candidate, la ligne indique «~Aucune correspondance~» et n'a pas de case à cocher.
\end{itemize}

Cochez les transactions à traiter, puis cliquez sur l'un des deux boutons placés au-dessus de la liste~:

\begin{itemize}
    \item \emph{Accepter (N)} --- N est le nombre de lignes cochées --- accepte, pour chaque ligne, \textbf{la candidate affichée}. Pour une facture, l'écriture de règlement est créée (Débit banque / Crédit débiteurs). Si la transaction solde une facture dont le reste porte plus de deux décimales, l'écart d'arrondi passe en troisième ligne, comme pour un règlement manuel (\S\ref{sec:reglement-client}) ; sans compte de différences d'arrondi, la ligne apparaît dans les \emph{Échecs partiels} avec le message «~aucun compte de différences d'arrondi utilisable n'est désigné~» --- désignez-en un dans \emph{Paramètres} → \emph{Facturation}. La proposition échoue aussi, avec son libellé, si le compte bancaire est lié au \textbf{compte débiteurs} de la facture (reliez-le à son propre compte de banque) ou si le compte de différences d'arrondi désigné \textbf{est} le compte débiteurs (un administrateur doit en désigner un autre dans \emph{Paramètres} → \emph{Facturation})~: l'écriture serait nulle. Pour une règle, l'écriture est passée entre la banque et le compte de contrepartie de la règle.
    \item \emph{Rejeter} écarte la proposition~: la transaction est marquée comme revue et \textbf{quitte la liste} de la réconciliation. Elle reste visible dans le \textbf{détail de son import} (\emph{Import bancaire}, puis l'import), avec le statut «~pending~», mais \textbf{sans aucune action}~: aucun geste de l'interface ne permet ensuite de la passer en comptabilité. Ne rejetez donc une transaction que si vous ne comptez pas la comptabiliser par ce chemin.
\end{itemize}

Pour imputer une transaction autrement qu'à la candidate affichée, utilisez les boutons de sa ligne~: \emph{Affecter manuellement} ou \emph{Éclater} (ci-dessous).

\paragraph{Un refus n'empêche pas les autres.} Chaque ligne cochée est traitée pour elle-même~: les transactions acceptées le sont, même si d'autres sont refusées. Kesh affiche le nombre d'opérations réussies puis, sous le titre \emph{Échecs partiels}, chaque transaction refusée --- désignée par sa date, son montant et sa contrepartie --- avec le motif du refus en clair~: période verrouillée, exercice fermé, facture déjà réglée entre-temps, compte non imputable, etc. Il en va de même pour \emph{Rejeter}. Ce bilan reste affiché même quand le lot a vidé la liste.

\subsection{Réconciliation manuelle}

\keshscreenshot{reconciliation-manuelle.png}{Affectation manuelle~: choix du compte de contrepartie d'une transaction.}  % TODO capture: ecran reconciliation manuelle

Si la candidate proposée ne convient pas, ou s'il n'y en a aucune~:

\begin{enumerate}
    \item Cliquer sur \emph{Affecter manuellement} sur la ligne de la transaction.
    \item Choisir le \textbf{compte de contrepartie}~: le sélecteur propose les comptes des classes 5, 6 et 7, et eux seuls, actifs et imputables. Compléter au besoin la description, la date de valeur et le projet analytique.
    \item Cliquer sur \emph{Affecter}.
\end{enumerate}

Kesh passe \textbf{une} écriture entre le compte de la banque et ce compte de contrepartie. L'affectation manuelle ne règle \textbf{pas} une facture~: le compte débiteurs n'est pas proposé. Pour un paiement de client que la réconciliation ne propose pas, enregistrez le règlement depuis la fiche de la facture (\S\ref{sec:reglement-client}).

\subsection{Éclatement d'une transaction (split)}

\keshscreenshot{reconciliation-split.png}{Éclatement d'une transaction bancaire~: ventilation de son montant sur plusieurs comptes, en une seule écriture.}  % TODO capture: ecran split

Cas d'usage : un prélèvement de carte de crédit de CHF 480 couvre des frais de déplacement, des fournitures et un abonnement, à imputer sur trois comptes de charges.

\begin{enumerate}
    \item Cliquer sur \emph{Éclater} sur la ligne de la transaction.
    \item Saisir les lignes~: pour chacune, un \textbf{compte} --- choisi dans le même sélecteur que l'affectation manuelle (classes 5, 6 et 7, actifs et imputables) ---, un montant, une description et, au besoin, un projet analytique. Deux lignes au moins, cinquante au plus~; la somme doit égaler le montant de la transaction.
    \item Valider.
\end{enumerate}

Kesh passe \textbf{une seule} écriture à plusieurs lignes~: la banque d'un côté, chaque compte de la ventilation de l'autre. Comme l'affectation manuelle, l'éclatement ne règle pas de facture.

\paragraph{La contrepartie n'est jamais le compte de la banque.} La contrepartie d'un rapprochement ne peut pas être le compte comptable du compte bancaire lui-même~: l'écriture passerait au débit et au crédit du même compte, et serait nulle. Kesh le \textbf{refuse} dans les trois cas --- affectation manuelle, éclatement et règle d'affectation (\S\ref{sec:regles-affectation})~; une règle refusée pour ce motif apparaît dans les \emph{Échecs partiels}, avec le libellé «~Erreur de validation~». En règle générale, le sélecteur de l'affectation manuelle ne propose pas ce compte, et une règle qui le vise n'est pas proposée --- la règle suivante qui correspond l'est à sa place. Ce filtrage n'est toutefois pas complet~: c'est le refus qui protège le grand livre.

\subsection{Règles d'affectation automatique}\label{sec:regles-affectation}

\keshscreenshot{reconciliation-regles.png}{Liste des règles d'affectation~: libellé, type de correspondance, valeur, compte de contrepartie, priorité et état.}  % TODO capture: ecran regles d'affectation

\emph{Administration} → \emph{Règles d'affectation}. Une règle reconnaît une transaction récurrente --- loyer, abonnement, frais bancaires --- et \textbf{propose} de l'imputer à un compte de charge ou de produit. Elle ne crée \textbf{aucune} écriture d'elle-même~: rien ne se passe à l'import, la proposition apparaît dans la réconciliation et c'est vous qui l'acceptez.

Exemple : le loyer est prélevé chaque mois par la régie. Une règle «~Loyer~», de type \emph{Contrepartie contient} et de valeur «~Régie~», vers le compte de charge du loyer, propose chaque mois d'imputer la transaction à ce compte.

\subsubsection{Création d'une règle}

\begin{enumerate}
    \item \emph{Administration} → \emph{Règles d'affectation} → \emph{Nouvelle règle}.
    \item Saisir le \textbf{libellé} de la règle (ex. «~Loyer~»).
    \item Choisir le \textbf{type} de correspondance --- \emph{Contrepartie contient}, \emph{Contrepartie exacte}, \emph{Référence contient} ou \emph{IBAN exact} --- et saisir la \textbf{valeur} cherchée. Une règle porte \textbf{une seule} condition. Le type ne se choisit qu'à la création~: il ne se modifie plus ensuite.
    \item Choisir le \textbf{compte de contrepartie}~: un compte de charge ou de produit, actif et imputable.
    \item Fixer la \textbf{priorité}, de 1 à 1000~: le \textbf{plus petit} nombre l'emporte.
    \item Au besoin, choisir un \textbf{projet analytique par défaut}.
    \item Cliquer sur \emph{Créer}.
\end{enumerate}

\subsubsection{Application des règles}

Une règle ne s'applique qu'aux transactions en CHF, et seulement quand \textbf{aucune facture} n'est candidate pour la transaction. Les règles actives sont essayées par \textbf{priorité croissante}, puis de la plus ancienne à la plus récente~; \textbf{seule la première} qui correspond est proposée. La colonne \emph{Appliquée} compte les acceptations.

\subsubsection{Désactiver, archiver, réactiver}

\emph{Désactiver} et \emph{Archiver} produisent le même état~: la règle n'est plus proposée, et la liste l'affiche «~Archivée~», grisée. \emph{Réactiver} la remet en service. Les écritures déjà passées par une règle ne sont pas touchées.

Une règle dont le \textbf{compte de contrepartie} a été archivé, ou est devenu \textbf{non imputable} (scindé en sous-comptes, par exemple), n'est \textbf{plus proposée}~: la règle suivante qui correspond l'est à sa place, s'il y en a une. La liste l'affiche pourtant toujours «~Active~» --- sa colonne d'état ne connaît que \emph{Active} et \emph{Archivée}~: si une règle «~Active~» ne propose plus rien, vérifiez son compte. Une règle dont le compte est non imputable se modifie, mais, une fois désactivée, ne se réactive pas sans changer de compte.

\subsection{Annuler un rapprochement}\label{sec:annuler-rapprochement}

Un rapprochement accepté par erreur --- mauvaise facture, mauvais compte de contrepartie ---
s'annule. On le fait depuis le \textbf{détail de l'import bancaire} (\emph{Import bancaire},
puis l'import concerné)~: chaque transaction rapprochée y porte un lien vers son écriture et un
bouton \emph{Annuler le rapprochement}. Quand le rapprochement a réglé une \textbf{facture
client}, le même bouton apparaît aussi sur la fiche de la facture, à côté du règlement qui ne
peut pas s'annuler seul.

Avant de confirmer, Kesh dit ce qui va se passer. À la confirmation~:

\begin{itemize}
    \item une \textbf{écriture inverse}, datée du jour, contre-passe l'écriture du rapprochement
          --- l'une et l'autre restent au grand livre~;
    \item si le rapprochement avait réglé une facture, son \textbf{règlement est retiré}~: la
          facture redevient à régler (ou partiellement réglée, si d'autres règlements
          subsistent)~;
    \item la transaction redevient \textbf{à rapprocher}, et réapparaît dans la réconciliation,
          prête à être rattachée à la bonne pièce.
\end{itemize}

Le journal d'audit nomme le geste. Quand l'annulation n'est pas possible, Kesh en donne la
raison~:

\begin{itemize}
    \item \textbf{la facture a été créditée} par un avoir~: son règlement reste ouvert au
          compte débiteurs, il ne s'annule pas --- cas hérité d'une version antérieure à la 0.12.1
          ou d'une sauvegarde restaurée~;
    \item \textbf{l'exercice de l'écriture du rapprochement est clôturé}~: un administrateur
          doit d'abord rouvrir les exercices clôturés jusqu'à celui-ci, en commençant par le
          plus récent. ⚠️ C'est l'exercice du \textbf{paiement} qui compte, jamais
          celui de la facture~: une facture de décembre dont l'exercice est clôturé, payée en
          janvier, se rapproche et se dé-rapproche normalement~;
    \item \textbf{le lettrage de la facture est figé par une période close}~: ce rapprochement
          est lettré avec sa facture, et toutes les lignes de ce lettrage sont dans une période
          close. Pour pouvoir l'annuler, prenez la date la plus récente du lettrage (en général
          celle du dernier règlement)~: si elle est sous le verrou de période, un administrateur
          doit faire reculer le verrou avant elle~; et si son exercice est suivi d'un exercice
          clôturé, ou clôturé lui-même, il doit rouvrir les exercices clôturés jusqu'à
          celui-ci, en commençant par le plus récent (\S\ref{sec:verrou-periode})~;
    \item \textbf{un compte de l'écriture a été archivé}~: réactivez-le (son numéro est
          indiqué)~;
    \item \textbf{aucun exercice ouvert ne couvre la date du jour}~: créez-le.
\end{itemize}

%%─────────────────────────────────────────────────────────────────────────────
\section{Lettrage et postes ouverts}\label{sec:lettrage}

Le solde d'un compte de tiers dit \emph{combien}~; il ne dit pas \emph{quoi}. Le
\textbf{lettrage} répond à la question~: il marque les lignes qui se soldent entre elles ---
une facture et son règlement, un acompte et sa reprise, une écriture et sa contre-passation
---, si bien que ce qui \textbf{reste ouvert} se lit ligne à ligne, et justifie le solde.

\subsection{Ce qu'est un lettrage}

Un lettrage est un \textbf{groupe de lignes d'un même compte}, dont la somme est
\textbf{exactement nulle}~: les débits égalent les crédits, au centime. Il réunit
\textbf{de 2 à 200 lignes}, et chaque ligne appartient à un groupe au plus. Kesh donne au
groupe un \textbf{code} de lettres (par exemple \texttt{AB}), qui s'affiche sur chacune
de ses lignes.

Seuls se lettrent les comptes d'\textbf{actif} et de \textbf{passif} qu'\textbf{aucun compte
bancaire ne désigne} --- typiquement les débiteurs, les créanciers, les comptes de passage et
d'acomptes. Les comptes de charges et de produits ne se lettrent pas, ni les comptes de banque~:
ceux-ci se rapprochent de leurs relevés (\S\ref{sec:lettrage-frontiere}).

\subsection{Qui lettre quoi}

Un groupe a l'une de trois origines, que sa vue nomme~:

\begin{itemize}
    \item \textbf{La pièce} («~lettrage de la pièce~\ldots~»). Kesh lettre \textbf{d'office} une
          facture client \textbf{soldée} --- par ses règlements, le solde de son reste ou son
          avoir (\S\ref{sec:reglement-client}) --- et une facture fournisseur
          \textbf{payée} (\S\ref{sec:regler-facture-fournisseur}). Un règlement partiel ne
          lettre rien~: la facture et son règlement restent ouverts côte à côte.
    \item \textbf{La contre-passation} («~contre-passation~»). Contre-passer une écriture lettre
          d'office chacune de ses lignes encore ouvertes sur un compte lettrable avec sa ligne
          miroir (\S\ref{sec:contre-passation}).
    \item \textbf{La main} («~lettrage manuel~»)~: tout le reste --- un acompte et sa reprise,
          une avance et son remboursement, un paiement fournisseur détaché de sa facture
          annulée et la ligne qui le solde --- se lettre à l'écran \emph{Postes ouverts},
          ci-dessous.
\end{itemize}

\subsection{L'écran \emph{Postes ouverts}}

\emph{Mensuel} → \emph{Postes ouverts}. Choisissez un \textbf{compte} --- seuls les comptes
lettrables sont proposés, archivés compris, signalés comme tels --- et une \textbf{date},
aujourd'hui par défaut. L'adresse de la page garde le compte, la date et le groupe ouvert~: un
signet, ou un rechargement, retrouve la même vue.

L'écran rappelle sa frontière en tête~: \emph{«~Cet écran solde entre elles les lignes des
comptes de tiers et de passage. Le rapprochement des relevés bancaires se fait dans Mensuel →
Réconciliation~; les comptes bancaires n'apparaissent pas ici.~»}

Chaque ligne ouverte dit sa date, son écriture (exercice et numéro), son journal, son libellé,
sa \textbf{pièce} --- facture, avoir, facture fournisseur, règlement, transaction bancaire, en
lien vers sa fiche quand elle en a une ---, son débit ou son crédit, et son \textbf{motif}, sur
deux lignes~:

\begin{itemize}
    \item \textbf{à la date choisie}~: «~non lettrée~», ou «~lettrée après cette date~», suivi
          du code du groupe, en lien, et du jour du lettrage~;
    \item \textbf{aujourd'hui}, l'état de la facture client qu'elle porte~: «~facture non
          réglée~» ou «~facture partiellement réglée~», avec le reste dû~; «~marquée payée avant
          la v0.12.0, sans écriture d'encaissement~: le compte porte encore cette créance~»~;
          ou «~pièce soldée, lettrage non posé~: période close ou avoir hérité --- rien à faire
          ici~» --- une pièce soldée dont le lettrage n'a pas été posé, parce que toutes ses
          lignes étaient en période close à la mise à jour, ou parce qu'un avoir antérieur à la
          version 0.13 a crédité un autre compte. Cette ligne n'a pas de case, et aucun geste de
          cet écran ne l'en sort.
\end{itemize}

\subsubsection{Les postes ouverts à une date}

Au pied de la liste, le \textbf{total des postes ouverts} et le \textbf{solde du compte} à la
date choisie --- et l'écran le dit~: \emph{«~Le total des postes ouverts au \ldots{} égale le
solde du compte au \ldots{} --- celui que la Balance montre pour ce compte à cette date, du côté
naturel du compte (solde cumulé depuis l'ouverture des livres).~»} Les deux montants
s'affichent en valeur absolue, suivie de leur sens, «~débiteur~» ou «~créditeur~»~; la Balance
et le Grand livre montrent le même nombre signé selon la nature du compte (positif au crédit
pour un passif). Une ligne lettrée \emph{après} la date choisie y figure encore~: à cette date,
elle était ouverte.

\begin{keshwarning}[title={\faExclamationTriangle\ «~Au 31 décembre~» n'est pas un instantané}]
La liste d'une date se calcule sur le lettrage \textbf{d'aujourd'hui}. Pour une date \textbf{en
période close} --- dans un exercice clôturé, ou au plus tard à la borne du verrou de période
(\S\ref{sec:verrou-periode}) ---, elle se reproduit à l'identique \textbf{tant qu'aucun
administrateur ne déverrouille la période, ne rouvre un exercice ni ne restaure une
sauvegarde}~: rien ne s'y lettre ni ne s'y délettre. Au-dessus de cette borne, un délettrage
ou l'annulation d'un règlement peut y faire \textbf{réapparaître} des lignes. Pour justifier un
solde de clôture, tirez la liste une fois l'exercice clôturé.
\end{keshwarning}

\subsubsection{Lettrer à la main}

Cochez les lignes qui se soldent~: l'écran affiche leur nombre et leur somme, et le bouton
\textbf{Lettrer} s'active quand la sélection réunit au moins deux lignes, au plus 200, de somme
nulle, dont \textbf{une au moins en période ouverte}~; sinon, il dit ce qui manque (l'écart, par
exemple). La sélection survit au changement de page. Seules les lignes lettrables à la main
ont une case~: une ligne déjà lettrée (son code est indiqué) ou une ligne d'une
\textbf{pièce} --- facture, avoir, facture fournisseur, règlement ---, dont le lettrage suit la
pièce et ses règlements, n'en ont pas. Une ligne seulement \textbf{rapprochée d'une transaction
bancaire}, elle, en a une~: la transaction ne la possède pas. Une fois le
groupe posé, son code s'affiche, en lien vers le groupe, et ses lignes quittent la liste. Si
une ligne a changé entre-temps (lettrée par quelqu'un d'autre, écriture modifiée ou
supprimée), Kesh refuse, le dit, et recharge la liste~: rien n'a été lettré.

\subsubsection{Les rapprochements proposés}

À côté de la liste, Kesh \textbf{propose} des paires évidentes~: une ligne au débit, une au
crédit, \textbf{du même montant exact}. Les contre-passations et leur origine viennent en tête,
puis les paires aux dates les plus proches (l'écart de jours est indiqué). Les propositions se
calculent sur les lignes ouvertes \textbf{aujourd'hui}, quelle que soit la date choisie.
\textbf{Kesh ne lettre jamais une proposition de lui-même}~: le bouton \emph{Lettrer} de la
paire le fait, à votre demande --- deux sociétés d'un même groupe, deux factures du même
montant, se ressemblent trop pour qu'un logiciel tranche seul. Une proposition devenue
périmée (une de ses lignes a été lettrée entre-temps) est refusée, et la liste se recharge.
Au-delà d'un certain nombre de lignes ouvertes sur le compte, Kesh ne propose rien et le dit~:
lettrez à la main.

Les lignes d'un règlement \textbf{annulé}, dont vous auriez ensuite délettré la paire de
contre-passation, reviennent ouvertes \textbf{sans pièce} --- le règlement n'est plus rattaché
à sa facture --- et sont de nouveau proposées en tête, comme toute contre-passation.

\subsubsection{Voir et défaire un groupe~: le délettrage}

Un groupe s'ouvre par son \textbf{code}, de partout où il s'affiche~: le message qui suit un
lettrage, le motif «~lettrée après cette date~» de la liste, la colonne \textbf{Lettrage} de la
fiche d'une écriture et du Grand livre (\S\ref{sec:grand-livre}), et le motif d'une écriture
que son lettrage empêche de modifier. La vue du groupe dit son compte, son origine et ses
lignes, avec leur pièce~; ouvert seul (depuis un lien), il offre le bouton \emph{Afficher les
postes ouverts de ce compte}, qui ramène à la liste.

Le bouton \textbf{Délettrer} rouvre toutes les lignes du groupe. Il n'est pas offert, et la vue
dit pourquoi, dans trois cas~:

\begin{itemize}
    \item \textbf{le groupe est celui d'une pièce}~: \emph{«~il suit la pièce et ses règlements,
          et ne se défait pas à la main~»}. Il se défait par les gestes de la pièce elle-même,
          que décrivent «~Enregistrer et annuler un règlement~» (\S\ref{sec:reglement-client}),
          «~Avoirs et notes de crédit~» (\S\ref{sec:avoirs}) et «~Factures fournisseurs et
          paiements~» (\S\ref{sec:factures-fournisseurs}~: règlement et annulation)~;
    \item \textbf{une de ses lignes appartient encore à une pièce}~: c'est la paire que forme
          l'annulation d'une facture fournisseur, \textbf{payée ou non} --- l'achat, qui reste la
          ligne de la facture, et sa contre-passation (\S\ref{sec:annuler-facture-fournisseur})~;
    \item \textbf{toutes ses lignes sont en période close}~: le lettrage n'y change plus
          (ci-dessous).
\end{itemize}

\subsubsection{La règle des périodes}

\textbf{Lettrer comme délettrer} exigent qu'\textbf{au moins une} ligne du groupe soit en
période ouverte~: dans un exercice ouvert, qu'aucun exercice postérieur clôturé ne suit, et
datée après le verrou de période. Un groupe dont \textbf{toutes} les lignes sont en période
close ne se crée ni ne se défait --- il se fige avec la période. Ce qui en découle pour
l'annulation d'un règlement, d'un solde, d'un rapprochement, d'un paiement ou d'une facture
fournisseur est décrit avec le verrou de période (\S\ref{sec:verrou-periode}) et avec chacun
de ces gestes.

\subsubsection{Ce que le lettrage ne fait pas}

\begin{itemize}
    \item \textbf{Pas de lettrage partiel.} Un groupe se solde exactement~: un acompte de 500
          repris pour 480 ne se lettre pas tel quel.
    \item \textbf{Pas de tolérance de montant.} Un écart de quelques centimes ne s'absorbe pas~:
          pour une facture client, \textbf{soldez son reste} depuis sa fiche
          (\S\ref{sec:reglement-client}) --- la facture se lettre alors d'office~; ailleurs,
          passez une écriture d'ajustement qui porte l'écart sur le compte, puis lettrez les trois
          lignes.
    \item \textbf{Une facture payée par une écriture saisie à la main ne se lettre pas avec
          elle}~: la ligne de la facture appartient à sa pièce. Enregistrez le paiement
          \textbf{sur la facture} (règlement, ou rapprochement bancaire), qui la solde et la
          lettre --- et corrigez l'écriture manuelle, qui ferait sinon double emploi.
    \item \textbf{Le code n'est pas dans les exports CSV et PDF du Grand livre}~: il ne s'y
          lit qu'à l'écran. L'export global (\S\ref{sec:exports}), lui, porte pour chaque ligne d'écriture
          la clé et l'origine de son groupe (colonnes \texttt{lettering\_key} et
          \texttt{lettering\_origin} de \texttt{journal\_entry\_lines.csv}).
    \item \textbf{Une écriture lettrée est figée}~: elle ne se modifie ni ne se supprime tant
          qu'une de ses lignes est lettrée. Quand le lettrage est ce qui la retient --- une
          écriture saisie à la main, d'un exercice ouvert ---, le motif de sa fiche porte le code
          du groupe, en lien~: délettrez-la là, puis modifiez (\S\ref{sec:modifier-ecriture}).
          Si ses lignes sont dans plusieurs groupes, le motif nomme le premier~: une fois celui-ci
          délettré, la fiche nomme le suivant. Une écriture qu'une autre cause fige (pièce,
          rapprochement bancaire, paiement détaché d'une facture fournisseur annulée,
          contre-passation, période close) affiche d'abord ce motif-là.
\end{itemize}

\subsubsection{La frontière avec la réconciliation, et les rôles}\label{sec:lettrage-frontiere}

Les \textbf{comptes bancaires ne se lettrent pas}~: le rapprochement des relevés se fait dans
\emph{Mensuel} → \emph{Réconciliation} (\S\ref{sec:reconciliation}), qui passe l'écriture d'une
transaction et, pour une facture, son règlement --- lequel lettre alors la facture d'office.

Tous les rôles lisent l'écran --- la liste, les propositions, les groupes. Le rôle
\textbf{Consultation} ne lettre ni ne délettre~: ni case, ni bouton. Les rôles
\textbf{Comptable} et \textbf{Administrateur} lettrent et délettrent, et chaque geste est inscrit
au journal d'audit (\S\ref{sec:tracabilite}), avec son auteur.

%%─────────────────────────────────────────────────────────────────────────────
\section{Rapports comptables}

\subsection{Bilan}

\keshscreenshot{rapport-bilan.png}{Bilan : actifs, passifs et capitaux propres à une date donnée.}  % TODO capture: rapport bilan

\emph{Mensuel} → \emph{Rapports} → \emph{Bilan}. Affiche, à une date donnée, les comptes
d'actif et de passif à solde non nul, \textbf{triés par numéro de compte}, avec les totaux
Actifs, Passifs et Capitaux propres.

\begin{keshwarning}[title=\faExclamationTriangle\ Le bilan n'est pas encore structuré selon l'art.~959a CO]
Kesh ne regroupe pas les comptes en \textbf{actif circulant} et \textbf{actif immobilisé}, ni
les dettes en \textbf{capitaux étrangers à court terme} et \textbf{à long terme}. Le plan
comptable porte bien une hiérarchie de comptes, mais les rapports ne l'utilisent pas encore.

\medskip
\noindent Seule la partition \textbf{capitaux propres / dettes} est faite, et elle l'est
correctement (voir ci-dessous). Pour des comptes annuels à remettre à un tiers --- fiduciaire,
organe de révision, banque ---, la mise en forme normée reste donc à produire hors de Kesh, à
partir des chiffres qu'il fournit.
\end{keshwarning}

La section \emph{Capitaux propres} est présentée \textbf{séparément des dettes} et \textbf{regroupée par rôle} de compte (capital, autres fonds propres, bénéfice / perte reporté), dans l'ordre légal de l'Art. 959a CO. Kesh distingue nettement deux grandeurs qu'il ne faut jamais confondre :

\begin{itemize}
    \item le \textbf{report d'ouverture physique} que vous avez saisi sur un compte (typiquement lors d'une reprise de comptabilité depuis un autre logiciel) — il apparaît sous son propre numéro de compte ;
    \item le \textbf{« Résultat reporté (calculé) »} que Kesh calcule automatiquement à partir des résultats de vos exercices antérieurs — ligne distincte, explicitement marquée \emph{calculé}, sans numéro de compte.
\end{itemize}

Ces deux lignes ne sont jamais additionnées ni fusionnées. Cette présentation se retrouve à l'identique dans les exports PDF et CSV du bilan.

Filtres : date de référence (par défaut fin d'exercice en cours), exercice à comparer (N-1).

\subsection{Compte de résultat}

\keshscreenshot{rapport-resultat.png}{Compte de résultat : produits, charges et résultat de la période.}  % TODO capture: rapport compte de resultat

\emph{Mensuel} → \emph{Rapports} → \emph{Compte de résultat}. Liste les comptes de produits
et les comptes de charges à solde non nul, et le résultat de la période.

\begin{keshwarning}[title=\faExclamationTriangle\ Le compte de résultat n'est pas encore structuré selon l'art.~959b CO]
Kesh ne propose \textbf{pas} les présentations \emph{par nature} et \emph{par fonction}, et ne
calcule aucun des sous-totaux légaux --- résultat brut, résultat d'exploitation, résultat
financier, résultat hors exploitation, impôts isolés. Vous obtenez deux listes et un résultat
net.

\medskip
\noindent Pour des comptes annuels destinés à un tiers, la mise en forme normée reste à
produire hors de Kesh, à partir de ces chiffres.
\end{keshwarning}

\subsection{Balance des comptes}

\emph{Mensuel} → \emph{Rapports} → \emph{Balance}. Pour chaque compte, sur la période choisie : le \textbf{solde
d'ouverture}, les \textbf{mouvements} débit et crédit, et le \textbf{solde de clôture} (ouverture + mouvements).

\begin{itemize}
    \item \textbf{L'ouverture suit la nature du compte.} Un compte de bilan (actif, passif) reporte son solde depuis
        l'origine, tous exercices confondus ; un compte de produits ou de charges repart de zéro au début de
        l'exercice — son ouverture ne compte que ce qui précède la période \emph{dans l'exercice}.
    \item \textbf{La balance concorde avec le bilan et le grand livre.} La clôture d'un compte de bilan est son solde au
        bilan à la fin de la période ; celle d'un compte de résultat, son montant \textbf{depuis le début de l'exercice} —
        égal au compte de résultat quand la période commence au premier jour de l'exercice, puisque le compte de
        résultat ne porte que les mouvements de la période choisie. Ouverture et clôture sont, dans tous les cas,
        celles du grand livre (\S\ref{sec:grand-livre}) sur la même période.
    \item \textbf{La ligne « Résultat reporté (calculé) ».} Kesh ne passe aucune écriture de clôture : le résultat des
        exercices antérieurs n'est porté par aucun compte. Une ligne calculée le montre, sous les comptes — la même
        valeur qu'au bilan, « Perte reportée » à l'écran si elle est négative.
    \item \textbf{Les montants sont signés selon la nature du compte} — positifs du côté naturel (débit pour un actif ou
        une charge, crédit pour un passif ou un produit). La somme visible des colonnes d'ouverture et de clôture n'est
        donc pas nulle ; le \textbf{contrôle} affiché sous l'ouverture (une coche) la vérifie en ramenant chaque solde au sens débit, résultat
        reporté compris. Un avertissement à sa place signale une écriture déséquilibrée antérieure à la période ; le PDF
        l'écrit en toutes lettres sous les totaux.
    \item \textbf{Premier exercice} : les soldes de départ saisis à la reprise sont une écriture datée du premier jour
        de l'exercice (\S\ref{sec:soldes-depart}) — ils apparaissent donc en \emph{mouvements}, et l'ouverture vaut zéro.
    \item \textbf{Complément des soldes de départ} : daté du jour (régularisation), il apparaît en \emph{mouvements} de
        l'exercice qui porte cette date, généralement un exercice ultérieur au premier.
    \item Un compte \textbf{archivé} figure tant qu'il porte un solde ou un mouvement.
\end{itemize}

Seul filtre : la période.

\subsection{Grand livre}\label{sec:grand-livre}

\emph{Rapports} → \emph{Grand livre}. C'est l'\textbf{extrait d'un compte} : ce qui a fait son
solde, ligne à ligne. Là où le bilan dit « clients : 84'320.15 » et où la balance clôt ce compte au
même nombre, le grand livre répond à la question qui suit toujours — \emph{de quoi ce solde
est-il fait ?}

Chaque compte y est encadré par trois lignes sans lesquelles ce n'est pas un extrait : le
\textbf{solde d'ouverture}, le \textbf{total des mouvements} (débit et crédit), le \textbf{solde
de clôture}. Entre les deux, une ligne par écriture, avec sa date, sa pièce, son journal, son
libellé, sa \textbf{contrepartie} — les autres comptes de la même écriture —, le \textbf{solde
progressif} après chaque mouvement et, en dernière colonne, son \textbf{lettrage}.

\begin{itemize}
    \item \textbf{Période libre, sans exercice.} C'est le seul rapport qui ne demande pas de
          choisir un exercice : il en franchit la borne. Quand la période en traverse deux, une
          ligne de rupture marque la clôture, et le solde d'un compte de charges ou de produits
          y \textbf{repart de zéro} — c'est ce que fait un bouclement.
    \item \textbf{Compte facultatif.} Sans sélection, tous les comptes mouvementés sont rendus.
    \item \textbf{Les comptes archivés apparaissent} s'ils portent encore un solde, signalés
          comme tels. Ce n'est pas un oubli de filtrage : un compte archivé à solde résiduel est
          précisément l'anomalie que l'on cherche.
    \item \textbf{Solde contre nature.} Un compte dont le solde tombe du côté opposé à sa nature
          — une caisse créditrice, un fournisseur débiteur — est signalé. C'est une alerte, pas
          une erreur : elle mérite un regard.
    \item \textbf{Accès direct depuis le bilan et la balance} : le numéro de compte y est un
          lien qui ouvre son grand livre sur la même période.
    \item \textbf{Lettrage.} La colonne \emph{Lettrage} porte le \textbf{code} du groupe d'une
          ligne lettrée, en lien vers ce groupe~; une ligne ouverte l'a vide
          (\S\ref{sec:lettrage}). L'en-tête d'un compte lettrable porte le lien
          \emph{Postes ouverts de ce compte}, qui ouvre l'écran des postes ouverts \textbf{au
          dernier jour de la période}~: leur total y égale le solde de clôture de l'extrait, au
          signe près --- l'extrait le donne du côté naturel du compte, l'écran en valeur absolue
          suivie de «~débiteur~» ou «~créditeur~». Le code ne figure \textbf{pas} dans les
          exports CSV et PDF du grand livre.
    \item \textbf{Exports PDF et CSV.} Le PDF répète l'en-tête du compte sur chaque page et
          porte « À reporter » au bas d'une page, « Report » en haut de la suivante, selon
          l'usage suisse. Le CSV écrit l'ouverture et la clôture comme des \textbf{lignes du
          tableau}, pour que le fichier reste additionnable dans un tableur.
\end{itemize}

\begin{keshnote}
L'écran s'arrête à \textbf{500 lignes par compte} et le dit quand il tronque. \textbf{Les
exports, eux, ne sont jamais tronqués} : c'est le fichier produit qui vaut présentation des
livres.
\end{keshnote}

\subsection{Balance âgée des créances}\label{sec:balance-agee}

Depuis la version \textbf{0.7}, l'onglet \emph{Rapports → Balance âgée} répartit l'\textbf{encours débiteur} --- le \textbf{reste dû} de chaque facture, TVA comprise, après ses règlements partiels et son avoir éventuel --- par \textbf{client} et par \textbf{tranche d'ancienneté}, arrêté à ce jour. C'est l'outil de pilotage du recouvrement : il montre d'un coup d'œil quelles créances vieillissent et lesquelles prioriser.

\keshscreenshot{balance-agee.png}{Balance âgée : reste dû par client et par tranche de retard, avec total général.}  % TODO capture: onglet Rapports/Balance âgée

\begin{itemize}
    \item \textbf{Colonnes} : \emph{Non échu} \textbar{} \emph{1-30} \textbar{} \emph{31-60} \textbar{} \emph{61-90} \textbar{} \emph{90+} jours de retard. Avec la colonne « Non échu », le \textbf{total général réconcilie} avec le solde du compte clients (débiteurs) du bilan \emph{tant que ce compte n'est mouvementé que par les factures, leurs règlements et leurs avoirs} (y compris leurs annulations). Toute autre écriture qui le touche l'en écarte : solde de départ, écriture saisie au journal, virement bancaire affecté à ce compte sans passer par une facture, paiement d'une facture fournisseur compensé sur ce compte, ou règlement enregistré avant la version 0.13.0 dont la contrepartie était ce compte lui-même. De même un changement du compte clients dans les réglages de facturation, ou des données antérieures à la version 0.12.1, restaurées ou mises à jour.
    \item \textbf{Détail} : une ligne par client, un total par tranche, et un total général. Chaque client renvoie d'un clic vers \textbf{ses factures}.
    \item \textbf{Export CSV} : réservé aux rôles \emph{Comptable} et \emph{Administrateur} (le rapport lui-même est consultable par tous les rôles).
\end{itemize}

\begin{keshnote}
Les factures dont les rappels sont \textbf{suspendus restent comptées} dans la balance âgée (elles ne sortent que de la liste « à rappeler »). Le montant réparti est le \textbf{reste dû} de chaque facture, jamais le montant hors taxe ni le total facturé : une facture de 1\,081.--- réglée de 900.--- y pèse 181.---, dans la tranche de \emph{son échéance} --- un règlement partiel ne rajeunit pas la créance.
\end{keshnote}

\subsection{Journal des écritures}

\emph{Mensuel} → \emph{Rapports} → \emph{Journal}. Liste chronologique de toutes les écritures de la période sélectionnée. Vue brute, idéale pour audit ou contrôle.

\subsection{Décompte TVA}\label{sec:decompte-tva}

\emph{Mensuel} → \emph{Rapports} → \emph{TVA}. Produit le décompte de la TVA pour la période sélectionnée. La TVA due est calculée \textbf{à partir des factures} validées de la période --- et non du grand livre ---, la TVA récupérable à partir du compte d'impôt préalable.

\begin{itemize}
    \item \textbf{TVA due} : la TVA facturée sur les ventes, ventilée par taux (8.1\%, 3.8\%, 2.6\%, 0\%). Elle est comptabilisée automatiquement lors de la validation de chaque facture de vente (compte \emph{TVA due}, p.~ex. 2200).
    \item \textbf{TVA récupérable} (impôt préalable) : la TVA payée sur les achats, lue du solde du compte d'impôt préalable (p.~ex. 1171). Elle s'alimente via l'assistant de saisie d'achat (cf. \emph{Saisie d'une écriture}), qui pré-remplit la ligne d'impôt préalable à partir d'un taux.
    \item \textbf{Diminutions de contre-prestation (soldes)} : la TVA corrigée par les \emph{soldes} (escompte accordé, perte sur débiteur --- \S\ref{sec:reglement-client}) datés dans la période, ventilée par taux. La section n'apparaît que s'il y a de tels soldes, suivie de la \textbf{TVA due nette} (TVA due − TVA des soldes).
    \item \textbf{Solde net dû à l'AFC} : TVA due \textbf{nette} − TVA récupérable. Un solde négatif signifie un crédit d'impôt en votre faveur.
\end{itemize}

\paragraph{Alerte d'écart --- ce qu'elle voit, et ce qu'elle ne voit pas.} Kesh compare le
montant de TVA due nette calculé depuis les factures et les soldes aux écritures comptables de ces mêmes factures et de ces mêmes soldes.
Si la ligne de TVA d'une facture a été \textbf{retouchée à la main} et ne correspond plus à la
TVA facturée, un \textbf{bandeau d'alerte} affiche l'écart. L'alerte est informative --- elle
ne bloque pas le rapport. Le décompte et l'écart figurent aussi dans les exports PDF et CSV.

\begin{keshwarning}[title=\faExclamationTriangle\ Cette alerte ne rapproche pas le décompte du grand livre]
Le contrôle porte sur les écritures \textbf{des factures et des soldes retenus par le décompte}, et non sur
le \textbf{solde du compte de TVA due}. Il applique donc le même filtre des deux côtés : une
facture qui sort du décompte sort aussi de la comparaison, et l'écart n'apparaît pas.

\medskip
\noindent Le cas le plus courant est celui de l'\textbf{avoir}. Émettre un avoir annule la
facture d'origine : elle quitte le décompte de sa période --- y compris si cette période a
déjà été déclarée --- et la contre-passation est datée du jour de l'avoir, où le décompte ne
la reprend pas non plus. Le grand livre, lui, est juste dans les deux périodes.
\textbf{L'alerte n'en dira rien}, et affichera même que tout concorde.

\medskip
\noindent Le \textbf{solde annulé} suit le même chemin. Annuler un escompte ou une perte retire
sa diminution de la période du solde --- y compris si elle a déjà été déclarée ---, et la
contre-passation, datée du jour, n'est reprise nulle part : \textbf{la TVA reprise n'apparaît
dans aucun décompte}, et l'alerte reste muette. Reportez-la à la main. N'apparaissent pas non
plus les parts d'un solde sans TVA (lignes exonérées, ou part trop petite pour porter un
centime de TVA).

\medskip
\noindent \textbf{Conséquence pratique.} Un décompte relancé plus tard peut donner un autre
chiffre que celui qui a été déposé, sans que Kesh le signale. \textbf{Conservez vos décomptes
déposés hors de Kesh} et rapprochez-les vous-même du solde du compte de TVA due --- c'est ce
que la concordance annuelle exige (art.~72 al.~1 LTVA), et Kesh ne l'outille pas aujourd'hui.
\end{keshwarning}

\begin{quote}
\small \emph{Limitation actuelle (v0.3)} : le format de décompte officiel AFC (e-décompte ESTV) et la méthode des taux de la dette fiscale nette (TDFN) ne sont pas encore couverts. Le décompte de Kesh fournit les montants (due, récupérable, solde net) à reporter manuellement sur le formulaire AFC.
\end{quote}

\subsection{Visualisation et exports}

Tous les rapports se \textbf{consultent à l'écran}. Pour les emporter :

\begin{itemize}
    \item \textbf{PDF et CSV} — bilan, compte de résultat, balance des comptes, grand livre,
          journal des écritures, décompte TVA, et les deux rapports par projet.
    \item \textbf{CSV seulement} — la \textbf{balance âgée des créances}, dont l'export est de
          surcroît réservé aux rôles \emph{Comptable} et \emph{Administrateur}.
\end{itemize}

\begin{keshnote}
Le PDF est paginé et destiné à être présenté ou archivé ; le CSV est destiné à être repris dans
un tableur. Aucun des deux n'est tronqué, même quand l'écran l'est.
\end{keshnote}

%%─────────────────────────────────────────────────────────────────────────────
\section{Comptabilité analytique par projet}

Kesh permet de rattacher chaque opération comptable à un \textbf{projet} pour l'analyser
isolément — typiquement pour capter \emph{toutes} les dépenses d'une rénovation
déductible, ou pour suivre le \emph{rendement} d'un investissement.

\subsection{Créer et organiser des projets}

Dans \emph{Administration} → \emph{Projets analytiques}, créez vos projets (code court + nom).
Un projet peut avoir des \textbf{sous-projets} (deux niveaux au maximum : projet →
sous-projet). Un projet dont des écritures dépendent ne se supprime pas mais
s'\textbf{archive} (il disparaît des sélecteurs mais reste dans l'historique et les rapports).

\subsection{Affecter un projet aux écritures}

L'affectation est \textbf{optionnelle} et disponible sur \textbf{tous les flux de saisie} :
\begin{itemize}
    \item \textbf{Écritures manuelles} — un sélecteur de projet \emph{par ligne} (une écriture
          de régularisation peut ainsi ventiler ses charges sur plusieurs projets).
    \item \textbf{Factures de vente et fournisseurs} — un projet au niveau du document,
          propagé à toutes les lignes de l'écriture générée.
    \item \textbf{Réconciliation bancaire} — projet lors d'une affectation manuelle,
          projet \emph{par ligne} lors d'une ventilation, et \textbf{projet par défaut}
          sur une règle d'affectation (chaque transaction réconciliée via la règle est
          alors taguée automatiquement).
\end{itemize}

\subsection{Les deux rapports par projet}

Dans \emph{Rapports}, deux onglets exploitent la dimension projet (exportables PDF/CSV,
en vue par \textbf{exercice} ou \textbf{cumulée depuis l'origine} du projet, avec
\textbf{agrégation des sous-projets}) :
\begin{itemize}
    \item \textbf{Dépenses par projet} — toutes les charges taguées, groupées par
          sous-projet puis par compte, avec sous-totaux et total. Chaque compte se
          \textbf{déplie} pour afficher les écritures qui le composent (drill-down).
          Idéal pour préparer une déclaration fiscale sans rien oublier.
    \item \textbf{Rendement par projet} — \textbf{coût investi} (charges + actifs immobilisés
          ; la trésorerie banque/caisse en est exclue pour ne pas gonfler le coût),
          \textbf{revenus}, \textbf{résultat net} et \textbf{rendement \%}
          (revenus / coût investi).
\end{itemize}

\begin{keshtip}
Les montants « Dépenses » sont hors TVA (la TVA récupérable étant portée sur un compte
d'actif séparé). Pour un usage privé non assujetti, il n'y a pas de TVA récupérable
distincte : le montant affiché correspond alors au décaissement réel.
\end{keshtip}

%%─────────────────────────────────────────────────────────────────────────────
\section{Exports de souveraineté}\label{sec:exports}

\subsection{Concept}

L'\emph{export ZIP de souveraineté} est l'export de vos données comptables dans un format
ouvert (CSV), accompagné d'un manifeste portant l'empreinte SHA-256 de chaque fichier.

\begin{keshwarning}[title={\faExclamationTriangle\ Ce que cet export couvre, et ce qu'il ne couvre pas}]
Il porte \textbf{toute votre comptabilité}~: ventes, achats, avoirs, règlements, banque,
grand livre, projets analytiques et journal d'audit --- la liste exacte suit. Ce qu'il ne
contient pas est dit plus bas, avec la raison~: des réglages techniques et des secrets, et,
point à retenir, les \textbf{fichiers} des pièces justificatives, dont seul le chemin figure.

\medskip
\noindent \textbf{Avant de compter sur cet export pour changer de logiciel}, vérifiez ce qu'il
contient réellement : la liste exacte figure sur l'écran d'export, et le manifeste
\texttt{metadata.json} énumère les fichiers produits. Pour une sauvegarde \emph{complète} de
l'installation, utilisez plutôt l'export \texttt{.keshbackup} de l'administration --- mais
c'est un format propriétaire, qui ne se relit que dans Kesh.

\medskip
\noindent \textbf{Sur la conservation à dix ans.} Le format ouvert et l'empreinte SHA-256 sont
une \textbf{bonne pratique d'archivage}, et ils facilitent la conservation exigée par
l'art.~958f CO. Ils ne constituent pas à eux seuls la garantie d'intégrité que l'Olico
(art.~9 al.~1 let.~b) demande pour un support modifiable : l'empreinte est calculée et
conservée par vous, sans horodatage ni tiers indépendant. Si vos livres doivent résister à une
contestation, complétez cet export par un dispositif d'archivage tiers.
\end{keshwarning}

\subsection{Contenu du ZIP}

L'archive contient \textbf{un fichier CSV par table}, plus un manifeste. Les trente tables
couvrent l'ensemble de votre comptabilité~:

\begin{itemize}
    \item \textbf{Écritures et plan comptable}~: \texttt{accounts.csv}, \texttt{journal\_entries.csv},
          \texttt{journal\_entry\_lines.csv}, \texttt{fiscal\_years.csv}.
    \item \textbf{Ventes}~: \texttt{invoices.csv} et \texttt{invoice\_lines.csv},
          \texttt{credit\_notes.csv} et \texttt{credit\_note\_lines.csv} (les avoirs),
          \texttt{invoice\_settlements.csv} (les règlements), \texttt{invoice\_reminders.csv}.
    \item \textbf{Achats}~: \texttt{supplier\_invoices.csv} et \texttt{supplier\_invoice\_lines.csv},
          \texttt{imported\_supplier\_invoices.csv} (les pièces reçues),
          \texttt{payment\_batches.csv} et \texttt{payment\_batch\_items.csv} (les ordres de virement).
    \item \textbf{Tiers et analytique}~: \texttt{contacts.csv}, \texttt{contact\_persons.csv},
          \texttt{products.csv}, \texttt{projects.csv}.
    \item \textbf{Banque}~: \texttt{bank\_accounts.csv}, \texttt{bank\_imports.csv},
          \texttt{bank\_transactions.csv}, \texttt{reconciliation\_rules.csv}, \texttt{bank\_profiles.csv}.
    \item \textbf{Paramètres et traçabilité}~: \texttt{company.csv}, \texttt{vat\_rates.csv},
          \texttt{company\_invoice\_settings.csv}, \texttt{dunning\_levels.csv},
          \texttt{company\_dunning\_settings.csv}, et \texttt{audit\_log.csv} --- le journal d'audit.
    \item \texttt{metadata.json}~: le manifeste (version de Kesh, date, auteur, et le hash
          SHA-256 de chaque fichier).
\end{itemize}

\begin{keshnote}
\textbf{Ce que l'export ne contient pas, et pourquoi.} Les \textbf{compteurs de numérotation}
(les numéros attribués figurent déjà sur les pièces), les \textbf{clés d'API et jetons de
session} (ce sont des secrets~: les exporter dans un fichier destiné à être transmis serait une
fuite), le \textbf{registre des utilisateurs} (sans perte~: chaque entrée du journal d'audit
porte le nom de son auteur \emph{au moment de l'écriture}), les \textbf{gabarits d'e-mail} et
l'\textbf{état d'installation} (configuration, non comptabilité).

⚠️ Les \textbf{pièces justificatives} --- les fichiers eux-mêmes --- ne sont \textbf{pas encore}
incluses. Le chemin de chaque pièce figure dans \texttt{imported\_supplier\_invoices.csv}, mais
le document doit être sauvegardé séparément. Il en va de même des \textbf{PDF des factures
émises} (\S\ref{sec:pdf-fige}) : \texttt{invoices.csv} porte leur empreinte, la date et la langue
du gel (colonnes \texttt{pdf\_sha256}, \texttt{pdf\_frozen\_at}, \texttt{pdf\_language}),
pas le document lui-même, qui se conserve avec le répertoire des documents de l'installation.
\end{keshnote}

\subsection{Génération}

\keshscreenshot{export-souverainete.png}{Génération d'un export de souveraineté (archive ZIP pour la conservation OLICo).}  % TODO capture: ecran export souverainete

\emph{Administration} → \emph{Export global}, puis \textbf{« Lancer l'export »}. L'export porte
\textbf{toute} votre comptabilité~: il n'y a ni exercice à choisir ni format à sélectionner. Le
ZIP est téléchargeable immédiatement (quelques secondes à une minute selon le volume).

\subsection{Vérification d'intégrité}

Le \texttt{metadata.json} contient un hash SHA-256 par fichier. Pour vérifier l'intégrité après
archivage~:

\begin{lstlisting}[language=bash]
unzip kesh-export-2026-05-20.zip -d kesh-export/
cd kesh-export/

for f in *.csv; do
    actual=$(sha256sum "$f" | cut -d' ' -f1)
    expected=$(jq -r ".tables[\"$f\"].sha256" metadata.json)
    if [ "$actual" = "$expected" ]; then
        echo "OK      $f"
    else
        echo "CORROMPU $f"
    fi
done
\end{lstlisting}

\subsection{Conservation OLICo 10 ans}

Pour la conservation légale 10 ans (CO Art. 958f al. 1) :

\begin{enumerate}
    \item Générer l'export ZIP à la clôture de chaque exercice.
    \item Calculer le SHA-256 du ZIP complet : \keshcommand{sha256sum kesh-export-2026.zip > kesh-export-2026.zip.sha256}.
    \item Stocker le ZIP + le fichier hash sur un support adapté (cloud avec immutabilité, NAS avec snapshot, DVD-R, etc.).
    \item Re-vérifier le hash périodiquement (recommandation : annuel) pour détecter toute corruption.
    \item Conserver le manuel utilisateur correspondant à la version Kesh utilisée (procédure documentée OLICo Art. 9 al. 1.b ch. 4).
\end{enumerate}

%%─────────────────────────────────────────────────────────────────────────────
\section{Isolation des données par société}

Kesh repose sur une architecture \textbf{multi-tenant} : le modèle de données
isole chaque société, et une installation en porte une. L'isolation
est appliquée au niveau de \textbf{chaque requête API} (portée par le
\texttt{company\_id}), garantissant qu'aucune donnée ne fuit d'une société à
l'autre.

\begin{keshnote}
Dans la version actuelle, \textbf{chaque utilisateur appartient à une seule
société} (son \texttt{company\_id} est porté par le jeton de connexion). Il n'y a
pas de sélecteur permettant de basculer entre plusieurs dossiers depuis un même
compte : un cabinet fiduciaire gère aujourd'hui chaque dossier client via une
instance dédiée --- un compte de la même installation appartient à la même société. Un sélecteur multi-dossiers est envisagé pour une version
ultérieure.
\end{keshnote}

%%─────────────────────────────────────────────────────────────────────────────
\section{Traçabilité (journal d'audit)}
\label{sec:tracabilite}

Kesh tient un \textbf{journal d'audit} (\texttt{audit\_log}) des actions
comptables significatives — création, modification et annulation d'entités,
imports, réconciliations, exports, gestion des utilisateurs, modification de
la société et la configuration initiale de l'installation — avec l'auteur,
l'horodatage et le contexte.

⚠️ \textbf{Une famille d'opérations n'y figure pas encore} : les \emph{gestes de
session} --- connexion, déconnexion, changement de votre propre mot de passe. C'est
un manque connu et suivi, pas un oubli.

\begin{keshnote}
Ce journal d'audit est \textbf{insert-only} : aucune entrée ne peut être modifiée,
et \textbf{aucune ne peut être supprimée sur une installation en service}. Deux
précisions, parce qu'une garantie vague ne vaut rien~: l'\textbf{import d'une
sauvegarde ne le remplace plus} --- les entrées de l'archive s'y ajoutent, les vôtres
restent ---, et la \textbf{réinitialisation des données de démonstration} reste
possible, mais elle est réservée à un administrateur \emph{et} refusée dès que
l'installation est finalisée~; elle efface le journal et y laisse \textbf{sa propre
entrée}, la \textbf{première} du journal neuf. Il \textbf{contribue} à la traçabilité attendue par
l'OLICo Art.~9 al.~1.b ch.~4 --- sans que cela vaille attestation de conformité~:
celle-ci se juge sur l'installation entière et son exploitation, non sur une table.
Il est \textbf{enregistré en base} à des fins de conformité et d'expertise, et se
\textbf{consulte à l'écran} (ci-dessous). Un \textbf{centre de notifications} (rappels d'échéance,
transactions à rapprocher) est prévu mais pas encore disponible.
\end{keshnote}

\subsection{Consulter le journal d'audit}

\emph{Administration} → \emph{Journal d'audit}. L'écran est réservé aux rôles
\textbf{Comptable} et \textbf{Administrateur}~; il n'est pas proposé au rôle
\emph{Consultation}, et le journal n'est pas lisible par une clé API.

\begin{itemize}
    \item \textbf{Les filtres}~: une période, un type d'entité, un numéro d'entité et une action,
          choisis dans des listes traduites dans la langue de votre installation. Le
          \textbf{numéro d'entité exige un type}~: le champ reste inactif tant qu'aucun type n'est
          choisi, et se vide quand vous changez de type --- un numéro ne désigne une pièce que dans
          son type. Un numéro nul, négatif ou non entier est ignoré.
    \item \textbf{Les dates filtrent sur des jours UTC}, alors que le tableau affiche l'heure
          \textbf{locale}. Une opération faite peu après minuit, heure suisse, appartient donc au
          jour UTC \emph{précédent}~: si elle manque, élargissez la période d'un jour. Une date de
          début postérieure à la date de fin est refusée par l'écran.
    \item \textbf{Le tableau}~: la date, l'auteur (tel qu'il s'appelait au moment de l'opération,
          et «~clé API~» si elle a été faite par une clé), l'action et le type d'entité
          \textbf{traduits}, le numéro de la pièce --- «~---~» quand l'opération ne vise pas une pièce
          précise --- et un bouton qui déplie le \textbf{détail} enregistré.
    \item \textbf{Une action ancienne}, écrite par une version antérieure et absente de la liste
          «~Action~», s'affiche sous son \textbf{code technique}~; on la retrouve en filtrant par le
          type et le numéro d'entité.
    \item \textbf{L'export CSV} reprend les filtres affichés, avec des heures en UTC et les libellés
          dans la langue de l'installation. Il est \textbf{plafonné à 10~000 lignes}~: au-delà, un
          message vous demande d'affiner les filtres. \emph{À l'export, une période inversée est
          refusée par le serveur avec un message qui n'est pas encore traduit.}
\end{itemize}

\emph{Ni la consultation ni l'export ne s'inscrivent eux-mêmes au journal.}

\paragraph{Les entrées de la configuration initiale.} L'assistant de
configuration inscrit \textbf{une entrée par étape franchie} (action
\emph{Étape d'installation franchie}, type \emph{Installation}). Les deux
premières --- la langue et le mode d'utilisation --- sont \textbf{communes aux deux
chemins}, production et démonstration~; les suivantes appartiennent à la
configuration de production. Chaque étape est précédée de ce qu'elle a changé~:
la société modifiée (langue, type d'organisation, langue comptable,
coordonnées), le compte bancaire créé ou modifié --- sans jamais l'IBAN en clair
---, puis, à la finalisation, les réglages de facturation, les taux de TVA et le
premier exercice. Le \textbf{plan comptable chargé s'inscrit en une seule entrée}
(\emph{Plan comptable chargé}, type \emph{Compte}), qui liste les comptes
créés~: filtrer sur le type \emph{Compte} et un numéro d'entité ne retrouve
\textbf{pas} la création d'un compte du plan livré --- ouvrez le détail de cette
entrée. Enfin, la société de départ créée par le \textbf{premier démarrage}
n'est pas inscrite~: seule la configuration qui suit l'est.

Sur le chemin de démonstration, le \textbf{chargement des données de
démonstration} s'inscrit en \textbf{une seule entrée} (\emph{Données de
démonstration chargées}, type \emph{Installation}), suivie de l'étape franchie~;
son détail donne la société, le nombre de comptes créés, l'exercice, le nombre de
taux de TVA créés et si les réglages de facturation l'ont été. Les données
elles-mêmes n'ont pas d'entrée chacune. ⚠️ \emph{Cette entrée n'est écrite que
lorsque le chargement aboutit.} Un chargement interrompu peut laisser une partie
des données de démonstration --- la société renommée, le plan comptable,
l'exercice --- \textbf{sans entrée au journal}, et un nouveau chargement échoue
alors jusqu'à une réinitialisation~: c'est une limite connue, suivie par
l'issue \#538. Il en va de même si la configuration de production franchit
l'étape pendant le chargement~: celui-ci est refusé, mais ces données restent sur
l'installation de production.

\paragraph{La réinitialisation.} Réinitialiser une installation (bouton
\emph{Réinitialiser pour la production} de la bannière de démonstration, qui n'agit
que si l'administrateur de l'installation l'a autorisée, voir «~Chemin A~») efface le
journal d'audit avec les données de la société, puis y inscrit \textbf{son propre
geste} --- une entrée \emph{Installation réinitialisée}, type \emph{Installation},
la seule du journal neuf. Son détail dit l'étape et le mode d'avant, le
\textbf{nombre d'entrées effacées} et la \textbf{plage de leurs numéros}, la liste des
tables vidées, et les utilisateurs ou clés d'API rattachés à la société. Les numéros
des entrées ne repartent pas de 1~: \textbf{le trou qui précède cette entrée
s'explique par elle}.

\paragraph{La réparation de l'installation.} Une entrée \emph{Réparation de
l'installation}, type \emph{Installation}, dit que Kesh a réparé de lui-même une
installation touchée par un défaut d'une version antérieure (\#528, \#542) ---
au démarrage, ou à l'import d'une sauvegarde. Personne ne l'a faite~: au
démarrage, elle est signée par le plus ancien administrateur actif. Son détail
nomme les utilisateurs rattachés, les clés d'API révoquées et les sociétés
provisoires supprimées.

%%─────────────────────────────────────────────────────────────────────────────
\section{FAQ et dépannage utilisateur}

\subsection{Mes écritures ne se balancent pas, que faire ?}

Vérifiez que la somme des débits égale exactement la somme des crédits (à 1 centime près). Kesh refuse d'enregistrer une écriture déséquilibrée. Si vous avez un écart d'arrondi sur de la TVA, utilisez le compte d'écart d'arrondi (6940 « Différences d'arrondi » dans les trois plans livrés).

\subsection{Mon QR Bill ne s'affiche pas correctement, que faire ?}

\begin{itemize}
    \item Vérifiez que votre compte bancaire est configuré avec un IBAN suisse valide.
    \item Si vous utilisez un QR-IBAN, assurez-vous qu'il est marqué comme tel dans \emph{Administration} → \emph{Comptes bancaires}.
    \item Testez le PDF généré avec une application e-banking suisse (PostFinance, Twint Pro, Raiffeisen) : le QR code doit être scanné en moins d'une seconde.
\end{itemize}

\subsection{Mon import CAMT.053 reporte 0 transactions, que faire ?}

\begin{itemize}
    \item Vérifiez que le fichier est bien au format CAMT.053 (extension \texttt{.xml} et balises ISO 20022 dedans).
    \item Vérifiez que le compte bancaire cible est correctement sélectionné.
    \item Vérifiez la date des transactions : si elles sont toutes antérieures à la date de début de votre exercice, elles sont ignorées.
    \item Consultez les logs d'import (lien \emph{Voir détails} sur l'import) pour identifier d'éventuelles erreurs de parsing.
\end{itemize}

\subsection{Une transaction bancaire ne propose aucun match automatique, est-ce normal ?}

Oui, si :

\begin{itemize}
    \item Aucune facture validée et non payée n'a un \textbf{reste dû} égal au montant de la transaction \textbf{à 5 centimes près}.
    \item La transaction tombe à plus de \textbf{30 jours de la date de la facture} --- de sa date, pas de son échéance~: une facture payée avec retard, ou après un rappel, n'est pas proposée. À l'inverse, une facture datée de plus d'un jour \textbf{après} le paiement peut être proposée, mais elle sera refusée à l'acceptation (\S\ref{sec:facture-posterieure}).
    \item La transaction est un débit, ou n'est pas en CHF~: seules les transactions créditrices en CHF reçoivent une facture candidate.
    \item Le score de la facture est nul~: la référence de la transaction ne contient pas le \textbf{numéro de la facture}, le nom de la contrepartie ne correspond pas au nom du client, et le montant n'est pas exactement le reste dû au centime.
    \item Aucune règle d'affectation ne correspond --- ou la règle qui correspondrait est écartée parce que son compte de contrepartie a été archivé ou est devenu non imputable. Une règle n'est d'ailleurs affichée que s'il n'y a \textbf{aucune} facture candidate.
\end{itemize}

Que faire ensuite dépend de la transaction~:

\begin{itemize}
    \item Pour un \textbf{paiement de client}, enregistrez le règlement \textbf{depuis la fiche de la facture} (\emph{Enregistrer un règlement}, \S\ref{sec:reglement-client}). N'utilisez ni \emph{Affecter manuellement} ni \emph{Éclater}~: ils ne proposent que les comptes des classes 5, 6 et 7, et la facture resterait due. Ne créez pas non plus de règle d'affectation pour un paiement de client~: elle porterait le paiement en produit --- le produit serait compté deux fois, et la facture resterait due, sans que rien le signale. La transaction bancaire reste ensuite dans la liste de la réconciliation~: Kesh ne la relie pas au règlement saisi depuis la fiche.
    \item Pour une \textbf{charge} diverse, utilisez \emph{Affecter manuellement} ou \emph{Éclater}~; pour une charge récurrente (loyer, abonnement, frais bancaires), créez une règle d'affectation.
\end{itemize}

\subsection{Comment changer la langue d'une facture ?}

La langue de la facture suit la langue de communication du client (champ \emph{Langue} du contact). Pour une facture ponctuelle dans une autre langue, vous pouvez forcer la langue dans le formulaire de création de facture.

\subsection{Comment exporter mes données pour migration ?}

Utilisez l'\emph{Export ZIP de souveraineté} (cf. section dédiée), en tenant compte de ce qu'il ne contient pas encore. Le format CSV ouvert se relit dans un tableur et s'importe dans la plupart des logiciels comptables.

\subsection{Comment puis-je revenir en arrière sur une action ?}

Selon l'action :

\begin{itemize}
    \item Écriture saisie à la main, d'un exercice ouvert~: le bouton «~Modifier~» de sa
          fiche --- chaque modification est tracée, avant et après~; une écriture dont une
          ligne est lettrée se délettre d'abord, depuis le lien du code que porte son motif
          (\S\ref{sec:lettrage}). Sinon, le bouton
          «~Contre-passer~»~: Kesh crée l'écriture inverse, garde l'origine, et lie les deux
          --- la correction reste visible. Une écriture saisie par erreur se
          \textbf{supprime} depuis la même fiche, dans le même cadre que la modification~:
          son numéro n'est pas réattribué, et la suppression est tracée avec son contenu.
          On atteint la fiche depuis la liste des écritures, en cliquant sur la ligne.
    \item Écriture appartenant à une pièce (facture, avoir, facture fournisseur,
          règlement, rapprochement) : elle n'est pas contre-passable directement --- passez
          par le chemin de la pièce, que l'écran vous indique.
    \item Écriture d'un exercice \textbf{clôturé}~: elle reste contre-passable, et la
          contre-passation tombe dans l'exercice \textbf{courant} --- c'est la pratique
          comptable, et il n'y a donc aucune raison de rouvrir l'exercice pour corriger une
          écriture. \textbf{Une exception}~: un \textbf{règlement} --- de facture client
          (\S\ref{sec:reglement-client}) comme de facture fournisseur ---, un
          \textbf{rapprochement bancaire} (\S\ref{sec:annuler-rapprochement}) et une
          \textbf{facture fournisseur} dont l'achat est dans un exercice clôturé
          (\S\ref{sec:annuler-facture-fournisseur}) ne s'annulent pas~: un administrateur doit
          d'abord rouvrir les exercices clôturés jusqu'à celui de la pièce, en commençant par
          le plus récent. De même, un règlement, un solde, un rapprochement, un paiement ou une
          facture fournisseur dont le \textbf{lettrage} est figé par une période close ne
          s'annule pas tant qu'un administrateur n'a pas fait reculer le verrou et/ou rouvert
          les exercices clôturés (\S\ref{sec:verrou-periode}).
    \item Facture validée : passage en avoir + génération d'une nouvelle facture.
    \item Facture fournisseur~: elle s'annule depuis sa fiche, ouverte ou payée
          (\S\ref{sec:annuler-facture-fournisseur})~; payée, son règlement reste au grand
          livre, détaché --- un paiement sans facture. La fiche de ce paiement n'offre ni «~Modifier~» ni «~Supprimer~»~: elle affiche le motif, et la contre-passation reste disponible.
    \item Rapprochement bancaire~: il s'annule depuis le détail de l'import bancaire, ou depuis
          la fiche de la facture qu'il a réglée (\S\ref{sec:annuler-rapprochement})~; son
          écriture est contre-passée, et un règlement de facture qu'il avait créé est retiré.
\end{itemize}

Le journal d'audit garde toujours la trace de l'action originale ET de l'action de retour.

%%─────────────────────────────────────────────────────────────────────────────
\section{Conseils d'utilisation au quotidien}

\subsection{Routine recommandée}

\paragraph{Quotidienne (5--10 min)}
\begin{itemize}
    \item Vérifier les notifications.
    \item Saisir les écritures du jour (s'il y en a).
\end{itemize}

\paragraph{Hebdomadaire (30--60 min)}
\begin{itemize}
    \item Importer les relevés bancaires de la semaine.
    \item Réconcilier les nouveaux mouvements (acceptation des propositions auto + manuelle pour le reste).
    \item Relire les écritures saisies dans la semaine.
    \item Émettre les factures de la semaine.
\end{itemize}

\paragraph{Mensuelle (2--4 h)}
\begin{itemize}
    \item Pointage du bilan et du compte de résultat.
    \item Génération des rapports d'analyse.
    \item Suivi des impayés (rappels de factures en retard).
\end{itemize}

\paragraph{Annuelle (1--2 jours)}
\begin{itemize}
    \item Clôture de l'exercice (génération bilan, P\&L, balance définitifs).
    \item Génération de l'export ZIP de souveraineté.
    \item Archivage légal 10 ans.
    \item Création du nouvel exercice.
    \item Test de restauration de sauvegarde (avec l'administrateur).
\end{itemize}

\subsection{Bonnes pratiques}

\begin{itemize}
    \item Saisissez les écritures \textbf{le plus tôt possible}, idéalement le jour même de la transaction.
    \item Utilisez des libellés \textbf{descriptifs et standardisés} : facilite la recherche et la lecture des rapports.
    \item Faites figurer le \textbf{numéro de la facture} dans la référence du paiement~: c'est lui que la réconciliation compare à la référence de la transaction.
    \item Vérifiez votre balance mensuellement : tout écart suspect doit être investigué immédiatement.
    \item Gardez votre plan comptable \textbf{simple et stable}. Évitez de créer des sous-comptes inutiles : ils alourdissent l'analyse.
\end{itemize}

%%─────────────────────────────────────────────────────────────────────────────
\section{Annexes}

\subsection{Glossaire comptable}

\begin{description}
    \item[Actif] Ce que possède l'entreprise (liquidités, créances, stocks, immobilisations). Classe 1 du plan comptable.
    \item[Amortissement] Constatation comptable de la perte de valeur d'une immobilisation sur sa durée d'utilisation.
    \item[Avoir (note de crédit)] Document qui annule tout ou partie d'une facture émise (l'inverse d'une facture).
    \item[Balance des comptes] Liste des comptes avec, sur une période, leur solde d'ouverture, leurs mouvements et leur solde de clôture.
    \item[Bilan] État du patrimoine d'une entreprise à une date donnée (actifs = passifs + capitaux propres).
    \item[Capitaux propres] Fonds appartenant aux propriétaires (capital, réserves, résultat reporté). Différence entre actifs et dettes.
    \item[Charge] Coût supporté par l'entreprise (achats, salaires, loyer). Classes 4 à 6.
    \item[Compte de résultat] État des produits et charges d'une période (bénéfice ou perte).
    \item[Créance] Somme due à l'entreprise par un tiers (typiquement un client). Compte « Débiteurs » (1100).
    \item[Débit / Crédit] Sens d'une opération comptable. Par convention, débit = entrée pour l'actif et sortie pour le passif.
    \item[Dette] Somme due à un tiers (fournisseur). Compte « Créanciers » (2000).
    \item[Échéance] Date à laquelle une facture devient exigible.
    \item[Écriture en partie double] Saisie respectant le principe somme(débits) = somme(crédits) (cf. \emph{Partie double}).
    \item[Exercice] Période sur laquelle s'établissent les comptes annuels.
    \item[Grand-livre] Ensemble des comptes individuels reprenant toutes les écritures qui les concernent.
    \item[IDE] Identifiant Unique des Entreprises (Suisse), format CHE-XXX.XXX.XXX.
    \item[Journal] Enregistrement chronologique de toutes les écritures comptables.
    \item[Journal d'audit (angl. audit log)] Registre des modifications comptables (qui, quand, quelle valeur), exigé par le CO Art. 957a. ⚠️ \emph{La couverture n'est pas totale --- voir la section «~Traçabilité~» pour la famille d'opérations qui n'y figure pas encore (les gestes de session).} Il se consulte à l'écran, \emph{Administration} → \emph{Journal d'audit} (\S\ref{sec:tracabilite}). À ne pas confondre avec le \emph{journal} des écritures, ci-dessus. Aucune entrée ne peut y être modifiée, ni supprimée sur une installation en service.
    \item[Lettrage] Marque posée sur des lignes d'un même compte d'actif ou de passif dont la
          somme est nulle, pour dire qu'elles se soldent entre elles (une écriture et sa
          contre-passation, un acompte et sa reprise\ldots). Le lettrage manuel, les
          rapprochements proposés et le délettrage se font à l'écran \emph{Postes ouverts}
          (\S\ref{sec:lettrage}). La contre-passation lettre d'elle-même une ligne et sa correction~; Kesh lettre de
          lui-même une facture client soldée avec ses règlements, son solde ou son avoir, et une
          facture fournisseur payée avec son paiement --- y compris, à la mise à jour vers la
          version 0.13, celles qui l'étaient déjà, sauf en période close~; un
          lettrage de pièce figé par une période close ne se défait pas~: l'annulation qui le
          défairait est refusée (\S\ref{sec:verrou-periode}). Une écriture dont une
          ligne est lettrée ne se modifie ni ne se supprime~: elle se délettre d'abord.
    \item[Liquidités] Avoirs immédiatement disponibles (caisse, comptes bancaires et postaux).
    \item[Partie double] Principe selon lequel toute opération est enregistrée sur au moins deux comptes, débit = crédit.
    \item[Passif] Origine des ressources de l'entreprise (dettes + capitaux propres). Classe 2.
    \item[Plan comptable] Liste structurée et hiérarchisée des comptes utilisés.
    \item[Postes ouverts] Lignes d'un compte lettrable qui ne sont pas lettrées à une date~: leur total égale le solde du compte à cette date (\S\ref{sec:lettrage}).
    \item[Produit] Revenu généré par l'activité (ventes, honoraires, cotisations). Classe 3.
    \item[Provision] Charge estimée et comptabilisée d'avance pour un risque ou une dépense future probable.
    \item[QR Bill] Standard suisse 2020+ de facture avec QR code (SIX Group).
    \item[QR-IBAN] IBAN spécial pour les paiements QR Bill avec référence structurée.
    \item[Réconciliation bancaire] Passage en comptabilité des transactions du relevé bancaire~: Kesh crée l'écriture, en réglant une facture client ou en imputant un compte de contrepartie.
    \item[Référence structurée] Référence numérique unique attachée à un paiement (ISO 11649 ou QRR).
    \item[Report à nouveau] Solde d'un compte de bilan reporté d'un exercice sur le suivant à la clôture.
    \item[Sous-compte (auxiliaire)] Compte de détail rattaché à un compte principal (ex. 1100.1 pour un client important).
    \item[TVA] Taxe sur la valeur ajoutée (Loi LTVA, RS 641.20).
\end{description}

\subsection{Raccourcis clavier}

\begin{tabularx}{\linewidth}{>{\bfseries}lX}
    \toprule
    \keshtableheader{Raccourci} & \keshtableheader{Action} \\
    \midrule
    \keshkey{Ctrl+S} & Sauvegarder / soumettre le formulaire en cours \\
    \keshkey{Ctrl+N} & Nouveau (en \emph{mode expert}, après l'onboarding : ouvre le formulaire de création adapté au contexte) \\
    \bottomrule
\end{tabularx}

\begin{keshnote}
La recherche globale (\keshkey{Ctrl+K}), l'export rapide et l'aide contextuelle
sont prévus pour une version ultérieure. \keshkey{Échap} ferme les boîtes de
dialogue qui le prennent en charge.
\end{keshnote}

\subsection{Nouveautés et limites (v0.3.0)}

Nouveautés de la version 0.3.0 :

\begin{itemize}
    \item TVA suisse comptabilisée : TVA due aux ventes (comptabilisée automatiquement à la validation des factures), impôt préalable sur les achats via l'assistant de saisie, décompte TVA (solde net dû à l'AFC, après les diminutions de contre-prestation des soldes) et contrôle de cohérence entre le décompte et les écritures des factures et des soldes qu'il retient (cf. \emph{Décompte TVA}, dont l'encadré dit ce que ce contrôle ne voit pas). Le décompte officiel AFC / e-décompte ESTV et le taux de la dette fiscale nette restent à venir.
\end{itemize}

Disponibles depuis la version 0.2.0 :

\begin{itemize}
    \item Récupération de mot de passe oublié par email en libre-service (cf. \emph{Premiers pas}).
    \item Export et import complets d'une installation depuis l'interface d'administration (format \keshcommand{.keshbackup} ; cf. manuel administrateur).
    \item Clés API (jetons d'accès personnels) pour les intégrations tierces (cf. \emph{Gestion de votre compte → Clés API}).
\end{itemize}

Fonctionnalités non encore disponibles :

\begin{itemize}
    \item Mono-devise CHF (multi-devises prévu pour une version ultérieure).
    \item Décompte TVA au format officiel AFC / e-décompte ESTV (les montants du décompte Kesh sont à reporter manuellement sur le formulaire AFC).
    \item Avoir \emph{partiel} et avoir sur facture \emph{déjà encaissée}, même en partie (l'avoir total sur une facture sans aucun règlement est, lui, disponible — cf. \S\ref{sec:avoirs}).
    \item Sélecteur de langue dans l'interface (la langue est fixée à l'onboarding ; le changement en cours de session est prévu pour une version ultérieure).
    \item Pas de signature électronique qualifiée (QES) sur les exports (verdict (b) de la recherche réglementaire, prévu Epic 14).
    \item Format PDF standard, non PDF/A (migration prévue selon ISO 19005).
\end{itemize}

\subsection{Liens externes utiles}

\begin{itemize}
    \item \href{https://www.kmu.admin.ch/kmu/fr/home/savoir-pratique/finances/comptabilite-et-revision.html}{Portail PME Confédération — Comptabilité}
    \item \href{https://www.estv.admin.ch/}{Administration fédérale des contributions (AFC)}
    \item \href{https://www.expertsuisse.ch/}{EXPERTsuisse}
    \item \href{https://www.treuhand-suisse.ch/}{TREUHAND|SUISSE}
    \item \href{https://www.six-group.com/en/products-services/banking-services/standardization/qr-bill.html}{QR Bill SIX (standard officiel)}
\end{itemize}

\subsection{Contact et support}

\begin{itemize}
    \item Documentation en ligne : \keshpath{https://github.com/guycorbaz/kesh/tree/main/docs}
    \item Issues et bugs : \keshpath{https://github.com/guycorbaz/kesh/issues}
    \item Discussions et questions : \keshpath{https://github.com/guycorbaz/kesh/discussions}
\end{itemize}

\clearpage

\keshversionnote

\end{document}
